%%%%%%%%%%%%%%%%%%%%%%%%%%%%%%%%%%%%%%%%%%%%%%%%%%%%%%%%%%%%%%%%%
\chapter{Unit 9: The Fundamental Theorem of Calculus}
\label{chap:unit9}
%%%%%%%%%%%%%%%%%%%%%%%%%%%%%%%%%%%%%%%%%%%%%%%%%%%%%%%%%%%%%%%%%

%===============================================================================
\section{The Fundamental Theorem of Calculus}
\label{sec:thef-1211}
%===============================================================================

The fundamental theorem of calculus is a critical portion of calculus because it links the concept of a derivative to that of an integral. As a result, we can use our knowledge of derivatives to find the area under the curve, which is often quicker and simpler than using the definition of the integral.


%-------------------------------------------------------------------------------
\subsubsection{Mean Value Theorem for Integration}
\label{sss:mean-b942}
%-------------------------------------------------------------------------------


We will need the following theorem in the discussion of the Fundamental

Theorem of Calculus.


Suppose $f(x)$ is continuous on $[a,b]$ . Then

$\frac{1}{b-a}\int\limits_a^b f(x)dx=f(c)$ for some $c\in[a,b]$.


%-------------------------------------------------------------------------------
\subsubsection{Proof of the Mean Value Theorem for Integration}
\label{sss:proo-439f}
%-------------------------------------------------------------------------------


$f(x)$ satisfies the requirements of the Extreme Value Theorem, so it

has a minimum $m$ and a maximum $M$ in $[a,b]$. Since


\[
\int_a^b f(x) dx = \lim _{n \to \infty} \sum _{k=1}^n f(x_k^*) \cdot \frac{b-a}{n} = \lim _{n \to \infty} \frac{b-a}{n} \cdot \sum _{k=1}^n f(x_k^{*})
\]


and since


\[
m\le f(x_k^{*})\le M
\] for all $x_k^{*}\in[a,b]$


we have

\begin{align}
& \lim _{n\to\infty}\frac{b-a}{n}\cdot\sum _{k=1}^n m &\le & \lim _{n\to\infty}\frac{b-a}{n}\cdot\sum _{k=1}^n f(x_k^{*})\le\lim _{n\to\infty}\frac{b-a}{n}\cdot\sum _{k=1}^n M \\& \lim _{n\to\infty} m n\cdot\frac{b-a}{n} &\le &  \int\limits_a^b f(x)\,dx\le\lim _{n\to\infty} M n\cdot\frac{b-a}{n} \\& \lim _{n\to\infty} m(b-a) &\le & \int\limits_a^b f(x)\,dx\le\lim _{n\to\infty} M(b-a) \\& m(b-a) &\le & \int\limits_a^b f(x)\,dx\le M(b-a) \\& m &\le & \frac{1}{b-a}\int\limits_a^b f(x)\,dx\le M
\end{align}


Since $f$ is continuous, by the Intermediate Value Theorem there is some

$f(c)$ with $c\in[a,b]$ such that


\[
\frac{1}{b-a}\int\limits_a^b f(x)dx=f(c)
\]


%-------------------------------------------------------------------------------
\subsubsection{Fundamental Theorem of Calculus}
\label{sss:fund-9421}
%-------------------------------------------------------------------------------


%-------------------------------------------------------------------------------
\subsubsection{Statement of the Fundamental Theorem}
\label{sss:stat-adcd}
%-------------------------------------------------------------------------------


Suppose that $f$ is continuous on $[a,b]$ . We can define a function $F$

by


\[
F(x)=\int\limits_a^x f(t)dt\quad\text{for }x\in[a,b]
\]


Suppose $f$ is continuous on $[a,b]$ and $F$ is defined by


\[
F(x)=\int\limits_a^x f(t)dt
\] Then $F$ is differentiable on $(a,b)$

and for all $x\in(a,b)$,


\[
F'(x)=f(x)
\]


When we have such functions $F$ and $f$ where $F'(x)=f(x)$ for every $x$

in some interval $I$ we say that $F$ is the \textbf{antiderivative} of $f$ on

$I$.


Suppose that $f$ is continuous on $[a,b]$ and that $F$ is any

antiderivative of $f$ . Then


\[
\int\limits_a^b f(x)dx=F(b)-F(a)
\]


% [Image: Figure 1]


Note: a minority of mathematicians refer to part one as two and part two

as one. All mathematicians refer to what is stated here as part 2 as The

Fundamental Theorem of Calculus.


%-------------------------------------------------------------------------------
\subsubsection{Proofs}
\label{sss:proo-dbb5}
%-------------------------------------------------------------------------------


%-------------------------------------------------------------------------------
\subsubsection{Proof of Fundamental Theorem of Calculus Part I}
\label{sss:proo-b3b2}
%-------------------------------------------------------------------------------


Suppose $x\in(a,b)$ . Pick $\Delta x$ so that $x+\Delta x\in(a,b)$ .

Then


\[
F(x)=\int\limits_a^x f(t)dt
\] and


\[
F(x+\Delta x)=\int\limits_a^{x+\Delta x}f(t)dt
\]


Subtracting the two equations gives


\[
F(x+\Delta x)-F(x)=\int\limits_a^{x+\Delta x}f(t)dt-\int\limits_a^x f(t)dt
\]


Now


\[
\int\limits_a^{x+\Delta x}f(t)dt=\int\limits_a^x f(t)dt+\int\limits_x^{x+\Delta x}f(t)dt
\]

so rearranging this we have


\[
F(x+\Delta x)-F(x)=\int\limits_x^{x+\Delta x}f(t)dt
\]


According to the Mean Value Theorem for Integration, there exists a

$c\in[x,x+\Delta x]$ such that


\[
\int\limits_x^{x+\Delta x}f(t)dt=f(c)\cdot\Delta x
\] Notice that $c$

depends on $\Delta x$ . Anyway what we have shown is that


\[
F(x+\Delta x)-F(x)=f(c)\cdot\Delta x
\] and dividing both sides by

$\Delta x$ gives


\[
\frac{F(x+\Delta x)-F(x)}{\Delta x}=f(c)
\]


Take the limit as $\Delta x\to0$ we get the definition of the derivative

of $F$ at $x$ so we have


\[
F'(x)=\lim _{\Delta x\to 0}\frac{F(x+\Delta x)-F(x)}{\Delta x}=\lim _{\Delta x\to 0}f(c)
\]


To find the other limit, we will use the \textbf{squeeze theorem}.

$c\in[x,x+\Delta x]$ , so $x\le c\le x+\Delta x$ . Hence,


\[
\lim _{\Delta x\to0}\Big[x+\Delta x\Big]=x\quad\Rightarrow\quad\lim _{\Delta x\to 0}c=x
\]


As $f$ is continuous we have


\[
F'(x)=\lim _{\Delta x\to0}f(c)=f\left(\lim _{\Delta x\to 0}c\right)=f(x)
\]

which completes the proof.


$\blacksquare$


%-------------------------------------------------------------------------------
\subsubsection{Proof of Fundamental Theorem of Calculus Part II}
\label{sss:proo-be95}
%-------------------------------------------------------------------------------


Define $P(x)=\int\limits_a^x f(t)dt$ . Then by the Fundamental Theorem of Calculus part I we know that $P$ is differentiable on $(a,b)$ and for all $x\in(a,b)$


\[
P'(x)=f(x)
\] So $P$ is an antiderivative of $f$ . Since we were assuming that $F$ was also an antiderivative for all $x\in(a,b)$ ,


\begin{align}&P'(x)=F'(x)\\\&P'(x)-F'(x)=0\\\&\Big(P(x)-F(x)\Big)'=0\end{align}


Let $g(x)=P(x)-F(x)$ . The Mean Value Theorem applied to $g(x)$ on

$[a,\xi]$ with $a<\xi$ \textbf{Power Rule of Integration I}

> \[
\int\limits_a^b x^ndx=\frac{x^{n+1}}{n+1}\Bigg| _a^b=\frac{b^{n+1}-a^{n+1}}{n+1}
\] as long as $n\ne-1$ and $0\notin[a,b]$ .


Notice that we allow all values of $n$ , even negative or fractional. If

$n>0$ then this works even if $[a,b]$ includes $0$ .


> \textbf{Power Rule of Integration II}

> \[
\int\limits_a^b x^ndx=\frac{x^{n+1}}{n+1}\Bigg| _a^b=\frac{b^{n+1}-a^{n+1}}{n+1}
\]

> as long as $n>0$ .


Examples


\begin{itemize}
    \item To find $\int\limits_1^2x^3dx$ we raise the power by 1 and have to
\end{itemize}
    divide by 4. So


\[
\int\limits_1^2x^3dx=\frac{x^4}{4}\Bigg| _1^2=\frac{2^4}{4}-\frac{1^4}{4}=\frac{15}{4}
\]


\begin{itemize}
    \item The power rule also works for negative powers. For instance
\end{itemize}

\begin{align}
\int\limits_1^3\frac{dx}{x^3} &= \int\limits_1^3 x^{-3}\,dx \\&= \frac{x^{-2}}{-2}\Bigg| _1^3 \\&= \frac{1}{-2}\left(3^{-2} - 1^{-2}\right) \\&= -\frac{1}{2}\left(\frac{1}{3^2} - 1\right) \\&= -\frac{1}{2}\left(\frac{1}{9} - 1\right) \\&= \frac{1}{2} \cdot \frac{8}{9} \\&= \frac{4}{9}
\end{align}


\begin{itemize}
    \item We can also use the power rule for fractional powers. For instance
\end{itemize}

\[
\int\limits_0^5\sqrt{x}dx=\int\limits_0^5x^\frac12dx=\frac{x\sqrt{x}}{\frac32}\Bigg| _0^5=\frac23\left(5^\frac32-0^\frac32\right)=\frac{10\sqrt5}{3}
\]


\begin{itemize}
    \item Using linearity the power rule can also be thought of as applying to
\end{itemize}
    constants. For example,


\[
=\int\limits_3^{11}7dx=\int\limits_3^{11}7x^0dx=7\int\limits_3^{11}x^0dx=7x\Bigg| _3^{11}=7(11-3)=56
\]


\begin{itemize}
    \item Using the linearity rule we can now integrate any polynomial. For
\end{itemize}
    example


\[
\int\limits_0^3(3x^2+4x+2)dx=(x^3+2x^2+2x)\Bigg| _0^3=3^3+2\cdot 3^2+2\cdot3-0=27+18+6=51
\]


%-------------------------------------------------------------------------------
\subsubsection{Exercises}
\label{sss:exer-ef8a}
%-------------------------------------------------------------------------------


%-------------------------------------------------------------------------------
\subsubsection{Questions}
\label{sss:ques-4786}
%-------------------------------------------------------------------------------


1.  Evaluate $\int\limits_0^1x^6dx$

2.  Evaluate $\int\limits_1^2x^6dx$

3.  Evaluate $\int\limits_0^2x^6dx$


%-------------------------------------------------------------------------------
\subsubsection{Solutions}
\label{sss:solu-8692}
%-------------------------------------------------------------------------------


1.  $\frac17=0.\overline{142857}$

2.  $\frac{127}{7}=18.\overline{142857}$

3.  $\frac{128}{7}=18.\overline{285714}$


%-------------------------------------------------------------------------------
\subsubsection{Resources}
\label{sss:reso-55b5}
%-------------------------------------------------------------------------------

\textbf{The Fundamental Theorem of Calculus, Part 1}

\begin{itemize}

\item \href{https://mathresearch.utsa.edu/wikiFiles/MAT1193/The\%20Fundamental\%20Theorem\%20of\%20Calculus/Presentation12_DefiniteIntegral\%20\&\%20Antiderivatives.pptx}{Definite Integral \& Antiderivatives (Slides 6\&7)}.  
PowerPoint file created by Professor Cynthia Roberts, UTSA.

\item \href{https://mathresearch.utsa.edu/wikiFiles/MAT1193/The\%20Fundamental\%20Theorem\%20of\%20Calculus/Presentation13_FTC\%20Part\%20I.pptx}{Fundamental Theorem of Calculus Part 1}.  
PowerPoint file created by Professor Cynthia Roberts, UTSA.

\item \href{https://mathresearch.utsa.edu/wikiFiles/MAT1214/The\%20Fundamental\%20Theorem\%20of\%20Calculus_/MAT1214-TheFundamentalTheoremPwPt.pptx}{The Fundamental Theorem of Calculus}.  
PowerPoint file created by Dr.\ Sara Shirinkam, UTSA.

\item \href{https://youtu.be/PGmVvIglZx8}{Fundamental Theorem of Calculus Part 1}.  
Video by Patrick JMT.

\item \href{https://youtu.be/yrQ6PnE6D8w}{Part 1 of the Fundamental Theorem of Calculus}.  
Video by Krista King.

\item \href{https://youtu.be/aeB5BWY0RlE}{Fundamental Theorem of Calculus Part 1}.  
Video by The Organic Chemistry Tutor.

\item \href{https://youtu.be/jyRdHbHeUuU}{The Second Fundamental Theorem of Calculus}.  
Video by James Sousa, Math is Power 4U.

\item \href{https://youtu.be/YE9jpfxEFYk}{Example 1: The Second Fundamental Theorem of Calculus}.  
Video by James Sousa, Math is Power 4U.

\item \href{https://youtu.be/gJtkTRjqM6I}{Example 2: The Second Fundamental Theorem of Calculus (Reverse Order)}.  
Video by James Sousa, Math is Power 4U.

\item \href{https://youtu.be/X5ke8bUiiQM}{Example 3: The Second Fundamental Theorem of Calculus}.  
Video by James Sousa, Math is Power 4U.

\item \href{https://youtu.be/ep52C-MUzDY}{Example 4: The Second Fundamental Theorem of Calculus with Chain Rule}.  
Video by James Sousa, Math is Power 4U.

\item \href{https://youtu.be/8IrK5JcjMvM}{Example 5: The Second Fundamental Theorem of Calculus with Chain Rule}.  
Video by James Sousa, Math is Power 4U.

\item \href{https://youtu.be/6p5NQwAeZRc}{Example 6: The Second Fundamental Theorem of Calculus with Chain Rule}.  
Video by James Sousa, Math is Power 4U.

\item \href{https://youtu.be/18qhfQ2IeLU}{Example 7: The Second Fundamental Theorem of Calculus with Chain Rule}.  
Video by James Sousa, Math is Power 4U.

\end{itemize}

\medskip

\textbf{The Fundamental Theorem of Calculus, Part 2}

\begin{itemize}

\item \href{https://youtu.be/7AteFyUCyZo}{The Fundamental Theorem of Calculus}.  
Video by James Sousa, Math is Power 4U.

\item \href{https://youtu.be/nHnZVFeQvNQ}{The Fundamental Theorem of Calculus Part 2}.  
Video by Patrick JMT.

\item \href{https://youtu.be/T-J7SkiE39Y}{Part 2 of the Fundamental Theorem of Calculus}.  
Video by Krista King.

\item \href{https://youtu.be/ns8N1UuXl4w}{The Fundamental Theorem of Calculus Part 2}.  
Video by The Organic Chemistry Tutor.

\end{itemize}

%-------------------------------------------------------------------------------
\subsubsection{Licensing}
\label{sss:lice-7352}
%-------------------------------------------------------------------------------

Content obtained and/or adapted from:

\begin{itemize}

\item \href{https://en.wikibooks.org/wiki/Calculus/Fundamental_Theorem_of_Calculus}{Fundamental Theorem of Calculus, Wikibooks: Calculus}.  
Licensed under CC BY-SA.

\end{itemize}


\subsection{Learning Objectives}

\begin{enumerate}
    \item \textbf{Understand and Apply the Fundamental Theorem of Calculus:}
    \item \textbf{Outcome:} Students will be able to explain the Fundamental Theorem of Calculus and use it to compute definite integrals.
    \item \textbf{Details:} This includes demonstrating how the theorem links differentiation and integration, and applying the theorem to solve problems involving the calculation of areas under curves.
    \item \textbf{Prove the Fundamental Theorem of Calculus:}
    \item \textbf{Outcome:} Students will be able to prove both parts of the Fundamental Theorem of Calculus.
    \item \textbf{Details:} This involves understanding the logical steps in the proofs, including the use of the Mean Value Theorem for Integration and the concept of antiderivatives.
    \item \textbf{Interpret and Apply the Mean Value Theorem for Integration:}
    \item \textbf{Outcome:} Students will be able to interpret the Mean Value Theorem for Integration and use it in the context of the Fundamental Theorem of Calculus.
    \item \textbf{Details:} This includes explaining the conditions under which the theorem holds and solving problems where the theorem is applied to find the average value of a continuous function over an interval.
\end{enumerate}

\subsection{Assessment Instrument}
\label{ssec:thef-0933}

Assessment verifies mastery of both parts of the Fundamental Theorem: computing derivatives of integral-defined functions and evaluating definite integrals via antiderivatives. Students must also connect the theorem to the Mean Value Theorem for Integration.

\textbf{Example problems.}

\begin{enumerate}
    \item Let $G(x) = \int_0^x \sin(t^2)\,dt$. 
    Find $G'(x)$ and $G'(\sqrt{\pi})$, citing the relevant part of the Fundamental Theorem.
    
    \item Evaluate 
    \[
    \int_1^4 \left(3\sqrt{x} - \frac{2}{x^2}\right)\,dx
    \]
    using Part~2 of the Fundamental Theorem, showing the antiderivative step.
    
    \item Find 
    \[
    \frac{d}{dx}\int_{x}^{x^2}\cos(t)\,dt
    \]
    by combining Part~1 of the Fundamental Theorem with the Chain Rule.
\end{enumerate}


%===============================================================================
\section{Substitution method for integrals}
\label{sec:subs-b663}
%===============================================================================


%-------------------------------------------------------------------------------
\subsubsection{Integration by Substitution}
\label{sss:inte-6921}
%-------------------------------------------------------------------------------


There is a theorem that will help you with substitution for integration.

It is called \textbf{Change of Variables for Definite Integrals}.


what the theorem looks like is this


\[
\int _{a}^{b} f(x)\operatorname {d}x = \int _{\alpha}^{\beta} f(g(u))g\prime (u)\operatorname {d}u
\]


In order to get $\alpha$ you must plug \textbf{a}

into the function \textbf{g} and to get $\beta$ you must plug

\textbf{b} into the function \textbf{g}.


The tricky part is trying to identify what you want to make your

\textbf{u} to be. Some times substitution will not

be enough and you will have to use the rules for integration by parts.

That will be covered in a different section


%-------------------------------------------------------------------------------
\subsubsection{Steps}
\label{sss:step-fca3}
%-------------------------------------------------------------------------------


%-------------------------------------------------------------------------------
\subsubsection{Steps}
\label{sss:step-fca3}
%-------------------------------------------------------------------------------


\begin{align}
\int _{x=a}^{x=b}f(x) \, dx
&= \int _{x=a}^{x=b} f(x) \frac{du}{du} \, dx \quad\quad\quad\quad\quad\quad \text{(i.e. } \frac{du}{du}=1\text{)} \\&= \int _{x=a}^{x=b} \left(f(x) \frac{dx}{du}\right) \left(\frac{du}{dx}\right) \, dx \quad\quad\quad \text{(i.e. } \frac{dx}{du} \cdot \frac{du}{dx}=1\text{)} \\&= \int _{x=a}^{x=b} \left(f(x) \frac{dx}{du}\right) g'(x) \, dx \quad\quad\quad \text{(i.e. } \frac{du}{dx}=g'(x)\text{)} \\&= \int _{x=a}^{x=b} h(g(x)) g'(x) \, dx \quad\quad\quad \text{(i.e. Now equate } \left(f(x) \frac{dx}{du}\right) \text{ with } h(g(x))\text{)} \\&= \int _{x=a}^{x=b} h(u) g'(x) \, dx \quad\quad\quad\quad \text{(i.e. } g(x)=u\text{)} \\&= \int _{u=g(a)}^{u=g(b)} h(u) \, du \quad\quad\quad\quad \text{(i.e. } du=\frac{du}{dx} \, dx=g'(x) \, dx\text{)} \\&= \int _{u=c}^{u=d} h(u) \, du \quad\quad\quad\quad \text{(i.e. We have achieved our desired result)}
\end{align}


%-------------------------------------------------------------------------------
\subsubsection{Example 1}
\label{sss:exam-acc3}
%-------------------------------------------------------------------------------


$\int _{0}^{2} x(x^2+1)^2 dx$


Instead of making this a big polynomial we will just use the

substitution method.


\textbf{Step 1}


> Identify your $u$


> Let $u = x^2+1$


\textbf{Step 2}


> Identify $du$


> $du=2x \text{ } dx$


\textbf{Step 3}


> Now we plug in our limits of integration to our $i$ to find our new limits of integration When $x = 0, u =0^2 + 1 = 1$ and when $x = 2, u = 2^2 + 1 = 5$


> Now our integration problem looks something like this


\[
\frac{1}{2} \int _{0}^{5} (x^2 + 1)^2 (2x)dx
\]


\textbf{Step 4}


> write your new integration problem


> When we plug in our \textbf{u} it looks like


\[
\frac {1}{2} \int _{0}^{5} (u)^2 \operatorname {d}u
\]


\textbf{Step 5}


> Evaluate the Integral

> \begin{align}
\frac{1}{2} \left[\frac{1}{3} u^3 \right] _{0}^{5} 
&= \frac{1}{2} \left[\left(\frac{1}{3} \cdot 5^3 \right) - \left(\frac{1}{3} \cdot 0^3 \right)\right] \\&= \frac{1}{2} \left[\frac{1}{3} \cdot 125 \right] \\&= \frac{1}{2} \left[\frac{125}{3}\right] \\&= \frac{125}{6}
\end{align}


As you can see this all simplified fairly nice. Using substitution will be hard, for most people, at first. Once you get the hang of doing this it should come to you faster and faster each time.


%-------------------------------------------------------------------------------
\subsubsection{Example 2}
\label{sss:exam-5bdd}
%-------------------------------------------------------------------------------


\[
\int 3x^2(x^3+1)^5dx
\]


we see that $3x^2$ is the derivative of $x^3+1$ . Letting


\[
u=x^3+1
\]


we have


\[
\frac{du}{dx}=3x^2
\]


or, in order to apply it to the integral,


\[
du=3x^2dx
\]


With this we may write


\[
\int 3x^2(x^3+1)^5dx=\int u^5du=\frac{u^6}{6}+C=\frac{(x^3+1)^6}{6}+C
\]


Note that it was not necessary that we had $\textit{exactly}$ the derivative of $u$ in our integrand. It would have been sufficient to have any constant multiple of the derivative.


For instance, to treat the integral


\[
\int x^4\sin(x^5)dx
\]


we may let $u=x^5$ . Then $du=5x^4dx$ and so


\[
\frac{du}{5}=x^4dx
\] the right-hand side of which is a factor of our integrand. Thus,


\[
\int x^4\sin(x^5)dx=\int\frac{\sin(u)}{5}du=-\frac{\cos(u)}{5}+C=-\frac{\cos(x^5)}{5}+C
\]


%-------------------------------------------------------------------------------
\subsubsection{Resources}
\label{sss:reso-55b5}
%-------------------------------------------------------------------------------


\begin{itemize}
    \item \href{https://en.wikibooks.org/wiki/Calculus/Integration_techniques/Recognizing_Derivatives_and_the_Substitution_Rule}{Recognizing Derivatives and Substitution Rules}, WikiBooks: Calculus
    \item \href{https://en.wikibooks.org/wiki/High_School_Calculus/Integration_by_Substitution}{Integration by Substitution}, WikiBooks: High School Calculus
    \item \href{https://www.youtube.com/watch?v=uoCW8S-I9Es}{Example 1}. Produced by Professor Zachary Sharon, UTSA
    \item \href{https://www.youtube.com/watch?v=zqMxMtjbaBE}{Example 2}. Produced by TA Catherine Sporer, UTSA
\end{itemize}


\begin{itemize}
    \item \href{https://www.youtube.com/watch?v=568hUU4beuk}{Indefinite Integration Using Substitution} by James Sousa, Math is Power 4U
    \item \href{https://youtu.be/VfvJ5C9FjcA}{Integration by Substitution Part 1} by James Sousa, Math is Power 4U
    \item \href{https://youtu.be/pmekyFVIorE}{Integration by Substitution Part 2} by James Sousa, Math is Power 4U
    \item \href{https://www.youtube.com/watch?v=3MbdmYeYB_I}{Ex 1: Indefinite Integration Using Substitution} by James Sousa, Math is Power 4U
    \item \href{https://www.youtube.com/watch?v=xLtPL1Td5f4}{Ex 2: Indefinite Integration Using Substitution} by James Sousa, Math is Power 4U
    \item \href{https://www.youtube.com/watch?v=Cdk-3H35TrY}{Ex 3: Indefinite Integration Using Substitution} by James Sousa, Math is Power 4U
    \item \href{https://www.youtube.com/watch?v=hauyUjzRTos}{Ex 4: Indefinite Integration Using Substitution} by James Sousa, Math is Power 4U
    \item \href{https://www.youtube.com/watch?v=902ta0Kshoc}{Ex 5: Indefinite Integration Using Substitution} by James Sousa, Math is Power 4U
    \item \href{https://www.youtube.com/watch?v=FJHHQdQEecE}{Ex 6: Indefinite Integration Using Substitution} by James Sousa, Math is Power 4U
    \item \href{https://www.youtube.com/watch?v=VY5r1-9apdc}{Ex 7: Indefinite Integration Using Substitution} by James Sousa, Math is Power 4U
    \item \href{https://www.youtube.com/watch?v=UnND2S4IQSA}{Ex 8: Indefinite Integration Using Substitution} by James Sousa, Math is Power 4U
    \item \href{https://www.youtube.com/watch?v=CVfe5vNy4y4}{Ex 9: Indefinite Integration Using Substitution} by James Sousa, Math is Power 4U
    \item \href{https://youtu.be/qclrs-1rpKI}{Integration using U-Substitution} by patrickJMT
    \item \href{https://youtu.be/x06V9xuLdqg}{U-Substitution - More Complicated Examples} by patrickJMT
    \item \href{https://youtu.be/HHSepCzP56M}{U-Substitution Example 1} by Krista King
    \item \href{https://youtu.be/iRIhO-w46L0}{U-Substitution Example 2} by Krista King
    \item \href{https://youtu.be/w1FHS-XV64Q}{U-Substitution Example 3} by Krista King
    \item \href{https://youtu.be/_r6xMeO7aEw}{U-Substitution Example 4} by Krista King
    \item \href{https://youtu.be/gIj3objG09Q}{U-Substitution Example 5} by Krista King
    \item \href{https://youtu.be/TlUW6ZenNXY}{U-Substitution Example 6} by Krista King
    \item \href{https://youtu.be/W1WVFJ8iMqQ}{U-Substitution Example 7} by Krista King
    \item \href{https://youtu.be/IAh00vU3FSY}{How To Integrate Using U-Substitution} by The Organic Chemistry Tutor
\end{itemize}

$\textbf{Definite Integrals Using Substitution}$


\begin{itemize}
    \item \href{https://www.youtube.com/watch?v=hjr27pv8pHQ}{Definite Integration Using Subsitution} by James Sousa, Math is Power 4U
    \item \href{https://www.youtube.com/watch?v=jUwNl_PDeDY}{Ex 1: Definite Integration Using Substitution - Change Limits of Integration?} by James Sousa, Math is Power 4U
    \item \href{https://www.youtube.com/watch?v=n3DGt7cnS70}{Ex 2: Definite Integration Using Substitution - Change Limits of Integration?} by James Sousa, Math is Power 4U
    \item \href{https://www.youtube.com/watch?v=8ofGv7TQL34}{Ex 1: Definite Integration Using Substitution} by James Sousa, Math is Power 4U
    \item \href{https://www.youtube.com/watch?v=M-IUhe-iQx0}{Ex 2: Definite Integration Using Substitution} by James Sousa, Math is Power 4U
    \item \href{https://youtu.be/3XOzg2-DdTA}{Integration by U-Substitution, Definite Integral} by patrickJMT
    \item \href{https://youtu.be/FJoyIAIC1Ag}{U-Substitution: When Do I Have to Change the Limits of Integration ?} by patrickJMT
    \item \href{https://youtu.be/0A2RlnutO8U}{U-Substitution Integration, Indefinite \& Definite Integral} by The Organic Chemistry Tutor
\end{itemize}


%-------------------------------------------------------------------------------
\subsubsection{Licensing}
\label{sss:lice-7352}
%-------------------------------------------------------------------------------


Content obtained and/or adapted from:


\begin{itemize}
    \item \href{https://en.wikibooks.org/wiki/Calculus/Integration_techniques/Recognizing_Derivatives_and_the_Substitution_Rule}{Recognizing Derivatives and Substitution Rules, WikiBooks: Calculus} under a CC BY-SA license
    \item \href{https://en.wikibooks.org/wiki/High_School_Calculus/Integration_by_Substitution}{Integration by Substitution, WikiBooks: High School Calculus} under a CC BY-SA license
\end{itemize}

\subsection{Learning Objectives}

\begin{enumerate}
    \item \textbf{Identify Appropriate Substitutions:}
    \item \textbf{Outcome:} Students will be able to recognize and select appropriate substitution variables ($u$) to simplify integrals using the method of substitution.
    \item \textbf{Details:} This involves understanding how to choose a substitution that simplifies the integral and correctly applying the change of limits in definite integrals.
    \item \textbf{Apply the Change of Variables Theorem:}
    \item \textbf{Outcome:} Students will be able to apply the Change of Variables theorem to transform and evaluate definite integrals.
    \item \textbf{Details:} This includes converting the integral bounds and integrand according to the substitution, then evaluating the transformed integral.
    \item \textbf{Solve Complex Integrals Using Substitution:}
    \item \textbf{Outcome:} Students will be able to solve complex integrals that require the method of substitution, including those that might initially seem non-trivial.
    \item \textbf{Details:} This includes integrating functions where the substitution involves polynomial, trigonometric, or exponential functions, and ensuring the correct simplification and integration process.
\end{enumerate}

\subsection{Assessment Instrument}
\label{ssec:subs-63bd}

Problems require identifying a suitable $u$-substitution, transforming the integrand and limits (for definite integrals), and verifying the result by differentiation. At least one problem involves a non-obvious substitution.

\textbf{Example problems.}

\begin{enumerate}
    \item Evaluate $\int 2x(x^2+3)^5\,dx$ using $u = x^2+3$. Show the substitution, transformed integral, and back-substitution.
    \item Compute $\int_0^1 x\,e^{x^2}\,dx$ by substitution, transforming the limits explicitly.
    \item Find $\int\dfrac{\cos x}{\sin^2 x}\,dx$ using $u = \sin x$ and verify by differentiating the result.
\end{enumerate}

%===============================================================================
\section{Area between curves}
\label{sec:area-3a2e}
%===============================================================================


%-------------------------------------------------------------------------------
\subsubsection{Introduction}
\label{sss:intr-8800}
%-------------------------------------------------------------------------------


Finding the area between two curves, usually given by two explicit

functions, is often useful in calculus.


In general the rule for finding the area between two curves is


$A=A _{\mathrm top}-A _{\mathrm bottom}$ or


If f(x) is the upper function and g(x) is the lower function


$A=\int\limits_a^b \bigl(f(x)-g(x)\bigr)dx$


This is true whether the functions are in the first quadrant or not.


%-------------------------------------------------------------------------------
\subsubsection{Area between two curves}
\label{sss:area-769e}
%-------------------------------------------------------------------------------


Suppose we are given two functions $y_1=f(x)$ and $y_2=g(x)$ and we want

to find the area between them on the interval $[a,b]$ . Also assume that

$f(x)\ge g(x)$ for all $x$ on the interval $[a,b]$ . Begin by

partitioning the interval $[a,b]$ into $n$ equal subintervals each

having a length of $\Delta x=\frac{b-a}{n}$ . Next choose any point in

each subinterval, $x_i^{*}$ . Now we can 'create' rectangles on each

interval. At the point $x_i\*$ , the height of each rectangle is

$f(x_i^{*})-g(x_i^{*})$ and the width is $\Delta x$ . Thus the area of each

rectangle is $\bigl[f(x_i^{*})-g(x_i^{*})\bigr]\Delta x$ . An

\textit{approximation} of the area, $A$ , between the two curves is


\[
A:=\sum _{i=1}^n \Big[f(x_i^{*})-g(x_i^{*})\Big]\Delta x
\] Now we take the limit as $n$ approaches infinity and get


\[
A=\lim _{n\to\infty}\sum _{i=1}^n \Big[f(x_i^{*})-g(x_i^{*})\Big]\Delta x
\]

which gives the exact area. Recalling the definition of the definite

integral we notice that


\[
A=\int\limits_a^b \bigl(f(x)-g(x)\bigr)dx
\]  This formula of finding

the area between two curves is sometimes known as applying integration

with respect to the \textit{x}-axis since the rectangles used to approximate

the area have their bases lying parallel to the \textit{x}-axis. It will be

most useful when the two functions are of the form $y_1=f(x)$ and

$y_2=g(x)$ . Sometimes however, one may find it simpler to integrate

with respect to the \textit{y}-axis. This occurs when integrating with respect

to the \textit{x}-axis would result in more than one integral to be evaluated.

These functions take the form $x_1=f(y)$ and $x_2=g(y)$ on the interval

$[c,d]$ . Note that $[c,d]$ are values of $y$ . The derivation of this

case is completely identical. Similar to before, we will assume that

$f(y)\ge g(y)$ for all $y$ on $[c,d]$ . Now, as before we can divide the

interval into $n$ subintervals and create rectangles to approximate the

area between $f(y)$ and $g(y)$ . It may be useful to picture each

rectangle having their 'width', $\Delta y$ , parallel to the \textit{y}-axis

and 'height', $f(y_i^{*})-g(y_i^{*})$ at the point $y_i^{*}$, parallel to

the \textit{x}-axis. Following from the work above we may reason that an

\textit{approximation} of the area, $A$ , between the two curves is


\[
A:=\sum _{i=1}^n \Big[f(y_i^{*})-g(y_i^{*})\Big]\Delta y
\] 

As before, we take the limit as $n$ approaches infinity to arrive at


\[
A=\lim _{n\to\infty}\sum _{i=1}^n \Big[f(y_i^{*})-g(y_i^{*})\Big]\Delta y
\]

which is nothing more than a definite integral, so


\[
A=\int\limits_c^d \bigl(f(y)-g(y)\bigr)dy
\] Regardless of the form of the functions, we basically use the same formula.


% [Image: Closed_path_integral_defined]


Figure: Closed Path Integral


%-------------------------------------------------------------------------------
\subsubsection{Resources}
\label{sss:reso-55b5}
%-------------------------------------------------------------------------------

\begin{itemize}
\item \href{https://en.wikibooks.org/wiki/Calculus/Area}{Area}.  
Wikibooks: Calculus.
\end{itemize}

%-------------------------------------------------------------------------------
\subsubsection{Videos}
\label{sss:vide-cce3}
%-------------------------------------------------------------------------------

\begin{itemize}

\item \href{https://www.youtube.com/watch?v=de-0e0nYf5I}{Determining Area Between Two Curves}.  
Video by James Sousa, Math is Power 4U.

\item \href{https://www.youtube.com/watch?v=E_1aDcOoDtE}{The Area Between Two Graphs}.  
Video by James Sousa, Math is Power 4U.

\item \href{https://www.youtube.com/watch?v=q7De3oOp4Ug}{Example 1: Area Between a Linear and Quadratic Function (with respect to $x$)}.  
Video by James Sousa, Math is Power 4U.

\item \href{https://www.youtube.com/watch?v=n2JqTztMrvQ}{Example 2: Area Between a Linear and Exponential Function (with respect to $x$)}.  
Video by James Sousa, Math is Power 4U.

\item \href{https://www.youtube.com/watch?v=OqMbRAEQ-cg}{Example 3: Area Between Two Exponential Functions (with respect to $x$)}.  
Video by James Sousa, Math is Power 4U.

\item \href{https://www.youtube.com/watch?v=58Esio4q4jk}{Example 4: Area Between Two Quadratic Functions (with respect to $x$)}.  
Video by James Sousa, Math is Power 4U.

\item \href{https://www.youtube.com/watch?v=CPhUolqwLNk}{Example 1: Area Bounded by Two Functions}.  
Video by James Sousa, Math is Power 4U.

\item \href{https://www.youtube.com/watch?v=lskEXsJaJV0}{Example 2: Area Bounded by Two Functions (Two Regions)}.  
Video by James Sousa, Math is Power 4U.

\item \href{https://www.youtube.com/watch?v=BH-V6Gx2n8E}{Example 3: Area Bounded by Two Trigonometric Functions}.  
Video by James Sousa, Math is Power 4U.

\item \href{https://www.youtube.com/watch?v=N4fEuEQLG5c}{Determine a Function Given the Area Between Two Functions}.  
Video by James Sousa, Math is Power 4U.

\item \href{https://youtu.be/DRFyNHdVgUA}{Finding Areas Between Curves}.  
Video by Patrick JMT.

\item \href{https://youtu.be/1yxSA-FZ0gk}{Areas Between Curves – Two Regions}.  
Video by Patrick JMT.

\item \href{https://youtu.be/iAvF1Rn7How}{Finding Area by Integrating with Respect to $y$}.  
Video by Patrick JMT.

\item \href{https://youtu.be/70NQ3ISYihw}{Area Between Curves – Integrating with Respect to $y$}.  
Video by Patrick JMT.

\item \href{https://youtu.be/Xne6Hv9useE}{Area Between Curves – Integrating with Respect to $y$ (Part 2)}.  
Video by Patrick JMT.

\item \href{https://www.youtube.com/watch?v=M2G68iQzYQE}{Finding Area Between Curves: Regions and Functions Given}.  
Video by Patrick JMT.

\item \href{https://youtu.be/8w9l99W8ucY}{Area of a Region Bounded by Three Curves}.  
Video by Patrick JMT.

\item \href{https://www.youtube.com/watch?v=2j38_za6ztc}{Area Between Curves – Which Curve is Which?}.  
Video by Krista King.

\item \href{https://youtu.be/AeInlCa-1Q8}{Area Between Curves – $dx$}.  
Video by Krista King.

\item \href{https://youtu.be/fGvSy6U209Y}{Area Between Curves – $dy$}.  
Video by Krista King.

\item \href{https://youtu.be/7lrn-8FiTW0}{Area Between Curves – Sketching}.  
Video by Krista King.

\item \href{https://www.youtube.com/watch?v=E2ArQGWX8J4}{Area Between Curves Example 1}.  
Video by Krista King.

\item \href{https://www.youtube.com/watch?v=UV-fbgHa85s}{Area Between Curves Example 2}.  
Video by Krista King.

\item \href{https://www.youtube.com/watch?v=86ivXd4-BsU}{Area Between Curves Example 3}.  
Video by Krista King.

\item \href{https://youtu.be/ltxdcegn8xc}{Area Between Two Curves}.  
Video by The Organic Chemistry Tutor.

\item \href{https://youtu.be/wXMr2YYDMiM}{Area Between Two Curves and Under a Curve (with respect to $x$ and $y$)}.  
Video by The Organic Chemistry Tutor.

\end{itemize}

%-------------------------------------------------------------------------------
\subsubsection{Licensing}
\label{sss:lice-7352}
%-------------------------------------------------------------------------------

Content obtained and/or adapted from:

\begin{itemize}
\item \href{https://en.wikibooks.org/wiki/Calculus/Area}{Area, Wikibooks: Calculus}.  
Licensed under CC BY-SA.
\end{itemize}


\subsection{Learning Objectives}

\begin{enumerate}
    \item \textbf{Understand the Concept of Area Between Two Curves:}
    \item \textbf{Outcome:} Students will be able to comprehend the mathematical principle behind finding the area between two curves by recognizing the difference between the upper and lower functions.
    \item \textbf{Details:} This includes identifying when $f(x) \geq g(x)$ on a given interval and understanding how the integral represents the area between these curves over the interval $[a, b]$.
    \item \textbf{Apply Integration to Calculate Area:}
    \item \textbf{Outcome:} Students will be able to calculate the area between two curves by applying the definite integral.
    \item \textbf{Details:} This involves using the formula $A = \int_a^b (f(x) - g(x)) \, dx$ for functions defined with respect to the x-axis, and $A = \int_c^d (f(y) - g(y)) \, dy$ for functions defined with respect to the y-axis.
    \item \textbf{Partition Intervals and Approximate Area:}
    \item \textbf{Outcome:} Students will be able to approximate the area between two curves by partitioning the interval into subintervals and using Riemann sums.
    \item \textbf{Details:} This includes understanding the process of dividing the interval $[a, b]$ or $[c, d]$ into $n$ subintervals, selecting points within each subinterval, and summing the areas of rectangles to approximate the total area between the curves.
\end{enumerate}

\subsection{Assessment Instrument}
\label{ssec:area-1946}

Students set up area integrals by identifying the upper and lower (or right and left) function on each sub-interval, finding intersection points when needed, and choosing the appropriate axis of integration.

\textbf{Example problems.}

\begin{enumerate}
    \item Find the area enclosed between $f(x)=x^2$ and $g(x)=2x$. Sketch the region, identify intersection points, and evaluate the integral.
    \item Compute the area between $y = \sin x$ and $y = \cos x$ on $[0, \pi]$. Identify where the curves cross on this interval.
    \item Find the area of the region bounded by $y^2 = x$ and $x - y = 2$ by integrating with respect to $y$.
\end{enumerate}

%===============================================================================
\section{Determining volumes by slicing}
\label{sec:dete-b39e}
%===============================================================================


%-------------------------------------------------------------------------------
\subsubsection{Volumes by Slices}
\label{sss:volu-7c52}
%-------------------------------------------------------------------------------


When we think about volume from an intuitive point of view, we typically

think of it as the amount of ``space'' an item occupies. Unfortunately

assigning a number that measures this amount of space can prove

difficult for all but the simplest geometric shapes. Calculus provides a

new tool that can greatly extend our ability to calculate volume. In

order to understand the ideas involved it helps to think about the

volume of a cylinder.


The volume of a cylinder is calculated using the formula $V=\pi r^2h$ .

The base of the cylinder is a circle whose area is given by $A=\pi r^2$

. Notice that the volume of a cylinder is derived by taking the area of

its base and multiplying by the height $h$ . For more complicated

shapes, we could think of approximating the volume by taking the area of

some cross section at some height $x$ and multiplying by some small

change in height $\Delta x$ then adding up the heights of all of these

approximations from the bottom to the top of the object. This would

appear to be a Riemann sum. Keeping this in mind, we can develop a more

general formula for the volume of solids in $\R^3$ (3 dimensional

space).


%-------------------------------------------------------------------------------
\subsubsection{Formal Definition}
\label{sss:form-5603}
%-------------------------------------------------------------------------------


Formally the ideas above suggest that we can calculate the volume of a

solid by calculating the integral of the cross-sectional area along some

dimension. In the above example of a cylinder, every cross section is

given by the same circle, so the cross-sectional area is therefore a

constant function, and the dimension of integration was vertical

(although it could have been any one we desired). Generally, if $S$ is a

solid that lies in $\R^3$ between $x=a$ and $x=b$ , let $A(x)$ denote

the area of a cross section taken in the plane perpendicular to the

$x$-axis, and passing through the point $x$ .


If the function $A(x)$ is continuous on $[a,b]$ , then the volume $V_S$

of the solid $S$ is given by:


\[
V_S = \int_a^b A(x) dx
\]


%-------------------------------------------------------------------------------
\subsubsection{Examples}
\label{sss:exam-bfeb}
%-------------------------------------------------------------------------------


%-------------------------------------------------------------------------------
\subsubsection{Example 1: A right cylinder}
\label{sss:exam-7dca}
%-------------------------------------------------------------------------------


% [Image: Figure

1]


Now we will calculate the volume of a right cylinder using our new ideas
about how to calculate volume. Since we already know the formula for the
volume of a cylinder this will give us a ``sanity check'' that our
formulas make sense. First, we choose a dimension along which to
integrate. In this case, it will greatly simplify the calculations to
integrate along the height of the cylinder, so this is the direction we
will choose. Thus we will call the vertical direction $x$ . Now
we find the function, $A(x)$ , which will describe the cross-sectional
area of our cylinder at a height of $x$ . The cross-sectional area of a
cylinder is simply a circle. Now simply recall that the area of a circle
is $\pi r^2$ , and so $A(x)=\pi r^2$ . Before performing the
computation, we must choose our bounds of integration. In this case, we
simply define $x=0$ to be the base of the cylinder, and so we will integrate from $x=0$ to $x=h$ , where $h$ is the height of the cylinder.


Finally, we integrate:
\begin{align} V _{\text{cylinder}} &= \int\limits_a^b A(x)\,dx \\ &= \int\limits_0^h \pi r^2\,dx \\ &= \pi r^2 \int\limits_0^h dx \\ &= \pi r^2 x \bigg| _{x=0}^h \\ &= \pi r^2 (h - 0) \\ &= \pi r^2 h \end{align} 
This is exactly the familiar formula for the volume of a cylinder.


%-------------------------------------------------------------------------------
\subsubsection{Example 2: A right circular cone}
\label{sss:exam-c28b}
%-------------------------------------------------------------------------------
For our next example we will look at an example where the cross
sectional area is not constant. Consider a right circular cone. Once
again the cross sections are simply circles. But now the radius varies
from the base of the cone to the tip. Once again we choose $x$ to be the
vertical direction, with the base at $x=0$ and the tip at $x=h$ , and we
will let $R$ denote the radius of the base. While we know the cross
sections are just circles we cannot calculate the area of the cross
sections unless we find some way to determine the radius of the circle
at height $x$ .

Luckily in this case it is possible to use some of what we know from
geometry. We can imagine cutting the cone perpendicular to the base
through some diameter of the circle all the way to the tip of the cone.
If we then look at the flat side we just created, we will see simply a
triangle, whose geometry we understand well. The right triangle from the
tip to the base at height $x$ is similar to the right triangle from the
tip to the base at height $h$ . This tells us that
$\frac{r}{h-x}=\frac{R}{h}$ . So that we see that the radius of the
circle at height $x$ is $r(x)=\frac{R}{h}(h-x)$ . Now using the familiar
formula for the area of a circle we see that
$A(x)=\pi\frac{R^2}{h^2}(h-x)^2$ .

Now we are ready to integrate.
\begin{align} V _{\text{cone}} &= \int\limits_a^b A(x)\,dx \\ &= \int\limits_0^h \pi\frac{R^2}{h^2}(h-x)^2\,dx \\ &= \pi\frac{R^2}{h^2}\int\limits_0^h (h-x)^2\,dx \end{align} 
By u-substitution we may let $u=h-x$ , then $du=-dx$ and our integral becomes
\begin{align} &= \pi\frac{R^2}{h^2}\left(-\int\limits_h^0 u^2\,du\right) \\ &= \pi\frac{R^2}{h^2}\left(-\frac{u^3}{3}\bigg| _h^0\right) \\ &= \pi\frac{R^2}{h^2}\left(-0 + \frac{h^3}{3}\right) \\ &= \frac{\pi}{3} R^2 h \end{align} 


%-------------------------------------------------------------------------------
\subsubsection{Example 3: A sphere}
\label{sss:exam-8e36}
%-------------------------------------------------------------------------------
In a similar fashion, we can use our definition to prove the well known
formula for the volume of a sphere. First, we must find our
cross-sectional area function, $A(x)$ . Consider a sphere of radius $R$
which is centered at the origin in $\R^3$ . If we again integrate
vertically then $x$ will vary from $-R$ to $R$ . In order to find the
area of a particular cross section it helps to draw a right triangle
whose points lie at the center of the sphere, the center of the circular
cross section, and at a point along the circumference of the cross
section. As shown in the diagram the side lengths of this triangle will
be $R$ , $|x|$ , and $r$ . Where $r$ is the radius of the circular cross
section. Then by the Pythagorean theorem $r=\sqrt{R^2-|x|^2}$ and find
that $A(x)=\pi(R^2-|x|^2)$ . It is slightly helpful to notice that
$|x|^2=x^2$ so we do not need to keep the absolute value.

So we have that
\begin{align} V _{\text{sphere}} &= \int\limits_a^b A(x)\,dx \\ &= \int\limits _{-R}^R \pi(R^2-x^2)\,dx \\ &= \pi \int\limits _{-R}^R R^2\,dx - \pi \int\limits _{-R}^R x^2\,dx \\ &= \pi R^2 x \Bigg| _{-R}^R - \pi \frac{x^3}{3} \Bigg| _{-R}^R \\ &= \pi R^2 \left(R - (-R)\right) - \pi \left(\frac{R^3}{3} - \frac{(-R)^3}{3}\right) \\ &= 2\pi R^3 - \frac{2\pi}{3} R^3 \\ &= \frac{4\pi}{3} R^3 \end{align} 


%-------------------------------------------------------------------------------
\subsubsection{Volume of Solids of Revolution}
\label{sss:volu-d284}
%-------------------------------------------------------------------------------
In this section we cover \textbf{solids of revolution} and how to calculate
their volume. A solid of revolution is a solid formed by revolving a
2-dimensional region around an axis. For example, revolving the semi-circular region bounded by the curve $y=\sqrt{1-x^2}$ and the line $y=0$ around the $x$-axis produces a sphere. There are two main methods of calculating the volume of a solid of revolution using calculus: the disk method and the shell method.


%-------------------------------------------------------------------------------
\subsubsection{Disk Method}
\label{sss:disk-f7dd}
%-------------------------------------------------------------------------------

Consider the solid formed by revolving the region bounded by the curve 
$y=f(x)$, which is continuous on $[a,b]$, and the lines $x=a$, $x=b$, and $y=0$ 
about the $x$-axis.

\medskip

We approximate the volume by replacing $f(x)$ with a step function. When the 
region is revolved, each step sweeps out a cylinder (disk), whose volume is
\[
V_{\text{cylinder}} = \pi r^2 h,
\]
where $r$ is the radius and $h$ is the thickness.

\medskip

Partition $[a,b]$ into $n$ subintervals of width
\[
\Delta x = \frac{b-a}{n}.
\]
Using a right-endpoint approximation, the sample points are
\[
x_k = a + k\Delta x.
\]

The volume of the $k$-th disk is
\[
V_k = \pi f(x_k)^2 \,\Delta x.
\]

Summing over all disks:
\[
V_{\text{approx}} = \sum_{k=1}^{n} \pi f(x_k)^2 \,\Delta x.
\]

Taking the limit as $n \to \infty$ gives the exact volume:
\[
V = \lim_{n\to\infty} \sum_{k=1}^{n} \pi f(x_k)^2 \,\Delta x 
= \int_a^b \pi f(x)^2 \, dx.
\]

\medskip

\textbf{Example: Volume of a Sphere}

Consider the semicircle
\[
f(x) = \sqrt{r^2 - x^2}
\]
on $[-r,r]$. Revolving this about the $x$-axis generates a sphere.

\begin{align}
V_{\text{sphere}} 
&= \int_{-r}^{r} \pi (r^2 - x^2)\,dx \\
&= \pi\left(\int_{-r}^{r} r^2\,dx - \int_{-r}^{r} x^2\,dx\right) \\
&= \pi\left(r^2 x \Big|_{-r}^{r} - \frac{x^3}{3} \Big|_{-r}^{r}\right) \\
&= \pi\left(2r^3 - \frac{2r^3}{3}\right) \\
&= \frac{4\pi}{3} r^3.
\end{align}

%-------------------------------------------------------------------------------
\subsubsection{Washer Method}
\label{sss:wash-235c}
%-------------------------------------------------------------------------------

The washer method is an extension of the disk method to solids of
revolution formed by revolving an area bounded between two curves around
the $x$-axis. Consider the solid of revolution formed by revolving the
region in \href{:File:Generating_fxs_washer_method.svg "wikilink"}{figure 3}
around the $x$-axis. The curve $f(x)$ is the same as that in [figure
1], but now our solid
has an irregularly shaped hole through its center whose volume is that
of the solid formed by revolving the curve $g(x)$ around the $x$-axis.
Our approximating region has the same upper boundary, $f _{\mathrm step}(x)$
as in [figure 2],
but now we extend only down to $g _{\mathrm step}(x)$ rather than all the way
down to the $x$-axis. Revolving each block around the $x$-axis forms a
washer-shaped solid with outer radius $f_ {\mathrm step}(x)$ and inner radius
$g_ {\mathrm step}(x)$ . The volume of the $k$-th hollow cylinder is
\begin{align}
V_k &= \pi\cdot f(x_k)^2 \Delta x - \pi\cdot g(x_k)^2 \Delta x \\&= \pi\bigl(f(x_k)^2 - g(x_k)^2\bigr)\Delta x
\end{align}
where 
\[
\Delta x = \frac{b-a}{n} \quad \text{and} \quad x_k = k\Delta x.
\]
The volume of the entire approximating solid is
\[
V _{\mathrm approx}=\sum _{k=1}^n \pi\bigl(f(x_k)^2-g(x_k)^2\bigr)\Delta x
\]
Taking the limit as $n$ approaches infinity gives the volume
\begin{align} V &= \lim_{n\to\infty}\sum _{k=1}^n \pi\bigl(f(x_k)^2 - g(x_k)^2\bigr)\Delta x \\ &= \int\limits_a^b \pi\bigl(f(x)^2 - g(x)^2\bigr)\,dx \end{align} 

%-------------------------------------------------------------------------------
\subsubsection{Resources}
\label{sss:reso-55b5}
%-------------------------------------------------------------------------------

%-------------------------------------------------------------------------------
\subsubsection{Volume by Slices}
\label{sss:volu-9ae4}
%-------------------------------------------------------------------------------

\begin{itemize}

\item \href{https://youtu.be/Q8c8Gs1OUgk}{Determine Volume of Solids by Slices}.  
Video by James Sousa, Math is Power 4U.

\item \href{https://youtu.be/_QhW_pXy0wE}{Example 1: Volume of a Solid with Known Cross Section Using Integration (Volume by Slices)}.  
Video by James Sousa, Math is Power 4U.

\item \href{https://youtu.be/ODNdwqJ7TAg}{Example 2: Volume of a Solid with Known Cross Section Using Integration (Volume by Slices)}.  
Video by James Sousa, Math is Power 4U.

\item \href{https://youtu.be/LvBZ0jnAMbo}{Example 3: Volume of a Solid with Known Cross Section Using Integration (Volume by Slices)}.  
Video by James Sousa, Math is Power 4U.

\item \href{https://youtu.be/BQ5sg21w-lA}{Example 4: Volume of a Solid with Known Cross Section Using Integration (Volume by Slices)}.  
Video by James Sousa, Math is Power 4U.

\item \href{https://youtu.be/2OgAluEqd8s}{Example 5: Volume of a Solid with Known Cross Section Using Integration (Volume by Slices)}.  
Video by James Sousa, Math is Power 4U.

\item \href{https://www.youtube.com/watch?v=4vLy5VoUcQE}{Volume with Cross Sections: Introduction}.  
Video by Khan Academy.

\item \href{https://www.youtube.com/watch?v=tf4C8x8e7HQ}{Volume with Cross Sections: Semicircle}.  
Video by Khan Academy.

\item \href{https://www.youtube.com/watch?v=S-agS4YaxxU}{Volume with Cross Sections: Triangle}.  
Video by Khan Academy.

\item \href{https://youtu.be/XasIRTb5WH4}{Volumes Using Cross Sections (Squares, Semicircles, Rectangles, Equilateral Triangles)}.  
Video by The Organic Chemistry Tutor.

\end{itemize}

%-------------------------------------------------------------------------------
\subsubsection{The Disk Method}
\label{sss:thed-913f}
%-------------------------------------------------------------------------------

\begin{itemize}

\item \href{https://www.youtube.com/watch?v=1CbZlM09zF8}{Volume of Revolution – The Disk Method}.  
Video by James Sousa, Math is Power 4U.

\item \href{https://www.youtube.com/watch?v=oYe4gkKQDE8}{Example 1: Volume of Revolution – Disk Method}.  
Video by James Sousa, Math is Power 4U.

\item \href{https://www.youtube.com/watch?v=e7qNtcWpg2s}{Example 2: Volume of Revolution – Disk Method}.  
Video by James Sousa, Math is Power 4U.

\item \href{https://www.youtube.com/watch?v=Ba_q-eu8n34}{Example 3: Volume of Revolution – Disk Method}.  
Video by James Sousa, Math is Power 4U.

\item \href{https://www.youtube.com/watch?v=yTElv1k33rk}{Example 4: Volume of Revolution Using the Disk Method}.  
Video by James Sousa, Math is Power 4U.

\item \href{https://www.youtube.com/watch?v=xoi-6CyL6cg}{Example 5: Volume of Revolution Using the Disk Method}.  
Video by James Sousa, Math is Power 4U.

\item \href{https://www.youtube.com/watch?v=2tpyuR-GypA}{Example 6: Volume of Revolution Using the Disk Method}.  
Video by James Sousa, Math is Power 4U.

\item \href{https://www.youtube.com/watch?v=RPpih8rD-5k}{Volume of Rotation: Disk Method About the $x$-axis or $y=c$}.  
Video by Krista King.

\item \href{https://www.youtube.com/watch?v=DNmHo68wYFY}{Volume of Rotation: Disk Method About the $y$-axis or $x=c$}.  
Video by Krista King.

\item \href{https://www.youtube.com/watch?v=btGaOTXxXs8}{Disk Method Around the $x$-axis}.  
Video by Khan Academy.

\item \href{https://www.youtube.com/watch?v=43AS7bPUORc}{Disk Method Around the $y$-axis}.  
Video by Khan Academy.

\item \href{https://www.youtube.com/watch?v=XdzcU5JbVcA}{Disk Method Around a Horizontal Line}.  
Video by Khan Academy.

\item \href{https://www.youtube.com/watch?v=jxf7XqvZWWg}{Disk Method Around a Vertical Line}.  
Video by Khan Academy.

\end{itemize}

%-------------------------------------------------------------------------------
\subsubsection{The Washer Method}
\label{sss:thew-7078}
%-------------------------------------------------------------------------------

\begin{itemize}

\item \href{https://www.youtube.com/watch?v=rbqWHbxmVUI}{Volume of Revolution – Washer Method about the $x$-axis}.  
Video by James Sousa, Math is Power 4U.

\item \href{https://www.youtube.com/watch?v=8gbphumzbSI}{Volume of Revolution – Washer Method about the $y$-axis}.  
Video by James Sousa, Math is Power 4U.

\item \href{https://www.youtube.com/watch?v=_FF3CM5MNe8}{Volume of Revolution – Washer Method about a general axis}.  
Video by James Sousa, Math is Power 4U.

\item \href{https://www.youtube.com/watch?v=6zaqRVD90mY}{Example: Washer Method about $y=3$}.  
Video by James Sousa, Math is Power 4U.

\item \href{https://www.youtube.com/watch?v=kivhB247H-c}{Example: Washer Method about $y=3$}.  
Video by James Sousa, Math is Power 4U.

\item \href{https://www.youtube.com/watch?v=IWoEivm5AqM}{Example: Washer Method about the $y$-axis}.  
Video by James Sousa, Math is Power 4U.

\item \href{https://www.youtube.com/watch?v=kkVVOStqRmI}{Example: Washer Method about the $y$-axis}.  
Video by James Sousa, Math is Power 4U.

\item \href{https://www.youtube.com/watch?v=2yd_Ae-7wls}{Example: Washer Method about $x=5$}.  
Video by James Sousa, Math is Power 4U.

\item \href{https://www.youtube.com/watch?v=zHi1qtf4pGg}{Volume of Rotation: Washer Method}.  
Video by Krista King.

\item \href{https://www.youtube.com/watch?v=Sg24zkhG0Vw}{Volume of Rotation: Washer Method}.  
Video by Krista King.

\item \href{https://www.youtube.com/watch?v=Thvc2s9aUP4}{The Washer Method}.  
Video by Khan Academy.

\item \href{https://www.youtube.com/watch?v=OFNGpKGg9IQ}{Washer Method Around a Horizontal Line (Part 1)}.  
Video by Khan Academy.

\item \href{https://www.youtube.com/watch?v=LKzpw_HUKNQ}{Washer Method Around a Horizontal Line (Part 2)}.  
Video by Khan Academy.

\item \href{https://www.youtube.com/watch?v=WAPZihVUmzE}{Washer Method Around a Vertical Line (Part 1)}.  
Video by Khan Academy.

\item \href{https://www.youtube.com/watch?v=i-Rb4_n929k}{Washer Method Around a Vertical Line (Part 2)}.  
Video by Khan Academy.

\end{itemize}

%-------------------------------------------------------------------------------
\subsubsection{Disk/Washer Methods}
\label{sss:thed-6219}
%-------------------------------------------------------------------------------

\begin{itemize}

\item \href{https://www.youtube.com/watch?v=E5OOMbz5jZk}{Volumes of Revolution – Disk/Washer Method Example 1}.  
Video by Patrick JMT.

\item \href{https://www.youtube.com/watch?v=ithgZfRKMHI}{Volumes of Revolution – Disk/Washer Method Example 2}.  
Video by Patrick JMT.

\item \href{https://www.youtube.com/watch?v=lE814ngFt3I}{Volumes of Revolution – Disk/Washer Method Example 3}.  
Video by Patrick JMT.

\end{itemize}

%-------------------------------------------------------------------------------
\subsubsection{Licensing}
\label{sss:lice-7352}
%-------------------------------------------------------------------------------

Content obtained and/or adapted from:

\begin{itemize}

\item \href{https://en.wikibooks.org/wiki/Calculus/Volume}{Volume, Wikibooks: Calculus}.  
Licensed under CC BY-SA.

\item \href{https://en.wikibooks.org/wiki/Calculus/Volume_of_solids_of_revolution}{Volume of solids of revolution, Wikibooks: Calculus}.  
Licensed under CC BY-SA.

\end{itemize}

\subsection{Learning Objectives}

\begin{enumerate}
    \item \textbf{Calculate Volumes of Solids Using Cross-Sectional Areas:}
    \item \textbf{Outcome:} Students will be able to calculate the volume of a solid by integrating the area of cross-sections perpendicular to a given axis.
    \item \textbf{Details:} This includes identifying the cross-sectional area function \( A(x) \), setting up the appropriate integral, and evaluating it to find the volume of solids such as cylinders, cones, and spheres.
    \item \textbf{Apply the Disk Method to Compute Volumes of Solids of Revolution:}
    \item \textbf{Outcome:} Students will be able to use the disk method to find the volume of a solid of revolution formed by revolving a 2-dimensional region around an axis.
    \item \textbf{Details:} This involves identifying the function \( f(x) \) describing the radius of the disks, setting up the integral \( V = \int_a^b \pi [f(x)]^2 dx \), and solving it to determine the volume.
    \item \textbf{Use the Washer Method to Determine Volumes of Solids with Holes:}
    \item \textbf{Outcome:} Students will be able to apply the washer method to calculate the volume of a solid of revolution with an inner radius, resulting in a hollow shape.
    \item \textbf{Details:} This includes identifying the outer function \( f(x) \) and inner function \( g(x) \), setting up the integral \( V = \int_a^b \pi [f(x)^2 - g(x)^2] dx \), and performing the integration to find the volume of solids such as those with an irregular hole.
\end{enumerate}

\subsection{Assessment Instrument}
\label{ssec:dete-251d}

Assessment targets the disk/washer method via cross-sectional area. Students must identify the area function $A(x)$ or $A(y)$, set up the integral, and interpret the geometry of each slice.

\textbf{Example problems.}

\begin{enumerate}
    \item Find the volume of the solid whose cross sections perpendicular to the $x$-axis are squares with side length $f(x) = \sqrt{x}$ on $[0,4]$.
    \item Use the disk method to find the volume of the solid obtained by revolving $y = x^2$ on $[0,2]$ about the $x$-axis.
    \item Use the washer method to find the volume of the solid formed by revolving the region between $y=\sqrt{x}$ and $y=x$ about the $x$-axis.
\end{enumerate}

%===============================================================================
\section{Volumes of Revolution and Shells}
\label{sec:volu-a692}
%===============================================================================

In this section we cover \textbf{solids of revolution} and how to calculate

their volume. A solid of revolution is a solid formed by revolving a

2-dimensional region around an axis. For example, revolving the

semi-circular region bounded by the curve $y=\sqrt{1-x^2}$ and the line

$y=0$ around the $x$-axis produces a sphere. There are two main methods

of calculating the volume of a solid of revolution using calculus: the

disk method and the shell method.


%-------------------------------------------------------------------------------
\subsubsection{Disk Method}
\label{sss:disk-f7dd}
%-------------------------------------------------------------------------------

Consider the solid formed by revolving the region bounded by the curve 
$y=f(x)$, which is continuous on $[a,b]$, and the lines $x=a$, $x=b$, and $y=0$ 
about the $x$-axis.

\medskip

We can approximate the volume by replacing $f(x)$ with a stepwise function 
$g(x)$ (see Figure~2), using a right-endpoint approximation. When the region 
is revolved, each step generates a cylinder (disk), whose volume is
\[
V_{\text{cylinder}} = \pi r^2 h,
\]
where $r$ is the radius and $h$ is the height (thickness).

\medskip

Partition the interval $[a,b]$ into $n$ subintervals of equal width
\[
\Delta x = \frac{b-a}{n}.
\]
Using a right-endpoint approximation, the sample points are
\[
x_k = a + k\Delta x.
\]

The volume of the $k$-th disk is
\[
V_k = \pi f(x_k)^2\,\Delta x.
\]

Summing over all subintervals gives the approximate volume
\[
V_{\text{approx}} = \sum_{k=1}^{n} \pi f(x_k)^2\,\Delta x.
\]

Taking the limit as $n \to \infty$ yields the exact volume
\[
V = \lim_{n\to\infty} \sum_{k=1}^{n} \pi f(x_k)^2\,\Delta x
= \int_a^b \pi f(x)^2\,dx.
\]

\medskip

\textbf{Example: Volume of a Sphere}

Let us compute the volume of a sphere using the disk method. The generating 
region is bounded by the curve
\[
f(x) = \sqrt{r^2 - x^2}
\]
and the line $y=0$, with limits of integration $x = -r$ to $x = r$.

\begin{align}
V_{\text{sphere}} 
&= \int_{-r}^{r} \pi (r^2 - x^2)\,dx \\
&= \pi\left(\int_{-r}^{r} r^2\,dx - \int_{-r}^{r} x^2\,dx\right) \\
&= \pi\left(r^2 x \Big|_{-r}^{r} - \frac{x^3}{3} \Big|_{-r}^{r}\right) \\
&= \pi\left(2r^3 - \frac{2r^3}{3}\right) \\
&= \frac{4\pi}{3} r^3.
\end{align}

%-------------------------------------------------------------------------------
\subsubsection{Exercises}
\label{sss:exer-disk}
%-------------------------------------------------------------------------------

\begin{enumerate}
\item Calculate the volume of the cone with radius $r$ and height $h$ 
generated by revolving the region bounded by
\[
y = r - \frac{r}{h}x,
\]
and the lines $y=0$ and $x=0$ about the $x$-axis.

\item Calculate the volume of the solid of revolution generated by 
revolving the region bounded by the curve $y = x^2$ and the lines 
$x=1$ and $y=0$ about the $x$-axis.
\end{enumerate}

%-------------------------------------------------------------------------------
\subsubsection{Washer Method}
\label{sss:wash-235c}
%-------------------------------------------------------------------------------

The washer method extends the disk method to solids formed by revolving 
a region bounded between two curves about the $x$-axis.

\medskip

Consider the region bounded by $y = f(x)$ (outer curve) and 
$y = g(x)$ (inner curve), where $f(x) \ge g(x) \ge 0$ on $[a,b]$.

\medskip

Partition $[a,b]$ into $n$ subintervals of width
\[
\Delta x = \frac{b-a}{n},
\]
and use right endpoints:
\[
x_k = a + k\Delta x.
\]

Each slice forms a washer with outer radius $f(x_k)$ and inner radius $g(x_k)$.

The volume of the $k$-th washer is
\begin{align}
V_k 
&= \pi f(x_k)^2 \Delta x - \pi g(x_k)^2 \Delta x \\
&= \pi\bigl(f(x_k)^2 - g(x_k)^2\bigr)\Delta x.
\end{align}

Summing over all washers:
\[
V_{\text{approx}} = \sum_{k=1}^{n} \pi\bigl(f(x_k)^2 - g(x_k)^2\bigr)\Delta x.
\]

Taking the limit as $n \to \infty$ gives the exact volume:
\begin{align}
V 
&= \lim_{n\to\infty} \sum_{k=1}^{n} \pi\bigl(f(x_k)^2 - g(x_k)^2\bigr)\Delta x \\
&= \int_a^b \pi\bigl(f(x)^2 - g(x)^2\bigr)\,dx.
\end{align}

%-------------------------------------------------------------------------------
\subsubsection{Exercises}
\label{sss:exer-washer}
%-------------------------------------------------------------------------------

\begin{enumerate}
\setcounter{enumi}{2}

\item Use the washer method to find the volume of a cone containing a 
central hole formed by revolving the region bounded by
\[
y = R - \frac{R}{h}x,
\]
and the lines $y=r$ and $x=0$ about the $x$-axis.

\item Calculate the volume of the solid of revolution generated by 
revolving the region bounded by the curves $y = x^2$ and $y = x^3$, 
and the lines $x=1$ and $y=0$ about the $x$-axis.
\end{enumerate}

%-------------------------------------------------------------------------------
\subsubsection{Shell Method}
\label{sss:shel-beba}
%-------------------------------------------------------------------------------

The shell method is another technique for finding the volume of a solid 
of revolution. In some cases, it simplifies the setup and evaluation of 
the integral.

\medskip

Consider a region bounded by $y=f(x)$ on $[a,b]$, revolved about the 
$y$-axis. Instead of disks, we approximate the volume using cylindrical 
shells.

\medskip

Partition $[a,b]$ into $n$ subintervals of width
\[
\Delta x = \frac{b-a}{n},
\]
with right endpoints
\[
x_k = a + k\Delta x.
\]

Each vertical strip generates a cylindrical shell with:
\begin{itemize}
\item radius: $x_k$,
\item height: $|f(x_k)|$,
\item thickness: $\Delta x$.
\end{itemize}

The volume of the $k$-th shell is approximately
\[
V_k = 2\pi x_k |f(x_k)|\,\Delta x.
\]

Thus, the total approximate volume is
\[
V_{\text{approx}} = \sum_{k=1}^{n} 2\pi x_k |f(x_k)|\,\Delta x.
\]

Taking the limit as $n \to \infty$ gives
\[
V = \lim_{n\to\infty} \sum_{k=1}^{n} 2\pi x_k |f(x_k)|\,\Delta x
= \int_a^b 2\pi x |f(x)|\,dx.
\]

\medskip

\textbf{Remark:} The derivation above can also be obtained by expanding
\[
\pi\big((x_k+\Delta x)^2 - x_k^2\big) = \pi(2x_k\Delta x + \Delta x^2),
\]
and noting that the $\Delta x^2$ term vanishes in the limit.

%-------------------------------------------------------------------------------
\subsubsection{Exercises}
\label{sss:exer-shell}
%-------------------------------------------------------------------------------

\begin{enumerate}
\setcounter{enumi}{4}

\item Find the volume of a cone with radius $r$ and height $h$ by using 
the shell method on an appropriate region that, when rotated about the 
$y$-axis, produces the cone.

\item Calculate the volume of the solid of revolution generated by 
revolving the region bounded by $y=x^2$, $x=1$, and $y=0$ about the 
$y$-axis.
\end{enumerate}

%-------------------------------------------------------------------------------
\subsubsection{Exercise Solutions}
\label{sss:exer-sol-shell}
%-------------------------------------------------------------------------------

\begin{enumerate}
\item $\dfrac{\pi r^2 h}{3}$
\item $\dfrac{\pi}{2}$
\end{enumerate}

%-------------------------------------------------------------------------------
\subsubsection{Resources}
\label{sss:reso-55b5}
%-------------------------------------------------------------------------------

\begin{itemize}
\item \href{https://en.wikibooks.org/wiki/Calculus/Volume_of_solids_of_revolution}{Volume of Solids of Revolution}.  
Wikibooks: Calculus.
\end{itemize}

%-------------------------------------------------------------------------------
\subsubsection{Videos}
\label{sss:vide-cce3}
%-------------------------------------------------------------------------------

\begin{itemize}

\item \href{https://youtu.be/pCMkHkprN0I}{Volume of Revolution – Shell Method about the $x$-axis}.  
Video by James Sousa, Math is Power 4U.

\item \href{https://youtu.be/3B2YQbEzshg}{Volume of Revolution – Shell Method about the $y$-axis}.  
Video by James Sousa, Math is Power 4U.

\item \href{https://youtu.be/lp3_rmjbxZ8}{Volume of Revolution – Shell Method about a general axis}.  
Video by James Sousa, Math is Power 4U.

\item \href{https://youtu.be/ZyFaaKhNPXo}{Comparing the Washer and Shell Methods}.  
Video by James Sousa.

\item \href{https://www.youtube.com/watch?v=DHCWM-Pg_Yw}{Example: Volume of Revolution Using the Shell Method}.  
Video by James Sousa, Math is Power 4U.

\item \href{https://www.youtube.com/watch?v=aW_3JoRHsPU}{Example: Volume of Revolution Using the Shell Method}.  
Video by James Sousa, Math is Power 4U.

\item \href{https://www.youtube.com/watch?v=m2VopSFCzZ4}{Example: Volume of Revolution Using the Shell Method}.  
Video by James Sousa, Math is Power 4U.

\item \href{https://www.youtube.com/watch?v=MdbsmztIAkA}{Example: Volume of Revolution Using the Shell Method}.  
Video by James Sousa, Math is Power 4U.

\item \href{https://www.youtube.com/watch?v=FANXpoGp-hQ}{Example: Volume of Revolution Using the Shell Method}.  
Video by James Sousa, Math is Power 4U.

\item \href{https://www.youtube.com/watch?v=SIHrWcLsNYs}{Example: Shell Method with Horizontal Axis}.  
Video by James Sousa, Math is Power 4U.

\item \href{https://www.youtube.com/watch?v=fAKmfe_5QWw}{Example: Shell Method with Vertical Axis}.  
Video by James Sousa, Math is Power 4U.

\item \href{https://www.youtube.com/watch?v=6Ozz3J-LRrY}{Shell Method for Rotation about a Vertical Line}.  
Video by Khan Academy.

\item \href{https://www.youtube.com/watch?v=R-Qu3QWOEiA}{Shell Method for Rotation about a Horizontal Line}.  
Video by Khan Academy.

\item \href{https://www.youtube.com/watch?v=SfWrVNyP9E8}{Shell Method with Two Functions of $x$}.  
Video by Khan Academy.

\item \href{https://www.youtube.com/watch?v=OelluIKIkCY}{Shell Method with Two Functions of $y$}.  
Video by Khan Academy.

\item \href{https://youtu.be/CiXME1u-oyU}{Volumes of Revolution Using Cylindrical Shells}.  
Video by Patrick JMT.

\item \href{https://youtu.be/DQihcRym9Ew}{Cylindrical Shells about the $y$-axis or $x=c$}.  
Video by Krista King.

\item \href{https://youtu.be/Sa2FDL3hmJo}{Cylindrical Shells about the $x$-axis or $y=c$}.  
Video by Krista King.

\item \href{https://youtu.be/fE6vmorV-rA}{Shell Method – Volume of Revolution}.  
Video by The Organic Chemistry Tutor.

\end{itemize}

%-------------------------------------------------------------------------------
\subsubsection{Licensing}
\label{sss:lice-7352}
%-------------------------------------------------------------------------------

Content obtained and/or adapted from:

\begin{itemize}

\item \href{https://en.wikibooks.org/wiki/Calculus/Volume_of_solids_of_revolution}{Volume of Solids of Revolution, Wikibooks: Calculus}.  
Licensed under CC BY-SA.

\end{itemize}


\subsection{Learning Objectives}

\begin{enumerate}
    \item Here are three student learning objectives (SLOs) related to solids of revolution and volume calculation methods:
    \item \textbf{Apply the Disk Method:}
    \item \textbf{Outcome:} Students will be able to calculate the volume of a solid of revolution using the disk method.
    \item \textbf{Details:} This includes setting up and evaluating the integral for a solid generated by revolving a function around the x-axis. Students should be able to identify the function to revolve, determine the limits of integration, and apply the formula \(V = \pi \int_a^b f(x)^2 \, dx\) to find the volume of the solid.
    \item \textbf{Utilize the Washer Method:}
    \item \textbf{Outcome:} Students will be able to use the washer method to determine the volume of a solid of revolution with a hollow center.
    \item \textbf{Details:} This includes setting up and solving integrals for solids formed by revolving the area between two curves around the x-axis. Students should be able to use the formula \(V = \pi \int_a^b \left[f(x)^2 - g(x)^2\right] dx\) and apply it to find volumes of solids with an inner radius defined by one function and an outer radius defined by another.
    \item \textbf{Implement the Shell Method:}
    \item \textbf{Outcome:} Students will be able to compute the volume of a solid of revolution using the shell method.
    \item \textbf{Details:} This involves setting up and evaluating the integral for a solid generated by revolving a region around the y-axis. Students should be able to derive the volume formula \(V = 2\pi \int_a^b x f(x) \, dx\), where \(x\) represents the radius of the cylindrical shell and \(f(x)\) is the height of the shell. Students should also be able to apply this method to various functions and limits of integration.
\end{enumerate}

\subsection{Assessment Instrument}
\label{ssec:volu-b407}

Problems compare the shell and disk/washer methods on the same solid, requiring students to choose the more convenient axis and set up the correct integrand for cylindrical shells ($2\pi r h \,dr$) or washers.

\textbf{Example problems.}

\begin{enumerate}
    \item Use the shell method to find the volume of the solid obtained by revolving the region bounded by $y = x^3$, $x=1$, and $y=0$ about the $y$-axis.
    \item Find the volume of the solid formed by revolving $y = 2x - x^2$ about the $y$-axis using cylindrical shells.
    \item Compare the volume of the solid of revolution of $y=e^x$ on $[0,1]$ about the $x$-axis obtained by the disk method and verify using the shell method.
\end{enumerate}

%===============================================================================
\section{Arc length and surface area}
\label{sec:arcl-0676}
%===============================================================================


%-------------------------------------------------------------------------------
\subsubsection{Arc Length}
\label{sss:arcl-e028}
%-------------------------------------------------------------------------------


Suppose that we are given a function $f$ that is continuous on an

interval $[a,b]$ and we want to calculate the length of the curve drawn

out by the graph of $f(x)$ from $x=a$ to $x=b$ . If the graph were a

straight line this would be easy: the formula for the length of the

line is given by Pythagoras' theorem. And if the graph were a piecewise

linear function we can calculate the length by adding up the length of

each piece.


The problem is that most graphs are not linear. Nevertheless we can

estimate the length of the curve by approximating it with straight

lines. Suppose the curve $C$ is given by the formula $y=f(x)$ for

$a\le x\le b$ . We divide the interval $[a,b]$ into $n$ subintervals with equal width $\Delta x$ and endpoints $x_0,x_1,\ldots,x_n$ . Now let

$y_i=f(x_i)$ so $P_i=(x_i,y_i)$ is the point on the curve above $x_i$ .

The length of the straight line between $P_i$ and $P _{i+1}$ is


\[
\bigl|P_iP _{i+1}\bigr|=\sqrt{(y _{i+1}-y_i)^2+(x _{i+1}-x_i)^2}
\]


So an estimate of the length of the curve $C$ is the sum


\[
\sum _{i=0}^{n-1}\bigl|P_iP _{i+1}\bigr|
\]


As we divide the interval $[a,b]$ into more pieces this gives a better

estimate for the length of $C$ . In fact, we make that a definition.


\textbf{Length of a Curve}


The length of the curve $y=f(x)$ for $a\le x\le b$ is defined to be


\[
L=\lim _{n\to\infty}\sum _{i=0}^{n-1}\bigl|P _{i+1}P_i\bigr|
\]


%-------------------------------------------------------------------------------
\subsubsection{The Arclength Formula}
\label{sss:thea-fcc9}
%-------------------------------------------------------------------------------


Suppose that $f'$ is continuous on $[a,b]$ . Then the length of the

curve given by $y=f(x)$ between $a$ and $b$ is given by


\[
L=\int\limits_a^b \sqrt{1+f'(x)^2}dx
\] And in Leibniz notation


\[
L=\int\limits_a^b \sqrt{1+\left(\tfrac{dy}{dx}\right)^2}dx
\]


\textbf{Proof:} Consider $y _{i+1}-y_i=f(x _{i+1})-f(x_i)$ . By the Mean Value Theorem there is a point $z_i$ in $(x _{i+1},x_i)$ such that


\[
y _{i+1}-y_i=f(x _{i+1})-f(x_i)=f'(z_i)(x _{i+1}-x_i)
\]


So


\begin{align}
\bigl|P_iP _{i+1}\bigr| &= \sqrt{(x _{i+1} - x_i)^2 + (y _{i+1} - y_i)^2} \\&= \sqrt{(x _{i+1} - x_i)^2 + f'(z_i)^2(x _{i+1} - x_i)^2} \\&= \sqrt{\bigl(1 + f'(z_i)^2\bigr)(x _{i+1} - x_i)^2} \\&= \sqrt{1 + f'(z_i)^2} \,\Delta x
\end{align}


Putting this into the definition of the length of $C$ gives


\[
L=\lim _{n\to\infty}\sum _{i=0}^{n-1}\sqrt{1+f'(z_i)^2}\Delta x
\]


Now this is the definition of the integral of the function

$g(x)=\sqrt{1+f'(x)^2}$ between $a$ and $b$ (notice that $g$ is

continuous because we are assuming that $f'$ is continuous). Hence


\[
L=\int\limits_a^b \sqrt{1+f'(x)^2}dx
\]


as claimed.


\textbf{Example}


Length of the curve $y=2x$ from $x=0$ to $x=1$


As a sanity check of our formula, let's calculate the length of the

``curve'' $y=2x$ from $x=0$ to $x=1$ . First let's find the answer

using the Pythagorean Theorem.


\[
P_0=(0,0)
\] and


\[
P_1=(1,2)
\] so the length of the curve, $s$ , is


\[
s=\sqrt{2^2+1^2}=\sqrt5
\] Now let's use the formula


\begin{align}
s &= \int\limits_0^1 \sqrt{1+\left(\tfrac{d(2x)}{dx}\right)^2}\,dx \\&= \int\limits_0^1 \sqrt{1+2^2}\,dx \\&= \sqrt{5}x\bigg| _0^1 \\&= \sqrt{5}
\end{align}


%-------------------------------------------------------------------------------
\subsubsection{Exercises}
\label{sss:exer-arclength}
%-------------------------------------------------------------------------------

\begin{enumerate}
\item Find the length of the curve $y = x\sqrt{x}$ from $x=0$ to $x=1$.

\item Find the length of the curve 
\[
y = \frac{e^x + e^{-x}}{2}
\]
from $x=0$ to $x=1$.
\end{enumerate}

%-------------------------------------------------------------------------------
\subsubsection{Arc Length of a Parametric Curve}
\label{sss:arcl-cc81}
%-------------------------------------------------------------------------------

For a parametric curve defined by $x=f(t)$ and $y=g(t)$, the arc length is
\[
L = \int_a^b \sqrt{f'(t)^2 + g'(t)^2}\,dt.
\]

\textbf{Proof:} Consider
\[
y_{i+1} - y_i = g(t_{i+1}) - g(t_i), \quad 
x_{i+1} - x_i = f(t_{i+1}) - f(t_i).
\]

By the Mean Value Theorem, there exist $c_i, d_i \in (t_i, t_{i+1})$ such that
\[
y_{i+1} - y_i = g'(c_i)(t_{i+1} - t_i), \quad
x_{i+1} - x_i = f'(d_i)(t_{i+1} - t_i).
\]

Thus,
\begin{align}
|P_i P_{i+1}|
&= \sqrt{(x_{i+1}-x_i)^2 + (y_{i+1}-y_i)^2} \\
&= \sqrt{f'(d_i)^2 (t_{i+1}-t_i)^2 + g'(c_i)^2 (t_{i+1}-t_i)^2} \\
&= \sqrt{f'(d_i)^2 + g'(c_i)^2}\,\Delta t.
\end{align}

Summing and taking limits:
\[
L = \lim_{n\to\infty} \sum_{i=0}^{n-1} 
\sqrt{f'(d_i)^2 + g'(c_i)^2}\,\Delta t
= \int_a^b \sqrt{f'(t)^2 + g'(t)^2}\,dt.
\]

%-------------------------------------------------------------------------------
\subsubsection{Exercises}
\label{sss:exer-parametric}
%-------------------------------------------------------------------------------

\begin{enumerate}
\setcounter{enumi}{2}

\item Find the circumference of the circle defined by
\[
x(t) = R\cos(t), \quad y(t) = R\sin(t),
\]
for $0 \le t \le 2\pi$.

\item Find the length of one arch of the cycloid defined by
\[
x(t) = R\bigl(t - \sin(t)\bigr), \quad
y(t) = R\bigl(1 - \cos(t)\bigr),
\]
for $0 \le t \le 2\pi$.
\end{enumerate}

%-------------------------------------------------------------------------------
\subsubsection{Exercise Solutions}
\label{sss:exer-sol-arclength}
%-------------------------------------------------------------------------------

\begin{enumerate}
\item $\dfrac{13\sqrt{13} - 8}{27}$
\item $\dfrac{e - \frac{1}{e}}{2}$
\item $2\pi R$
\item $8R$
\end{enumerate}

%-------------------------------------------------------------------------------
\subsubsection{Surface Area}
\label{sss:surf-f4a9}
%-------------------------------------------------------------------------------

Suppose we are given a function $f(x)$ and wish to compute the surface 
area generated by revolving the curve about a line.

\medskip

If $f$ is linear, geometric formulas may be used. Otherwise, we approximate 
the surface using conical frusta.

\medskip

Divide $[a,b]$ into subintervals of width $\Delta x$. Each segment generates 
a frustum with surface area approximately
\[
A_i \approx 2\pi r_i \, l_i,
\]
where $r_i$ is the average radius and $l_i$ is the slant height.

Thus,
\[
A \approx \sum_{i=1}^n A_i.
\]

%-------------------------------------------------------------------------------
\subsubsection{Definition (Surface of Revolution)}
\label{sss:defi-cf5f}
%-------------------------------------------------------------------------------

The surface area of revolution is defined as
\[
A = \lim_{n\to\infty} \sum_{i=1}^n A_i.
\]

%-------------------------------------------------------------------------------
\subsubsection{The Surface Area Formula}
\label{sss:thes-528f}
%-------------------------------------------------------------------------------

Let $f$ be continuous on $[a,b]$, and let $r(x)$ be the distance from 
$f(x)$ to the axis of rotation. Then
\[
A = 2\pi \int_a^b r(x)\sqrt{1 + f'(x)^2}\,dx.
\]

In Leibniz notation:
\[
A = 2\pi \int_a^b r(x)\sqrt{1 + \left(\frac{dy}{dx}\right)^2}\,dx.
\]

\textbf{Proof:}
\begin{align}
A &= \lim_{n\to\infty} \sum_{i=1}^n A_i \\
  &= \lim_{n\to\infty} \sum_{i=1}^n 2\pi r_i l_i \\
  &= 2\pi \lim_{n\to\infty} \sum_{i=1}^n r_i l_i.
\end{align}

As $\Delta x \to 0$:
\begin{itemize}
\item $r_i \to r(x)$,
\item $l_i \approx \sqrt{1 + f'(x_i)^2}\,\Delta x$.
\end{itemize}

Thus,
\begin{align}
A &= 2\pi \lim_{n\to\infty} \sum_{i=1}^n 
r_i \sqrt{1 + f'(x_i)^2}\,\Delta x \\
&= 2\pi \int_a^b r(x)\sqrt{1 + f'(x)^2}\,dx.
\end{align}

If expressed in terms of $y$ on $[c,d]$, then
\[
A = 2\pi \int_c^d r(y)\sqrt{1 + \left(\frac{dx}{dy}\right)^2}\,dy.
\]

%-------------------------------------------------------------------------------
\subsubsection{Resources}
\label{sss:reso-55b5}
%-------------------------------------------------------------------------------

\begin{itemize}
\item \href{https://en.wikibooks.org/wiki/Calculus/Arc_length}{Arc Length}.  
Wikibooks: Calculus.

\item \href{https://en.wikibooks.org/wiki/Calculus/Surface_area}{Surface Area}.  
Wikibooks: Calculus.
\end{itemize}

\medskip

\textbf{Arc Length}

\begin{itemize}

\item \href{https://youtu.be/seoFxrNL85c}{Arc Length – Part 1 of 2}.  
Video by James Sousa, Math is Power 4U.

\item \href{https://youtu.be/NbnTw0opE_0}{Arc Length – Part 2 of 2}.  
Video by James Sousa, Math is Power 4U.

\item \href{https://www.youtube.com/watch?v=QttlclJGxJQ}{Example: Arc Length of a Linear Function}.  
Video by James Sousa, Math is Power 4U.

\item \href{https://www.youtube.com/watch?v=cw76RZ_fsZI}{Example: Arc Length of a Radical Function}.  
Video by James Sousa, Math is Power 4U.

\item \href{https://www.youtube.com/watch?v=0j227VZN0X8}{Example: Arc Length of a Quadratic Function}.  
Video by James Sousa, Math is Power 4U.

\item \href{https://youtu.be/yfJB4n-IzBE}{Deriving the Arc Length Formula}.  
Video by Patrick JMT.

\item \href{https://youtu.be/PwmCZAWeRNE}{Arc Length}.  
Video by Patrick JMT.

\item \href{https://youtu.be/tfn4cpkPHUI}{Arc Length for $y=f(x)$}.  
Video by Krista King.

\item \href{https://youtu.be/Mz3ELMAhMxk}{Arc Length for $x=g(y)$}.  
Video by Krista King.

\item \href{https://www.youtube.com/watch?v=8Y-snjheI9M}{Arc Length Introduction}.  
Video by Khan Academy.

\item \href{https://www.youtube.com/watch?v=OhISsmqv4_8}{Arc Length Example}.  
Video by Khan Academy.

\item \href{https://www.youtube.com/watch?v=MtRXjXdXDow}{Arc Length Example}.  
Video by Khan Academy.

\item \href{https://youtu.be/DNDAwWIL5FY}{Arc Length}.  
Video by The Organic Chemistry Tutor.

\end{itemize}

\medskip

\textbf{Surface Area}

\begin{itemize}

\item \href{https://youtu.be/4XLq-BWK5NY}{Surface Area of Revolution – Part 1 of 2}.  
Video by James Sousa, Math is Power 4U.

\item \href{https://youtu.be/u-kEdDCno44}{Surface Area of Revolution – Part 2 of 2}.  
Video by James Sousa, Math is Power 4U.

\item \href{https://www.youtube.com/watch?v=nRWNLlsziN4}{Example: Surface Area of Revolution (Linear Function)}.  
Video by James Sousa, Math is Power 4U.

\item \href{https://www.youtube.com/watch?v=hvOC6u26yK4}{Example: Surface Area of Revolution (Sine Function)}.  
Video by James Sousa, Math is Power 4U.

\item \href{https://www.youtube.com/watch?v=csWPkm-gJ8E}{Example: Surface Area of Revolution (Cubic Function about the $x$-axis)}.  
Video by James Sousa, Math is Power 4U.

\item \href{https://www.youtube.com/watch?v=NA1PpNLVGzw}{Example: Surface Area of Revolution (Square Root Function about the $x$-axis)}.  
Video by James Sousa, Math is Power 4U.

\item \href{https://www.youtube.com/watch?v=lZ9cEnagXBw}{Example: Surface Area of Revolution (Quadratic Function about the $y$-axis)}.  
Video by James Sousa, Math is Power 4U.

\item \href{https://www.youtube.com/watch?v=hGjaiwcEO9E}{Example: Surface Area of Revolution (Cube Root Function about the $y$-axis)}.  
Video by James Sousa, Math is Power 4U.

\item \href{https://youtu.be/-j2eKo84Ef8}{Finding Surface Area – Part 1}.  
Video by Patrick JMT.

\item \href{https://youtu.be/Jxf_XeKsiyY}{Finding Surface Area – Part 2}.  
Video by Patrick JMT.

\item \href{https://youtu.be/WlaFF-OgwJM}{Surface Area of Revolution Example 1}.  
Video by Krista King.

\item \href{https://youtu.be/VdDitAOifsY}{Surface Area of Revolution Example 2}.  
Video by Krista King.

\item \href{https://youtu.be/2fhKcatexdw}{Surface Area of Revolution Example 3}.  
Video by Krista King.

\item \href{https://youtu.be/lQM-0Nqs9Pg}{Surface Area of Revolution by Integration}.  
Video by The Organic Chemistry Tutor.

\end{itemize}

%-------------------------------------------------------------------------------
\subsubsection{Licensing}
\label{sss:lice-7352}
%-------------------------------------------------------------------------------

Content obtained and/or adapted from:

\begin{itemize}

\item \href{https://en.wikibooks.org/wiki/Calculus/Arc_length}{Arc Length, Wikibooks: Calculus}.  
Licensed under CC BY-SA.

\item \href{https://en.wikibooks.org/wiki/Calculus/Surface_area}{Surface Area, Wikibooks: Calculus}.  
Licensed under CC BY-SA.

\end{itemize}

\subsection{Learning Objectives}

\begin{enumerate}
    \item Here are three student learning objectives (SLOs) for the concept of arc length, formatted according to your specifications:
    \item \textbf{Calculate Arc Length for a Function:}
    \item \textbf{Outcome:} Students will be able to compute the arc length of a curve defined by a function $y = f(x)$ over a given interval $[a, b]$.
    \item \textbf{Details:} This involves applying the arc length formula \[
L = \int_a^b \sqrt{1 + \left(\frac{dy}{dx}\right)^2} \, dx
\] and accurately evaluating definite integrals to find the length of the curve.
    \item \textbf{Apply the Arc Length Formula to Parametric Curves:}
    \item \textbf{Outcome:} Students will be able to determine the length of a parametric curve given by $x = f(t)$ and $y = g(t)$ over a specific range of $t$.
    \item \textbf{Details:} This includes using the formula \[
L = \int_a^b \sqrt{\left(\frac{dx}{dt}\right)^2 + \left(\frac{dy}{dt}\right)^2} \, dt
\] to compute the arc length of parametric curves and interpreting the results within a given context.
    \item \textbf{Estimate Arc Length Using Approximations:}
    \item \textbf{Outcome:} Students will be able to estimate the arc length of a curve by approximating it with piecewise linear segments and understanding the effect of increasing the number of segments.
    \item \textbf{Details:} This involves calculating the length of each segment using \[
\bigl|P_iP_{i+1}\bigr| = \sqrt{(y_{i+1} - y_i)^2 + (x_{i+1} - x_i)^2}
\] and understanding how the approximation improves as the number of segments increases, thereby approaching the exact arc length.
\end{enumerate}

\subsection{Assessment Instrument}
\label{ssec:arcl-c62c}

Students apply the arc-length formula $\int\sqrt{1+(f\'(x))^2}\,dx$ and the surface-area formula $2\pi\int f(x)\sqrt{1+(f(x))^2}\,dx$, paying careful attention to algebraic simplification of the square root.

\textbf{Example problems.}

\begin{enumerate}
    \item Find the arc length of $y = \dfrac{2}{3}x^{3/2}$ from $x=0$ to $x=3$.
    \item Compute the surface area generated by revolving $y = \sqrt{x}$ on $[1, 4]$ about the $x$-axis.
    \item Find the arc length of the curve $x = \dfrac{y^3}{6} + \dfrac{1}{2y}$ from $y=1$ to $y=2$.
\end{enumerate}

%===============================================================================
\section{Applications in physics}
\label{sec:appl-8ba4}
%===============================================================================


%-------------------------------------------------------------------------------
\subsubsection{Work}
\label{sss:work-67e9}
%-------------------------------------------------------------------------------


For moving objects, the quantity of work/time (power) is integrated

along the trajectory of the point of application of the force. Thus, at

any instant, the rate of the work done by a force (measured in

joules/second, or \textbf{watts}) is the scalar product of the force (a

vector), and the velocity vector of the point of application. This

scalar product of force and velocity is known as instantaneous power.

Just as velocities may be integrated over time to obtain a total

distance, by the fundamental theorem of calculus, the total work along a

path is similarly the time-integral of instantaneous power applied along

the trajectory of the point of application.


Work is the result of a force on a point that follows a curve \textbf{X},

with a velocity \textbf{v}, at each instant. The small amount of work \textit{dW}

that occurs over an instant of time \textit{dt} is calculated as


\begin{align}
\delta W &= \mathbf{F}\cdot d\mathbf{s} \\&= \mathbf{F}\cdot\mathbf{v}\,dt
\end{align}


where the \textbf{F} · \textbf{v} is the power over the instant \textit{dt}. The sum of these small amounts of work over the trajectory of the point yields the work,


\begin{align}
W &=  \int _{t_1}^{t_2}\mathbf{F} \cdot \mathbf{v}\,dt \\&=  \int _{t_1}^{t_2}\mathbf{F} \cdot \tfrac{d\mathbf{s}}{dt}\,dt \\&= \int_C \mathbf{F} \cdot d\mathbf{s},
\end{align}


where \textit{C} is the trajectory from \textbf{x}($t_1$) to \textbf{x}($t_2$). This

integral is computed along the trajectory of the particle, and is

therefore said to be \textit{path dependent}.


If the force is always directed along this line, and the magnitude of

the force is \textit{F}, then this integral simplifies to


\[
W = \int_C F\,ds
\] where \textit{s} is displacement along the line. If \textbf{F}

is constant, in addition to being directed along the line, then the

integral simplifies further to


\begin{align}
W &= \int_C F\,ds \\&= F\int_C ds \\&= Fs
\end{align}


where \textit{s} is the displacement of

the point along the line.


This calculation can be generalized for a constant force that is not

directed along the line, followed by the particle. In this case the dot

product $ \mathbf{F} \cdot ds = F \cos \theta ds$, where $\theta$ is the angle

between the force vector and the direction of movement, that is


\begin{align}
W &= \int_C \mathbf{F} \cdot d\mathbf{s} \\&= Fs\cos\theta.
\end{align}


When a force component is perpendicular to the displacement of the

object (such as when a body moves in a circular path under a central

force), no work is done, since the cosine of 90° is zero. Thus, no work

can be performed by gravity on a planet with a circular orbit (this is

ideal, as all orbits are slightly elliptical). Also, no work is done on

a body moving circularly at a constant speed while constrained by

mechanical force, such as moving at constant speed in a frictionless

ideal centrifuge.


%-------------------------------------------------------------------------------
\subsubsection{Work done by a variable force}
\label{sss:work-2473}
%-------------------------------------------------------------------------------


Calculating the work as ``force times straight path segment'' would only

apply in the most simple of circumstances, as noted above. If force is

changing, or if the body is moving along a curved path, possibly

rotating and not necessarily rigid, then only the path of the

application point of the force is relevant for the work done, and only

the component of the force parallel to the application point velocity is

doing work (positive work when in the same direction, and negative when

in the opposite direction of the velocity). This component of force can be described by the scalar quantity called \textit{scalar tangential component}

($F \cos (\theta)$, where $\theta$ is the angle between the force and the velocity). And then the most

general definition of work can be formulated as follows:


$\quad$ *Work of a force is the line integral of its scalar tangential

    component along the path of its application point.*


$\quad$ If the force varies (e.g. compressing a spring) we need to use

    calculus to find the work done. If the force is given by

    $F(x)$ (a function of

    $x$ then the work done by the force along the

    x-axis from $a$ to $b$ is:


\[
W = \int _{a}^{b} \mathbf{F(s)} \cdot d\mathbf{s}
\]


%-------------------------------------------------------------------------------
\subsubsection{Work and potential energy}
\label{sss:work-305a}
%-------------------------------------------------------------------------------


The scalar product of a force \textbf{F} and the velocity \textbf{v} of its point

of application defines the power input to a system at an instant of

time. Integration of this power over the trajectory of the point of

application, \textit{C} = \textbf{x}(\textit{t}), defines the work input to the system by

the force.


%-------------------------------------------------------------------------------
\subsubsection{Path dependence}
\label{sss:path-5c74}
%-------------------------------------------------------------------------------


Therefore, the work done by a force \textbf{F} on an object that travels along a curve \textit{C} is given by the line integral:


\begin{align}
W &= \int_C \mathbf{F} \cdot d\mathbf{x} \\&=  \int _{t_1}^{t_2}\mathbf{F}\cdot \mathbf{v}\,dt,
\end{align}


where \textit{dx}(\textit{t}) defines the trajectory \textit{C} and \textbf{v} is the velocity

along this trajectory. In general this integral requires the path along

which the velocity is defined, so the evaluation of work is said to be

path dependent.


The time derivative of the integral for work yields the instantaneous

power,


\[
\frac{dW}{dt} = P(t) = \mathbf{F}\cdot \mathbf{v} .
\]


%-------------------------------------------------------------------------------
\subsubsection{Path independence}
\label{sss:path-59ea}
%-------------------------------------------------------------------------------


If the work for an applied force is independent of the path, then the

work done by the force, by the gradient theorem, defines a potential

function which is evaluated at the start and end of the trajectory of

the point of application. This means that there is a potential function

\textit{U}(\textbf{x}), that can be evaluated at the two points \textbf{x}($t_1$) and

\textbf{x}($t_2$) to obtain the work over any trajectory between these two

points. It is tradition to define this function with a negative sign so

that positive work is a reduction in the potential, that is


\begin{align}
W &= \int_C \mathbf{F} \cdot \mathrm{d}\mathbf{x} \\&=  \int _{\mathbf{x}(t_1)}^{\mathbf{x}(t_2)} \mathbf{F} \cdot \mathrm{d}\mathbf{x} \\&=  U(\mathbf{x}(t_1)) - U(\mathbf{x}(t_2)).
\end{align}


The function \textit{U}(\textbf{x}) is called the potential energy associated with

the applied force. The force derived from such a potential function is

said to be conservative. Examples of forces that have potential energies

are gravity and spring forces.


In this case, the gradient of work yields


$\quad$ $\nabla W =  -\nabla U= -\left(\frac{\partial U}{\partial x}, \frac{\partial U}{\partial y}, \frac{\partial U}{\partial z}\right) = \mathbf{F},$


and the force \textbf{F} is said to be ``derivable from a potential.''


Because the potential \textit{U} defines a force \textbf{F} at every point \textbf{x} in

space, the set of forces is called a force field. The power applied to a

body by a force field is obtained from the gradient of the work, or

potential, in the direction of the velocity \textbf{V} of the body, that is


$\quad$ $P(t) = -\nabla U \cdot \mathbf{v} = \mathbf{F}\cdot\mathbf{v}.$


%-------------------------------------------------------------------------------
\subsubsection{Work by gravity}
\label{sss:work-7f98}
%-------------------------------------------------------------------------------


% [Image: Gravity \textit{F} = \textit{mg} does work \textit{W} = \textit{mgh} along any descending

path]


In the absence of other forces, gravity results in a constant downward

acceleration of every freely moving object. Near Earth's surface the

acceleration due to gravity is $g = 9.8 m \cdot s^{-2}$ and the gravitational

force on an object of mass $m$ is $F' g = mg$. It is convenient to

imagine this gravitational force concentrated at the center of mass of

the object.


If an object with weight \textit{mg} is displaced upwards or downwards a

vertical distance $y_2 - y_1$, the work $W$ done on the object is:


\begin{align}
W &= F_g (y_2 - y_1) \\&= F_g \Delta y \\&= mg \Delta y
\end{align}


where $F_g$ is

weight (pounds in imperial units, and newtons in SI units), and $\Delta y$ is

the change in height $y$. Notice that the work done by gravity depends

only on the vertical movement of the object. The presence of friction

does not affect the work done on the object by its weight.


%-------------------------------------------------------------------------------
\subsubsection{Work by gravity in space}
\label{sss:work-670c}
%-------------------------------------------------------------------------------


The force of gravity exerted by a mass \textit{M} on another mass \textit{m} is given

by


\[
\mathbf{F}=-\frac{GMm}{r^3}\mathbf{r},
\] where \textbf{r} is the position

vector from \textit{M} to \textit{m}.


Let the mass \textit{m} move at the velocity \textbf{v}; then the work of gravity on this mass as it moves from position $\mathbf{r}(t_1$) to $\mathbf{r}(t_2$) is given by


\begin{align}
W &= -\int^{\mathbf{r}(t_2)} _{\mathbf{r}(t_1)}\frac{GMm}{r^3}\mathbf{r}\cdot d\mathbf{r} \\&= -\int^{t_2} _{t_1}\frac{GMm}{r^3}\mathbf{r}\cdot\mathbf{v}\,dt.
\end{align}


Notice that the position and velocity of the mass \textit{m} are given by


\[
\mathbf{r} = r\mathbf{e} _r, \qquad\mathbf{v} =  \frac{d\mathbf{r}}{dt} = \dot{r}\mathbf{e} _r + r\dot{\theta}\mathbf{e} _t,
\]

where $e_r$ and $e_t$ are the radial and tangential unit vectors directed relative to the vector from \textit{M} to \textit{m}, and we use the

fact that $d \mathbf{e} _r/dt  = \dot{\theta}\mathbf{e} _t$. Use this to simplify the formula for work of gravity to,


\begin{align}
W &= -\int^{t_2} _{t_1}\frac{GmM}{r^3}(r\mathbf{e} _r)\cdot(\dot{r}\mathbf{e} _r + r\dot{\theta}\mathbf{e} _t)\,dt \\&= -\int^{t_2} _{t_1}\frac{GmM}{r^3}r\dot{r}\,dt \\&= \frac{GMm}{r(t_2)} - \frac{GMm}{r(t_1)}.
\end{align}


This calculation uses the fact that


\[
\frac{d}{dt}r^{-1}=-r^{-2}\dot{r}=-\frac{\dot{r}}{r^2}.
\] The function


\[
U=-\frac{GMm}{r},
\] is the gravitational potential function, also

known as gravitational potential energy. The negative sign follows the

convention that work is gained from a loss of potential energy.


%-------------------------------------------------------------------------------
\subsubsection{Work by a spring}
\label{sss:work-1088}
%-------------------------------------------------------------------------------


% [Image: Forces in springs assembled in

parallel]


Consider a spring that exerts a horizontal force $\mathbf{F} = (-kx, 0, 0)$

that is proportional to its deflection in the \textit{x} direction independent

of how a body moves. The work of this spring on a body moving along the

space with the curve $\mathbf{X}(t) = (x(t)$, $y(t)$, $z(t))$, is

calculated using its velocity, $\mathbf{v} = (v_x, v_y, v_z)$, to

obtain


\[
W=\int_0^t\mathbf{F}\cdot\mathbf{v}dt =-\int_0^tkx v_x dt = -\frac{1}{2}kx^2.
\]

For convenience, consider contact with the spring occurs at $t = 0$,

then the integral of the product of the distance $x$ and the x-velocity,

$xv_x dt$, over time $t$ is $(1/2)x^2$. The work is the product of the distance times the spring force, which is also dependent on

distance; hence the $x^2$ result.


%-------------------------------------------------------------------------------
\subsubsection{Hooke's law for linear springs}
\label{sss:hook-9335}
%-------------------------------------------------------------------------------


Consider a simple helical spring that has one end attached to some fixed

object, while the free end is being pulled by a force whose magnitude is

$F_s$. Suppose that the spring has

reached a state of equilibrium, where its length is not changing anymore. Let $x$ be the amount by which the free end of the spring was displaced from its ``relaxed'' position (when it is not being stretched). Hooke's law states that


\[
F_s = kx
\]


or, equivalently,


\[
x = \frac{F_s}{k}
\]


where $k$ is a positive real number, characteristic of the spring. Moreover, the same formula holds when the spring is

compressed, with $F_s$ and $x$ both negative in that case. According to this formula, the graph of the applied force

$F_s$ as a function of the displacement

$x$ will be a straight line passing through the

origin, whose slope is $k$.


Hooke's law for a spring is sometimes, but rarely, stated under the

convention that $F_s$ is the restoring

force exerted by the spring on whatever is pulling its free end. In that

case, the equation becomes


\[
F_s = -kx
\] since the direction of the restoring force is opposite to

that of the displacement.


%-------------------------------------------------------------------------------
\subsubsection{Work by a gas}
\label{sss:work-2efb}
%-------------------------------------------------------------------------------


\[
W=\int_a^b{P}dV
\] Where $P$ is pressure, \textit{V} is volume, and \textit{a} and

\textit{b} are initial and final volumes.


%-------------------------------------------------------------------------------
\subsubsection{Hydrostatic pressure}
\label{sss:hydr-20bf}
%-------------------------------------------------------------------------------


In a fluid at rest, all frictional and inertial stresses vanish and the

state of stress of the system is called \textit{hydrostatic}. When this

condition of $V = 0$ is applied to the

Navier-Stokes equations, the gradient of pressure becomes a function of

body forces only. For a barotropic fluid in a conservative force field

like a gravitational force field, the pressure exerted by a fluid at

equilibrium becomes a function of force exerted by gravity.


The hydrostatic pressure can be determined from a control volume

analysis of an infinitesimally small cube of fluid. Since pressure is

defined as the force exerted on a test area ($p = \frac{F}{A}$, with

$p$: pressure, $F$: force

normal to area $A$, $A$:

area), and the only force acting on any such small cube of fluid is the

weight of the fluid column above it, hydrostatic pressure can be

calculated according to the following formula:


\begin{align}
p(z) - p(z_0) &= \frac{1}{A}\int _{z_0}^z dz' \iint_A dx' dy'\, \rho (z') g(z') \\&= \int _{z_0}^z dz'\, \rho (z') g(z').
\end{align}


where:


\begin{itemize}
    \item $p$ is the hydrostatic pressure (Pa),
\end{itemize}

\begin{itemize}
    \item $\rho$ is the fluid density (kg/$\text{m}^3$),
\end{itemize}

\begin{itemize}
    \item $g$ is gravitational acceleration (m/$\text{s}^2$),
\end{itemize}

\begin{itemize}
    \item $A$ is the test area ($\text{m}^2$),
\end{itemize}

\begin{itemize}
    \item $z$ is the height (parallel to the direction of gravity) of the test
\end{itemize}
    area ($m$),


\begin{itemize}
    \item $z_0$ is the height of the zero reference point of the pressure ($m$).
\end{itemize}

For water and other liquids, this integral can be simplified

significantly for many practical applications, based on the following

two assumptions: Since many liquids can be considered incompressible, a

reasonable good estimation can be made from assuming a constant density

throughout the liquid. (The same assumption cannot be made within a

gaseous environment.) Also, since the height $h$ of the fluid column between $z$ and

$z_0$ is often reasonably small compared to the radius of the Earth, one can neglect the variation of $g$. Under these circumstances, the integral is simplified into the formula:


\[
p - p_0 = \rho g h,
\]


where $h$ is the height

$z - z_0$ of the liquid column

between the test volume and the zero reference point of the pressure.

This formula is often called Stevin's law. Note that this reference point should lie at or below the surface of the liquid. Otherwise, one has to split the integral into two (or more) terms with the constant

$\rho _\text{liquid}$ and $\rho (z') _\text{above}$. For example, the

absolute pressure compared to vacuum is:


\[
p = \rho g H + p _\mathrm{atm},
\]


where $H$ is the total height of the liquid column above the test area to the surface, and

$p _\text{atm}$ is the atmospheric pressure,

i.e., the pressure calculated from the remaining integral over the air

column from the liquid surface to infinity. This can easily be

visualized using a pressure prism.


%-------------------------------------------------------------------------------
\subsubsection{Hydrostatic force on submerged surfaces}
\label{sss:hydr-12b6}
%-------------------------------------------------------------------------------


The horizontal and vertical components of the hydrostatic force acting

on a submerged surface are given by the following:


\[
F _\mathrm{h} = p _\mathrm{c}A
\]

\[
F _\mathrm{v} = \rho g V 
\]


where:


\begin{itemize}
    \item $p_c$ is the pressure at the centroid of the vertical projection of the submerged surface
\end{itemize}

\begin{itemize}
    \item $A$ is the area of the same vertical projection of the surface
\end{itemize}

\begin{itemize}
    \item $\rho$ is the density of the fluid
\end{itemize}

\begin{itemize}
    \item $g$ is the acceleration due to gravity
\end{itemize}

\begin{itemize}
    \item $V$ is the volume of fluid directly above the curved surface
\end{itemize}


%-------------------------------------------------------------------------------
\subsubsection{Resources}
\label{sss:reso-55b5}
%-------------------------------------------------------------------------------


\begin{itemize}
    \item \href{https://en.wikipedia.org/wiki/Work_(physics}{Work (physics)}),
\end{itemize}
    Wikipedia

\begin{itemize}
    \item \href{https://en.wikipedia.org/wiki/Hydrostatics}{Hydrostatics},
\end{itemize}
    Wikipedia

\begin{itemize}
    \item \href{https://en.wikipedia.org/wiki/Hooke\%27s_law}{Hooke's Law},
\end{itemize}
    Wikipedia


$\textbf{Work Done by a Variable Force}$


\begin{itemize}
    \item \href{https://youtu.be/KRKWJre--6U}{Work Done by a Variable Force} by
\end{itemize}
    Krista King

\begin{itemize}
    \item \href{https://youtu.be/LbAxcMQ7J6c}{Work Done By a Variable Force Physics
\end{itemize}
    Problems} by The Organic Chemistry

    Tutor

\begin{itemize}
    \item \href{https://youtu.be/TLw8xbmnY3c}{Work Problems - Calculus} by The
\end{itemize}
    Organic Chemistry Tutor


$\textbf{Work (Rope/Cable Problems)}$


\begin{itemize}
    \item \href{https://youtu.be/YMaR5KJj2zQ}{Ex 1: Integration Application - Work Lifting an
\end{itemize}
    Object} by James Sousa, Math is Power

    4U

\begin{itemize}
    \item \href{https://youtu.be/pKZzfYIgq1Q}{Ex 2: Integration Application - Work Lifting an Object and
\end{itemize}
    Cable} by James Sousa, Math is Power

    4U

\begin{itemize}
    \item \href{https://youtu.be/XEwdZqz93z4}{Ex: Find the Work Lifting a Leaking Bucket of Sand Given
\end{itemize}
    Mass} by James Sousa, Math is Power 4U

\begin{itemize}
    \item \href{https://youtu.be/XEwdZqz93z4}{Ex: Find the Work Lifting a Leaking Bucket of Sand and Rope Given
\end{itemize}
    Mass} by James Sousa, Math is Power 4U

\begin{itemize}
    \item \href{https://youtu.be/oWcN96kcbkM}{Finding Work using Calculus - The Cable/Rope Problem - Part
\end{itemize}
    b} by patrickJMT

\begin{itemize}
    \item \href{https://youtu.be/UJF_34TSAwQ}{Work done using a rope to lift a
\end{itemize}
    weight} by patrickJMT

\begin{itemize}
    \item \href{https://youtu.be/TLw8xbmnY3c}{Work Problems - Calculus} by The
\end{itemize}
    Organic Chemistry Tutor


$\textbf{Work (Spring Problem)}$


\begin{itemize}
    \item \href{https://youtu.be/3_-7IJO6EVc}{Ex: Find the Work Required to Stretch a
\end{itemize}
    Spring} by James Sousa, Math is Power

    4U

\begin{itemize}
    \item \href{https://youtu.be/Fd65ZxqpktE}{Ex: Find the Force Required to Stretch a
\end{itemize}
    Spring} by James Sousa, Math is Power

    4U

\begin{itemize}
    \item \href{https://youtu.be/x_0YWeHXZFE}{Work and Hooke's Law - Ex 1} by
\end{itemize}
    patrickJMT

\begin{itemize}
    \item \href{https://youtu.be/_YXA8JeGLjo}{Work and Hooke's Law - Ex 2} by
\end{itemize}
    patrickJMT

\begin{itemize}
    \item \href{https://youtu.be/ZKfYUxzlLx8}{Work Done on Elastic Springs} by
\end{itemize}
    Krista King

\begin{itemize}
    \item \href{https://youtu.be/D5vnqohxEOI}{Hooke's Law Physics, Basic Introduction, Restoring Force, Spring
\end{itemize}
    Constant} by The Organic Chemistry

    Tutor

\begin{itemize}
    \item \href{https://youtu.be/TLw8xbmnY3c}{Work Problems} by The Organic
\end{itemize}
    Chemistry Tutor


$\textbf{Work (Pumping Fluid Out of a Tank)}$


\begin{itemize}
    \item \href{https://youtu.be/1x-xMf4TFNM}{Ex: Determine the Work Required to Pump Water Out of a Circular
\end{itemize}
    Cylinder} by James Sousa, Math is

    Power 4U

\begin{itemize}
    \item \href{https://youtu.be/1x-xMf4TFNM}{Ex: Determine the Work Required to Pump Water Out of Trough
\end{itemize}
    (Isosceles Triangle)} by James Sousa,

    Math is Power 4U

\begin{itemize}
    \item \href{https://youtu.be/RPXnZl3przM}{Ex: Determine the Work Required to Pump Water Out of Trough
\end{itemize}
    (Quadratic Cross Section)} by James

    Sousa, Math is Power 4U

\begin{itemize}
    \item \href{https://youtu.be/fJtxJv5sdqo}{Calculating the Work Required to Drain a
\end{itemize}
    Tank} by patrickJMT

\begin{itemize}
    \item \href{https://youtu.be/TLw8xbmnY3c}{Work Problems - Calculus} by The
\end{itemize}
    Organic Chemistry Tutor


$\textbf{Hydrostatic Pressure and Force}$


\begin{itemize}
    \item \href{https://youtu.be/zSyh3IzQMdc}{Ex: Find the Hydrostatic Force on a Horizontal Plate (No
\end{itemize}
    Calculus)} by James Sousa, Math is

    Power 4U

\begin{itemize}
    \item \href{https://youtu.be/7exmlrEbNXI}{Ex: Find the Hydrostatic Force on a Vertical Plane in the Shape of
\end{itemize}
    a Isosceles Triangle} by James Sousa,

    Math is Power 4U

\begin{itemize}
    \item \href{https://youtu.be/Ntk9RpCCSeM}{Ex: Find the Hydrostatic Force on a Semicircle Window Submerged in
\end{itemize}
    Water} by James Sousa, Math is Power

    4U

\begin{itemize}
    \item \href{https://youtu.be/8i7xOtAo6ws}{Ex: Find the Hydrostatic Force on a Dam in the Shape of a Degree 4
\end{itemize}
    Polynomial} by James Sousa, Math is

    Power 4U

\begin{itemize}
    \item \href{https://youtu.be/12MhraQo0TY}{Hydrostatic Force - Basic Idea / Deriving the
\end{itemize}
    Formula} by patrickJMT

\begin{itemize}
    \item \href{https://youtu.be/cXNmCaTod58}{Hydrostatic Force - Complete Example
\end{itemize}

%-------------------------------------------------------------------------------
\subsubsection{1}
\label{sss:1-c4ca}
%-------------------------------------------------------------------------------
 by patrickJMT}

\begin{itemize}
    \item \href{https://youtu.be/h8kMaW2q9EM}{Hydrostatic Force - Complete Example #2, Part 1 of
\end{itemize}
    2} by patrickJMT

\begin{itemize}
    \item \href{https://youtu.be/F2poHPZZBhE}{Hydrostatic Force - Complete Example #2, Part 2 of
\end{itemize}
    2} by patrickJMT

\begin{itemize}
    \item \href{https://youtu.be/xTDA3H6-FAY}{Hydrostatic Pressure} by Krista King
    \item \href{https://youtu.be/FgwWPlYiZgY}{Hydrostatic Force} by Krista King
    \item \href{https://youtu.be/kkq8ruV8_Jw}{Introduction to Pressure \& Fluids} by
\end{itemize}
    The Organic Chemistry Tutor

\begin{itemize}
    \item \href{https://youtu.be/3jG-hWgUJko}{Hydrostatic Force Problems} by The
\end{itemize}
    Organic Chemistry Tutor


%-------------------------------------------------------------------------------
\subsubsection{Licensing}
\label{sss:lice-7352}
%-------------------------------------------------------------------------------


Content obtained and/or adapted from:


\begin{itemize}
    \item \href{https://en.wikipedia.org/wiki/Work_(physics}{Work (physics),
\end{itemize}
    Wikipedia}) under a CC

    BY-SA license

\begin{itemize}
    \item \href{https://en.wikipedia.org/wiki/Hydrostatics}{Hydrostatics,
\end{itemize}
    Wikipedia} under a CC

    BY-SA license

\begin{itemize}
    \item \href{https://en.wikipedia.org/wiki/Hooke\%27s_law}{Hooke's Law,
\end{itemize}
    Wikipedia} under a CC

    BY-SA license

\subsection{Learning Objectives}

\begin{enumerate}
    \item Here are three Student Learning Objectives (SLOs) for the topic of ``Work'' based on the provided content:
    \item \textbf{Calculate Work Done by a Constant Force:}
    \item \textbf{Outcome:} Students will be able to compute the work done by a constant force acting along a straight line.
    \item \textbf{Details:} This includes applying the formula $W = F \int_C ds = Fs$, where $F$ is the constant force and $s$ is the displacement along the line. Students should be able to solve problems involving forces that are directed along the displacement path and compute the total work done over that path.
    \item \textbf{Determine Work Done by a Variable Force Using Integration:}
    \item \textbf{Outcome:} Students will be able to calculate the work done by a variable force along a curved trajectory using the line integral.
    \item \textbf{Details:} This includes evaluating the integral $W = \int_C \mathbf{F} \cdot d\mathbf{s}$, where $\mathbf{F}$ is the force vector and $d\mathbf{s}$ is the differential displacement vector. Students should be able to handle situations where the force varies along the path or when the trajectory is not a straight line, applying integration techniques to find the total work done.
    \item \textbf{Relate Work to Potential Energy in Conservative Fields:}
    \item \textbf{Outcome:} Students will be able to connect the concept of work done by a conservative force to changes in potential energy.
    \item \textbf{Details:} This involves using the relationship $W = U(\mathbf{x}(t_1)) - U(\mathbf{x}(t_2))$ where $U$ is the potential energy associated with the force field. Students should understand how to compute work in scenarios where the force is conservative, and how to apply the gradient theorem to find potential energy differences between two points in the field.
\end{enumerate}

\subsection{Assessment Instrument}
\label{ssec:appl-d806}

Assessment covers work, hydrostatic force, and energy integrals. Students must set up the physical model (force function, pressure at depth) before integrating, and attach correct units throughout.

\textbf{Example problems.}

\begin{enumerate}
    \item A spring has natural length $0.5$\,m and requires $10$\,N to stretch it $0.1$\,m. Find the work done in stretching it from $0.5$\,m to $0.8$\,m.
    \item A cable of length $10$\,m and mass $3$\,kg/m hangs vertically. Find the work required to wind the entire cable onto a drum at the top.
    \item A triangular plate with base $2$\,m and height $3$\,m is submerged vertically with its apex at the surface. Find the hydrostatic force on one face ($\rho g \approx 9800$\,N/m$^3$).
\end{enumerate}

%===============================================================================
\section{Moments and Center of mass}
\label{sec:mome-3593}
%===============================================================================


%-------------------------------------------------------------------------------
\subsubsection{Definition of center of mass}
\label{sss:defi-7289}
%-------------------------------------------------------------------------------


The center of mass is the unique point at the center of a distribution

of mass in space that has the property that the weighted position

vectors relative to this point sum to zero. In analogy to statistics,

the center of mass is the mean location of a distribution of mass in

space.


%-------------------------------------------------------------------------------
\subsubsection{A system of particles}
\label{sss:asys-cc14}
%-------------------------------------------------------------------------------


In the case of a system of particles $P_i$, $i=1,$ $\dots$, $n$, each with mass $m_i$ that are located in space

with coordinates $\mathbf{r} _j$, $i=1$, $\dots$, $n,$ the coordinates $\mathbf{R}$ of the center of mass satisfy the condition


\[
\sum _{i=1}^n m_i(\mathbf{r} _i - \mathbf{R}) = \mathbf{0}.
\]


Solving this equation for \textbf{R} yields the formula


\[
\mathbf{R} = \frac 1M \sum _{i=1}^n m_i \mathbf{r} _i,
\]


where $M = \sum _{i = 1}^n m_i$ is the total mass of all of the particles.


%-------------------------------------------------------------------------------
\subsubsection{A continuous volume}
\label{sss:acon-26a8}
%-------------------------------------------------------------------------------


If the mass distribution is continuous with the density $\rho(\mathbf{r})$ within a solid $Q$, then the integral of the weighted position coordinates of the points in this volume relative to the center of mass $\mathbf{R}$ over the volume $\mathbf{V}$ is zero, that is


\[
\iiint\limits _{Q} \rho(\mathbf{r})(\mathbf{r}-\mathbf{R})dV = 0.
\]


Solve this equation for the coordinates $\mathbf{R}$ to obtain


\[
\mathbf R = \frac 1M \iiint\limits _{Q}\rho(\mathbf{r}) \mathbf{r} dV,
\]

where $M$ is the total mass in the volume.


If a continuous mass distribution has uniform density, which means $\rho$ is constant, then the center of mass is the same as the centroid of the volume.


%-------------------------------------------------------------------------------
\subsubsection{Locating the center of mass}
\label{sss:loca-0ff0}
%-------------------------------------------------------------------------------


% [Image: Plumb line method] 


Figure: Plumb line method


The experimental determination of a body's centre of mass makes use of gravity forces on the body and is based on the fact that the centre of mass is the same as the centre of gravity in the parallel gravity field near the earth's surface.


The center of mass of a body with an axis of symmetry and constant density must lie on this axis. Thus, the center of mass of a circular cylinder of constant density has its center of mass on the axis of the cylinder. In the same way, the center of mass of a spherically symmetric body of constant density is at the center of the sphere. In general, for any symmetry of a body, its center of mass will be a fixed point of that

symmetry.


%-------------------------------------------------------------------------------
\subsubsection{In two dimensions}
\label{sss:intw-f157}
%-------------------------------------------------------------------------------


An experimental method for locating the center of mass is to suspend the object from two locations and to drop plumb lines from the suspension points. The intersection of the two lines is the center of mass.


The shape of an object might already be mathematically determined, but it may be too complex to use a known formula. In this case, one can subdivide the complex shape into simpler, more elementary shapes, whose centers of mass are easy to find. If the total mass and center of mass can be determined for each area, then the center of mass of the whole is the weighted average of the centers. This method can even work for objects with holes, which can be accounted for as negative masses.


A direct development of the planimeter known as an integraph, or integerometer, can be used to establish the position of the centroid or center of mass of an irregular two-dimensional shape. This method can be applied to a shape with an irregular, smooth or complex boundary where other methods are too difficult. It was regularly used by ship builders to compare with the required displacement and center of buoyancy of a ship, and ensure it would not capsize.


%-------------------------------------------------------------------------------
\subsubsection{In three dimensions}
\label{sss:inth-92e6}
%-------------------------------------------------------------------------------


An experimental method to locate the three-dimensional coordinates of the center of mass begins by supporting the object at three points and measuring the forces, $\mathbf{F} _1$, $\mathbf{F} _2$, and $\mathbf{F} _3$ that resist the weight of the object, $\mathbf{W} = -W\mathbf{\hat{k}}$, where $\mathbf{\hat{k}}$ is the unit vector in the vertical direction. Let $\mathbf{r} _1$, $\mathbf{r} _2$, and $\mathbf{r} _3$ be the position coordinates of the support points, then the coordinates $\mathbf{R}$ of the center of mass satisfy the condition that the resultant torque is zero,


\[
\mathbf{T} = (\mathbf{r} _1 - \mathbf{R}) \times \mathbf{F} _1 + (\mathbf{r} _2 - \mathbf{R}) \times \mathbf{F} _2 + (\mathbf{r} _3 - \mathbf{R}) \times \mathbf{F} _3 = 0,
\]

or


\[
\mathbf{R} \times \left(-W\mathbf{\hat{k}}\right) = \mathbf{r} _1 \times \mathbf{F} _1 + \mathbf{r} _2 \times \mathbf{F} _2 + \mathbf{r} _3 \times \mathbf{F} _3.
\]


This equation yields the coordinates of the center of mass $\mathbf{R} ^{*}$ in the horizontal plane as,


\[
\mathbf{R}^{*} = -\frac{1}{W} \mathbf{\hat{k}} \times (\mathbf{r} _1 \times \mathbf{F} _1 + \mathbf{r} _2 \times\mathbf{F} _2 + \mathbf{r} _3 \times \mathbf{F} _3).
\]


The center of mass lies on the vertical line $\mathbf{L}$, given by


\[
\mathbf{L}(t) = \mathbf{R}^* + t\mathbf{\hat{k}}.
\]


The three-dimensional coordinates of the center of mass are determined by performing this experiment twice with the object positioned so that these forces are measured for two different horizontal planes through the object. The center of mass will be the intersection of the two lines $L_1$ and $L_2$ obtained from the two experiments.


The moment of inertia can be defined as the second moment about an axis and is usually designated the symbol \textit{I}. The moment of inertia is very useful in solving a number of problems in mechanics. For example, the moment of inertia can be used to calculate angular momentum, and angular energy. Moment of inertia is also important in beam design.


%-------------------------------------------------------------------------------
\subsubsection{Linear and angular momentum}
\label{sss:line-9ff6}
%-------------------------------------------------------------------------------


The linear and angular momentum of a collection of particles can be simplified by measuring the position and velocity of the particles relative to the center of mass. Let the system of particles $P_i$, $i=1$, $\dots$, $n$ of masses $m_i$ be located at the coordinates $\mathbf{r} _i$ with velocities $\mathbf{v} _i$. Select a reference point $\mathbf{R}$ and compute the relative position and velocity vectors,


\[
\mathbf{r} _i = (\mathbf{r} _i - \mathbf{R}) + \mathbf{R}, \quad \mathbf{v} _i = \frac{d}{dt}(\mathbf{r} _i - \mathbf{R}) + \mathbf{v}.
\]


The total linear momentum and angular momentum of the system are


$\mathbf{p} = \frac{d}{dt}\left(\sum _{i=1}^n m_i (\mathbf{r} _i - \mathbf{R})\right) + \left(\sum _{i=1}^n m_i\right) \mathbf{v},$


and


$\mathbf{L} = \sum _{i=1}^n m_i (\mathbf{r} _i - \mathbf{R}) \times \frac{d}{dt}(\mathbf{r} _i - \mathbf{R}) + \left(\sum _{i=1}^n m_i \right) \left[\mathbf{R} \times \frac{d}{dt}(\mathbf{r} _i - \mathbf{R}) + (\mathbf{r} _i - \mathbf{R}) \times \mathbf{v} \right] + \left(\sum _{i=1}^n m_i \right)\mathbf{R} \times \mathbf{v}$


If $\mathbf{R}$ is chosen as the center of mass these equations simplify to


\[
\mathbf{p} = m\mathbf{v},\quad \mathbf{L} = \sum _{i=1}^n m_i (\mathbf{r} _i - \mathbf{R}) \times \frac{d}{dt}(\mathbf{r} _i - \mathbf{R}) + \sum _{i=1}^n m_i \mathbf{R} \times \mathbf{v}
\]

where $m$ is the total mass of all the particles, $\mathbf{p}$ is the linear momentum, and $\mathbf{L}$ is the angular momentum.


The law of conservation of momentum predicts that for any system not subjected to external forces the momentum of the system will remain constant, which means the center of mass will move with constant velocity. This applies for all systems with classical internal forces, including magnetic fields, electric fields, chemical reactions, and so on. More formally, this is true for any internal forces that cancel in accordance with Newton's Third Law.


%-------------------------------------------------------------------------------
\subsubsection{Moment of Inertia}
\label{sss:mome-0a81}
%-------------------------------------------------------------------------------


%-------------------------------------------------------------------------------
\subsubsection{Shape moment of inertia for flat shapes}
\label{sss:shap-34d4}
%-------------------------------------------------------------------------------


The area moment of inertia takes only shape into account, not mass.


It can be used to calculate the moment of inertia of a flat shape about the x or y axis when $I$ is only important at one cross-section.


$I_x=\int y^2\,dA$


$I_y=\int x^2\,dA$


$I_z =\int r^2\,dA$


Because, for flat shapes, $r^2=x^2 + y^2$ the following is true

\[
I_z = I_x + I_y
\]


The shape moment of inertia of the cross-section of a beam is used in Structural Engineering in order to find the stress and deflection of the beam.


%-------------------------------------------------------------------------------
\subsubsection{Shape moment of inertia for 3D shapes}
\label{sss:shap-a7a8}
%-------------------------------------------------------------------------------


% [Image: The moment of inertia]


Figure: The moment of inertia $I=\int r2dm$ for a hoop, disk, cylinder, box, plate, rod, and spherical shell or solid can be found from this figure.


%-------------------------------------------------------------------------------
\subsubsection{Mass moment of inertia}
\label{sss:mass-ff43}
%-------------------------------------------------------------------------------


The mass moment of inertia takes mass into account. The mass moment of inertia of a point mass about a reference axis is equal to mass multiplied by the square of the distance from that point mass to the reference axis:


$I _\text{point mass} = r^2 m \,$


The metric units are kg*m^2.


The mass moment of inertia of any body of mass rotating around any axis is equal to the sum of the mass moment of inertia of each of the particles of that body:


$I_m = \sum r_i^2 m_i$


Rather than adding up each particle individually, sometimes we can take a mathematical shortcut by integrating over all the particles:


$I_m = \int r^2 \,dm$


%-------------------------------------------------------------------------------
\subsubsection{Radius of gyration}
\label{sss:radi-d7ff}
%-------------------------------------------------------------------------------


The radius of gyration is the radius at which you could concentrate the entire mass to make the moment of inertia equal to the actual moment of inertia. If the mass of an object was 2kg, and the moment of inertia was $18kg\*m^2$, then the radius of gyration would be 3m. In other words, if all of the mass was concentrated at a distance of 2m from the axis, then the moment of inertia would still be $18kg\*m^2$. Radius of gyration is represented with a $k$.


$k=\sqrt{I/m}\,$


The formula for the area radius of gyration replaces the mass with area.


$k=\sqrt{I/A}\,$


%-------------------------------------------------------------------------------
\subsubsection{Parallel axis theorem}
\label{sss:para-6dbb}
%-------------------------------------------------------------------------------


% [Image: x runs through the center of mass, while x' is parallel to

x]


Figure: $x$ runs through the center of mass, while $x'$ is parallel to $x$


If the moment of inertia is known about an axis that runs through the center of mass, then the moment of inertia about any parallel axis is given by,


$I _\text{parallel-axis}=I _\text{center\,of\,mass}+md^2$


where $d$ is the distance between the two axes of rotation.


%-------------------------------------------------------------------------------
\subsubsection{Resources}
\label{sss:reso-55b5}
%-------------------------------------------------------------------------------


\begin{itemize}
    \item \href{https://en.wikipedia.org/wiki/Center_of_mass}{Center of mass},
\end{itemize}
    Wikipedia

\begin{itemize}
    \item \href{https://en.wikibooks.org/wiki/Statics/Moment_of_Inertia_(contents}{Moment of
\end{itemize}
    Inertia}),

    Wikibooks: Statics


%-------------------------------------------------------------------------------
\subsubsection{Videos}
\label{sss:vide-cce3}
%-------------------------------------------------------------------------------


\href{https://youtu.be/-ICxOk6KINQ}{Moments and Center of Mass of a Discrete Set of

Objects} by patrickJMT


\href{https://youtu.be/NYUyHj3c1Xg}{Centroids / Centers of Mass - Part 1 of

2} by patrickJMT


\href{https://youtu.be/H5RcfMIZ_yw}{Centroids / Centers of Mass - Part 2 of

2} by patrickJMT


\href{https://youtu.be/H46E5yKFfNE}{Center of mass of the system, x-axis} by

Krista King


\href{https://youtu.be/xf2sa5Fr1fM}{Center of mass of the system} by Krista

King


\href{https://youtu.be/zuSkFdzyCO4}{Centroids of Plane Regions} by Krista

King


\href{https://youtu.be/Z3pYtlhLXBk}{Moments of the System} by Krista King


\href{https://youtu.be/JSGlBHAGvy4}{Moment, Center of Mass, and Centroid} by

The Organic Chemistry Tutor


\href{https://youtu.be/SWu_i-19Rn0}{Center of Mass \& Centroid Problems} by

The Organic Chemistry Tutor


\href{https://youtu.be/nSD2FRE6RTs}{Ex: Determine the Center of Mass of Three Point Masses in the

Plane} by James Sousa, Math is Power 4U


\href{https://youtu.be/fu7kNr6CR3g}{Ex: Find the Centroid of a Region Consisting of Three

Rectangles} by James Sousa, Math is Power

4U


\href{https://youtu.be/4s8DsmbECuI}{Ex: Find the Centroid of a Triangular Region in the

Plane} by James Sousa, Math is Power 4U


\href{https://youtu.be/k10APzmjKIA}{Ex: Find the Centroid of a Bounded Region Involving Two Quadratic

Functions} by James Sousa, Math is Power

4U


%-------------------------------------------------------------------------------
\subsubsection{Licensing}
\label{sss:lice-7352}
%-------------------------------------------------------------------------------


Content obtained and/or adapted from:


\begin{itemize}
    \item \href{https://en.wikipedia.org/wiki/Center_of_mass}{Center of mass,
\end{itemize}
    Wikipedia} under a CC

    BY-SA license

\begin{itemize}
    \item \href{https://en.wikibooks.org/wiki/Statics/Moment_of_Inertia_(contents}{Moment of Inertia, Wikibooks:
\end{itemize}
    Statics})

    under a CC BY-SA license

\subsection{Learning Objectives}

\begin{enumerate}
    \item Here are three Student Learning Objectives (SLOs) related to the definition of the center of mass:
    \item \textbf{System of Particles:}
    \item \textbf{Outcome:} Students will be able to determine the center of mass of a system of discrete particles using the formula for the center of mass.
    \item \textbf{Details:} This includes applying the formula \[
\mathbf{R} = \frac{1}{M} \sum_{i=1}^n m_i \mathbf{r}_i
\] where \( M = \sum_{i=1}^n m_i \) is the total mass. Students should be able to solve for the coordinates of the center of mass given the masses and positions of the particles.
    \item \textbf{Continuous Mass Distribution:}
    \item \textbf{Outcome:} Students will be able to calculate the center of mass for a continuous mass distribution by evaluating the integral of the weighted position coordinates.
    \item \textbf{Details:} This involves using the formula \[
\mathbf{R} = \frac{1}{M} \iiint\limits_{Q} \rho(\mathbf{r}) \mathbf{r} \, dV
\] where \( M \) is the total mass and \( \rho(\mathbf{r}) \) is the density function. Students should be able to handle cases where the density is uniform or non-uniform and solve integrals accordingly.
    \item \textbf{Experimental Determination in Two Dimensions:}
    \item \textbf{Outcome:} Students will be able to experimentally locate the center of mass of a two-dimensional object using suspension and plumb lines.
    \item \textbf{Details:} This includes suspending the object from two points, dropping plumb lines, and finding the intersection point, which represents the center of mass. Students should also be able to handle complex shapes by subdividing them into simpler components and calculating the center of mass as a weighted average of the component centers.
\end{enumerate}

\subsection{Assessment Instrument}
\label{ssec:mome-9a2a}

Problems require computing moments about each axis and the coordinates $(\bar{x}, \bar{y})$ of the centroid for planar regions of uniform density. Students must identify the correct moment formula and evaluate the integrals.

\textbf{Example problems.}

\begin{enumerate}
    \item Find the centroid of the region bounded by $y = x^2$ and $y = x$ on $[0,1]$.
    \item Find $\bar{x}$ and $\bar{y}$ for the triangular region with vertices $(0,0)$, $(4,0)$, and $(0,3)$.
    \item A lamina occupies the region under $y = \sin x$ on $[0,\pi]$ with uniform density. Find its centroid.
\end{enumerate}

%===============================================================================
\section{Integration by parts}
\label{sec:inte-5441}
%===============================================================================

Continuing on the path of reversing derivative rules in order to make them useful for integration, we reverse the product rule.


%-------------------------------------------------------------------------------
\subsubsection{Integration by parts}
\label{sss:inte-5441}
%-------------------------------------------------------------------------------


If $y=uv$ where $u$ and $v$ are functions of $x$ , then


\[
y'=(uv)'=u'v+uv'
\] Rearranging,


\[
uv'=(uv)'-u'v
\] Therefore,


\[
\int uv'dx=\int (uv)'\ dx-\int u'v\ dx
\] Therefore,


\[
\int uv'\ dx=uv-\int vu'\ dx
\] or


\[
\int u\ dv=uv-\int v\ du
\] This is the integration by parts formula.

It is very useful in many integrals involving products of functions, as

well as others.


For instance, to treat


\[
\int x\sin(x)dx
\] we choose $u=x$ and $dv=\sin(x)dx$ . With these choices, we have $du=dx$ and $v=-\cos(x)$ , and we have


\begin{align}
\int x\sin(x)\,dx &= -x\cos(x) - \int\big(-\cos(x)\big)\,dx \\&= -x\cos(x) + \int\cos(x)\,dx \\&= \sin(x) - x\cos(x) + C
\end{align}


Note that the choice of $u$ and $dv$ was critical. Had we chosen the

reverse, so that $u=\sin(x)$ and $dv=x\ dx$ , the result would have been


\[
\frac{x^2\sin(x)}{2}-\int\frac{x^2\cos(x)}{2}dx
\] The resulting

integral is no easier to work with than the original; we might say that

this application of integration by parts took us in the wrong direction.


So the choice is important. One general guideline to help us make that

choice is, if possible, to choose $u$ to be the factor of the integrand

which \textit{becomes simpler} when we differentiate it. In the last example,

we see that $\sin(x)$ does not become simpler when we differentiate it:

$\cos(x)$ is no simpler than $\sin(x)$ .


An important feature of the integration by parts method is that we often

need to apply it more than once. For instance, to integrate


\[
\int x^2e^xdx
\] we start by choosing $u=x^2$ and $dv=e^xdx$ to get


\[
\int x^2e^xdx=x^2e^x-2\int xe^xdx
\] Note that we still have an

integral to take care of, and we do this by applying integration by parts again, with $u=x$ and $dv=e^xdx$, which gives us


\begin{align}
\int x^2e^x\,dx &= x^2e^x - 2\int xe^x\,dx \\&= x^2e^x - 2(xe^x - e^x) + C \\&= x^2e^x - 2xe^x + 2e^x + C
\end{align}


So, two applications of integration by parts were necessary, owing to

the power of $x^2$ in the integrand.


Note that \textit{any power of x} does become simpler when we differentiate it,

so when we see an integral of the form


\[
\int x^nf(x)dx
\] one of our first thoughts ought to be to consider

using integration by parts with $u=x^n$ . Of course, in order for it to

work, we need to be able to write down an antiderivative for $dv$ .


%-------------------------------------------------------------------------------
\subsubsection{Example}
\label{sss:exam-1a79}
%-------------------------------------------------------------------------------

Use integration by parts to evaluate the integral
\[
\int e^x \sin(x)\,dx.
\]

\textbf{Solution:} Let $u = \sin(x)$ and $dv = e^x\,dx$. Then $du = \cos(x)\,dx$ and $v = e^x$. Using integration by parts,
\[
\int e^x \sin(x)\,dx = e^x \sin(x) - \int e^x \cos(x)\,dx.
\]

Now apply integration by parts again with $u = \cos(x)$ and $dv = e^x\,dx$, so that $du = -\sin(x)\,dx$ and $v = e^x$:
\[
\int e^x \cos(x)\,dx = e^x \cos(x) - \int e^x(-\sin(x))\,dx
= e^x \cos(x) + \int e^x \sin(x)\,dx.
\]

Substituting back,
\[
\int e^x \sin(x)\,dx = e^x \sin(x) - \left(e^x \cos(x) + \int e^x \sin(x)\,dx\right).
\]

Simplifying,
\[
\int e^x \sin(x)\,dx = e^x \sin(x) - e^x \cos(x) - \int e^x \sin(x)\,dx.
\]

Adding $\int e^x \sin(x)\,dx$ to both sides:
\[
2\int e^x \sin(x)\,dx = e^x(\sin(x) - \cos(x)).
\]

Hence,
\[
\int e^x \sin(x)\,dx = \frac{e^x(\sin(x) - \cos(x))}{2} + C.
\]

%-------------------------------------------------------------------------------
\subsubsection{With Definite Integrals}
\label{sss:with-3ba0}
%-------------------------------------------------------------------------------

For definite integrals, the rule is essentially the same, but the limits must be retained.

\textbf{Integration by Parts for Definite Integrals:}  
If $f$ and $g$ are differentiable functions, then
\begin{align}
\int_a^b f(x) g'(x)\,dx 
&= \bigl(f(x)g(x)\bigr)\Big|_a^b - \int_a^b f'(x) g(x)\,dx \\
&= f(b)g(b) - f(a)g(a) - \int_a^b f'(x) g(x)\,dx.
\end{align}

In Leibniz notation:
\[
\int_a^b u\,dv = (uv)\Big|_a^b - \int_a^b v\,du.
\]

%-------------------------------------------------------------------------------
\subsubsection{More Examples}
\label{sss:more-cdfa}
%-------------------------------------------------------------------------------

\begin{itemize}
\item \href{http://en.wikiversity.org/wiki/Application_of_Integration_by_Parts}{Examples Set 1: Integration by Parts}.
\end{itemize}

%-------------------------------------------------------------------------------
\subsubsection{Exercises}
\label{sss:exer-ef8a}
%-------------------------------------------------------------------------------

Evaluate the following using integration by parts:

\begin{enumerate}
\item $\int -4\ln(x)\,dx$
\item $\int (38-7x)\cos(x)\,dx$
\item $\int_0^{\tfrac{\pi}{2}} (-6x+45)\cos(x)\,dx$
\item $\int (5x+1)(x-6)^4\,dx$
\item $\int_{-1}^1 (2x+8)^3(2-x)\,dx$
\end{enumerate}

%-------------------------------------------------------------------------------
\subsubsection{Exercise Solutions}
\label{sss:exer-1604}
%-------------------------------------------------------------------------------

\begin{enumerate}
\item $4x - 4x\ln(x) + C$
\item $(38 - 7x)\sin(x) - 7\cos(x) + C$
\item $51 - 3\pi$
\item $\dfrac{(5x+1)(x-6)^5}{5} - \dfrac{(x-6)^6}{6} + C$
\item $1916.8$
\end{enumerate}

%-------------------------------------------------------------------------------
\subsubsection{Resources}
\label{sss:reso-55b5}
%-------------------------------------------------------------------------------

\begin{itemize}

\item \href{https://www.khanacademy.org/math/integral-calculus/integration-techniques/integration_by_parts/}{Videos and exercises on integration by parts}.  
Khan Academy.

\end{itemize}

%-------------------------------------------------------------------------------
\subsubsection{Videos}
\label{sss:vide-cce3}
%-------------------------------------------------------------------------------

\begin{itemize}

\item \href{https://youtu.be/815s-YB9X44}{Integration by Parts: The Basics}.  
Video by James Sousa.

\item \href{https://youtu.be/hfJ1zPizj_s}{Integration by Parts (After the Basics)}.  
Video by James Sousa.

\item \href{https://youtu.be/_kJNhoTJ4D4}{Integration by Parts – Additional Examples}.  
Video by James Sousa.

\item \href{https://youtu.be/q_BP4LyvdOA}{Deriving the Integration by Parts Formula}.  
Video by Patrick JMT.

\item \href{https://youtu.be/dqaDSlYdRcs}{Integration by Parts Made Easy}.  
Video by Patrick JMT.

\item \href{https://youtu.be/97RasvYIvl4}{Integration by Parts – Indefinite Integral}.  
Video by Patrick JMT.

\item \href{https://youtu.be/_MuKOm8IHBA}{Integration by Parts – Using IBPs Twice}.  
Video by Patrick JMT.

\item \href{https://youtu.be/L8HoALsls-E}{Integration by Parts – A Loopy Example}.  
Video by Patrick JMT.

\item \href{https://youtu.be/zGGI4PkHzhI}{Integration by Parts – Definite Integral}.  
Video by Patrick JMT.

\item \href{https://youtu.be/Eb5EDMHDnx8}{How Does Integration by Parts Work?}.  
Video by Krista King.

\item \href{https://youtu.be/SlBp9hZBqaQ}{Integration by Parts}.  
Video by Krista King.

\item \href{https://youtu.be/s8l1vs4gjtY}{Integration by Parts – Example 2}.  
Video by Krista King.

\item \href{https://youtu.be/0clmZGVl0ss}{Integration by Parts – Example 3}.  
Video by Krista King.

\item \href{https://youtu.be/CxV5HE6chSY}{Integration by Parts – Example 4}.  
Video by Krista King.

\end{itemize}

%-------------------------------------------------------------------------------
\subsubsection{Licensing}
\label{sss:lice-7352}
%-------------------------------------------------------------------------------

Content obtained and/or adapted from:

\begin{itemize}

\item \href{https://en.wikibooks.org/wiki/Calculus/Integration_techniques/Integration_by_Parts}{Integration by Parts, Wikibooks}.  
Licensed under CC BY-SA.

\end{itemize}

\subsection{Learning Objectives}

\begin{enumerate}
    \item Here are three student learning objectives (SLOs) related to integration by parts:
    \item \textbf{Apply Integration by Parts:}
    \item \textbf{Outcome:} Students will be able to apply the integration by parts formula to evaluate integrals involving the product of two functions.
    \item \textbf{Details:} This includes selecting appropriate functions for $u$ and $dv$, computing their derivatives and antiderivatives, and applying the formula \[
\int u \, dv = uv - \int v \, du
\] to solve integrals.
    \item \textbf{Choose Appropriate Functions for Integration by Parts:}
    \item \textbf{Outcome:} Students will be able to make informed choices for $u$ and $dv$ in the integration by parts formula to simplify the integral.
    \item \textbf{Details:} This involves identifying which function should be $u$ and which should be $dv$ by considering which function simplifies when differentiated and which simplifies when integrated, using examples such as \[
\int x \sin(x) \, dx
\] and \[
\int e^x \sin(x) \, dx
\].
    \item \textbf{Solve Multiple Applications of Integration by Parts:}
    \item \textbf{Outcome:} Students will be able to solve integrals that require multiple applications of integration by parts.
    \item \textbf{Details:} This includes performing integration by parts more than once when necessary, as demonstrated in integrals like \[
\int x^2 e^x \, dx
\] and recognizing when the method needs to be reapplied to solve the integral completely.
\end{enumerate}

\subsection{Assessment Instrument}
\label{ssec:inte-4b01}

Assessment requires applying the integration-by-parts formula $\int u\,dv = uv - \int v\,du$ with deliberate choice of $u$ and $dv$ guided by LIATE. At least one problem requires applying the formula twice or solving for the integral algebraically.

\textbf{Example problems.}

\begin{enumerate}
    \item Evaluate $\int x\,e^x\,dx$ by parts, stating your choices of $u$ and $dv$.
    \item Compute $\int x^2\cos x\,dx$ by applying integration by parts twice.
    \item Find $\int e^x\sin x\,dx$ using integration by parts, then solve algebraically for the integral on both sides.
\end{enumerate}

%===============================================================================
\section{Trigonometric Integrals}
\label{sec:trigonometric-integrals}
%===============================================================================

When the integrand is primarily or exclusively based on trigonometric

functions, the following techniques are useful.


%-------------------------------------------------------------------------------
\subsubsection{Powers of Sine and Cosine}
\label{sss:powe-2468}
%-------------------------------------------------------------------------------


We will give a general method to solve generally integrands of the form

$\cos^m(x)\cdot\sin^n(x)$ . First let us work through an example.


\[
\int\cos^3(x)\sin^2(x)dx
\]


Notice that the integrand contains an odd power of cos. So rewrite it as


\[
\int\cos^2(x)\sin^2(x)\cos(x)dx
\]


We can solve this by making the substitution $u=\sin(x)$ so

$du=\cos(x)dx$ . Then we can write the whole integrand in terms of $u$

by using the identity


\[
\cos^2(x)=1-\sin^2(x)=1-u^2
\]


So


\begin{align}
\int\cos^3(x)\sin^2(x)\,dx &= \int\cos^2(x)\sin^2(x)\cos(x)\,dx \\&= \int (1 - u^2)u^2\,du \\&= \int u^2\,du - \int u^4\,du \\&= \frac{u^3}{3} - \frac{u^5}{5} + C \\&= \frac{\sin^3(x)}{3} - \frac{\sin^5(x)}{5} + C
\end{align}


This method works whenever there is an odd power of sine or cosine.


> To evaluate $\int\cos^m(x)\sin^n(x)dx$ when \textbf{either} $m$ or $n$ is

> \textbf{odd}.

>

> -   If $m$ is odd substitute $u=\sin(x)$ and use the identity

>     $\cos^2(x)=1-\sin^2(x)=1-u^2$ .

> -   If $n$ is odd substitute $u=\cos(x)$ and use the identity

>     $\sin^2(x)=1-\cos^2(x)=1-u^2$ .


%-------------------------------------------------------------------------------
\subsubsection{Example}
\label{sss:exam-1a79}
%-------------------------------------------------------------------------------


Find $\int\limits_0^\frac{\pi}{2} \cos^{40}(x)\sin^3(x)dx$ .


As there is an odd power of $\sin$ we let $u=\cos(x)$ so $du=-\sin(x)dx$

. Notice that when $x=0$ we have $u=\cos(0)=1$ and when

$x=\frac{\pi}{2}$ we have $u=\cos(\tfrac{\pi}{2})=0$ .


\begin{align} \int\limits_0^\frac{\pi}{2} \cos^{40}(x)\sin^3(x)\,dx  &= \int\limits_0^\frac{\pi}{2} \cos^{40}(x)\sin^2(x)\sin(x)\,dx \\ &= -\int\limits_1^0 u^{40}(1-u^2)\,du \quad \text{(using $u = \cos(x)$, $du = -\sin(x)\,dx$)} \\ &= \int\limits_0^1 u^{40}(1-u^2)\,du \\ &= \int\limits_0^1 \left(u^{40} - u^{42}\right)\,du \\ &= \left(\frac{u^{41}}{41} - \frac{u^{43}}{43}\right)\Bigg| _0^1 \\ &= \frac{1}{41} - \frac{1}{43} \end{align}


When both $m$ and $n$ are even, things get a little more complicated.


> To evaluate $\int\cos^m(x)\sin^n(x)dx$ when both $m$ and $n$ are

> \textbf{even}.

>

> > Use the identities $\sin^2(x)=\frac{1-\cos(2x)}{2}$ and

> $\cos^2(x)=\frac{1+\cos(2x)}{2}$ .


\textbf{Preparation for Example 1}


Problem: Solve $\int \cos(2x) \cos^2(2x)dx$ by substitution.


We'll use the substitution $u = \frac{1}{2} \sin(2x)$ and compute $du=\cos(2x)dx$.


Note that $2u = \sin(2x)$ so $4u^2 = \sin^2(2x)$


The identity implies

\[
\cos^2(y) + \sin^2(y) = 1 \text{ for } y = 2x \text{ } \Rightarrow \cos^2(2x) = 1 - \sin^2(2x)
\]

So, we compute


\begin{align}
\int \cos(2x) \cos^2(2x)\,dx &= \int \cos^2(2x) \cos(2x)\,dx \\&= \int \left( 1 - \sin^2(2x) \right) \cos(2x)\,dx \\&= \int \left( 1 - 4u^2 \right) du \quad \text{(using } 4u^2 = \sin^2(2x),\, du = \cos(2x)\,dx\text{)} \\&= u - \frac{4u^3}{3} + C \\&= \frac{1}{2} \sin(2x) - \frac{\sin^3(2x)}{6} + C
\end{align}


%-------------------------------------------------------------------------------
\subsubsection{Example}
\label{sss:exam-1a79}
%-------------------------------------------------------------------------------


Find $\int\sin^2(x)\cos^4(x)dx$ .


As $\sin^2(x)=\frac{1-\cos(2x)}{2}$ and $\cos^2(x)=\frac{1+\cos(2x)}{2}$

we have


\begin{align}
& \int \sin^2(x) \cos^4(x) \, dx \\& = \int \left(\frac{1 - \cos(2x)}{2}\right) \cdot \frac{(1 + \cos(2x))^2}{4} \, dx \\& = \frac{1}{8} \int (1 - \cos(2x))(1 + 2\cos(2x) + \cos^2(2x)) \, dx \\& = \frac{1}{8} \int \left(1 + \cos(2x) - \cos^2(2x) - \cos^3(2x)\right) \, dx \\& = \frac{1}{8} \left[\int 1 \, dx + \int \cos(2x) \, dx - \int \cos^2(2x) \, dx - \int \cos^3(2x) \, dx\right]
\end{align}


We evaluate $\int \cos^2 (2x) dx$ as

\begin{align}
\int \cos^2 (2x)\, dx &= \int \frac{1 + \cos(4x)}{2}\, dx \\&= \frac{1}{2} \int (1 + \cos(4x))\, dx \\&= \frac{x}{2} + \frac{1}{8}\sin(4x) + C
\end{align}


We now proceed with the original problem

\begin{align}
\int \sin^2(x) \cos^4(x) \, dx &= \frac{1}{8} \left[\int 1 \, dx + \int \cos(2x) \, dx - \int \cos^2(2x) \, dx - \int \cos^3(2x) \, dx\right] \\&= \frac{1}{8} \left[x + \frac{1}{2}\sin(2x) - \left( \frac{x}{2} + \frac{1}{8}\sin(4x) \right) - \left( \frac{1}{2} \sin(2x) - \frac{\sin^3(2x)}{6} \right) \right] \\&= \frac{1}{8} \left[x - \left( \frac{x}{2} + \frac{1}{8}\sin(4x) \right) + \frac{\sin^3(2x)}{6} \right] \\&= \frac{1}{8} \left[\frac{x}{2} - \frac{1}{8}\sin(4x) + \frac{\sin^3(2x)}{6} \right] + C  \\&= \frac{x}{16} - \frac{\sin(4x)}{64} + \frac{\sin^3(2x)}{48} + C
\end{align}


%-------------------------------------------------------------------------------
\subsubsection{Powers of Tan and Secant}
\label{sss:powe-35bd}
%-------------------------------------------------------------------------------


> To evaluate $\int\tan^m(x)\sec^n(x)dx$ .

>

> 1.  If $n$ is even and $n\ge 2$ then substitute $u=tan(x)$ and use the

>     identity $\sec^2(x)=1+\tan^2(x)$ .

> 2.  If $n$ and $m$ are both odd then substitute $u=\sec(x)$ and use

>     the identity $\tan^2(x)=\sec^2(x)-1$ .

> 3.  If $n$ is odd and $m$ is even then use the identity

>     $\tan^2(x)=\sec^2(x)-1$ and apply a reduction formula to integrate

>     $\sec^j(x)dx$ , using the examples below to integrate when $j=1,2$

>     .


%-------------------------------------------------------------------------------
\subsubsection{Example 1}
\label{sss:exam-acc3}
%-------------------------------------------------------------------------------


Find $\int\sec^2(x)dx$ .


There is an even power of $\sec(x)$ . Substituting $u=\tan(x)$ gives

$du=\sec^2(x)dx$ so


$\int\sec^2(x)dx=\int du=u+C=\tan(x)+C.$


%-------------------------------------------------------------------------------
\subsubsection{Example 2}
\label{sss:exam-5bdd}
%-------------------------------------------------------------------------------


Find $\int\tan(x)dx$ .


Let $u=\cos(x)$ so $du=-\sin(x)dx$ . Then


\begin{align}
\int\tan(x)\,dx 
&= \int\frac{\sin(x)}{\cos(x)}\,dx \\&= \int -\frac{du}{u} \quad \text{where } u = \cos(x), \, du = -\sin(x)\,dx \\&= -\ln|u| + C \\&= -\ln\Big|\cos(x)\Big| + C \\&= \ln\Big|\sec(x)\Big| + C
\end{align}


%-------------------------------------------------------------------------------
\subsubsection{Example 3}
\label{sss:exam-ad87}
%-------------------------------------------------------------------------------


Find $\int\sec(x)dx$ .


The trick to do this is to multiply and divide by the same thing like

this:


\begin{align}
\int\sec(x)\,dx 
&= \int\sec(x)\frac{\sec(x) + \tan(x)}{\sec(x) + \tan(x)}\,dx \\&= \int\frac{\sec^2(x) + \sec(x)\tan(x)}{\sec(x) + \tan(x)}\,dx
\end{align}


Making the substitution $u=\sec(x)+\tan(x)$ so

$du=\sec(x)\tan(x)+\sec^2(x)dx$ ,


\[
\int\sec(x)dx=\int\frac{du}{u}
\]

 \[
=\ln|u|+C
\]

\[
\ln\Big|\sec(x)+\tan(x)\Big|+C
\]


%-------------------------------------------------------------------------------
\subsubsection{More trigonometric combinations}
\label{sss:more-dab8}
%-------------------------------------------------------------------------------


> For the integrals $\int\sin(nx)\cos(mx)dx$ or $\int\sin(nx)\sin(mx)dx$

> or $\int\cos(nx)\cos(mx)dx$ use the identities

>

> -   $\sin(a)\cos(b)=\frac{\sin(a+b)+\sin(a-b)}{2}$

> -   $\sin(a)\sin(b)=\frac{\cos(a-b)-\cos(a+b)}{2}$

> -   $\cos(a)\cos(b)=\frac{\cos(a-b)+\cos(a+b)}{2}$


%-------------------------------------------------------------------------------
\subsubsection{Example 1}
\label{sss:exam-acc3}
%-------------------------------------------------------------------------------


Find $\int\sin(3x)\cos(5x)dx$ .


We can use the fact that $\sin(a)\cos(b)=\frac{\sin(a+b)+\sin(a-b)}{2}$, so


\[
\sin(3x)\cos(5x)=\frac{\sin(8x)+\sin(-2x)}{2}
\] Now use the oddness

property of $\sin(x)$ to simplify


\[
\sin(3x)\cos(5x)=\frac{\sin(8x)-\sin(2x)}{2}
\] And now we can

integrate


\[
\int\sin(3x)\cos(5x)dx
\]

\[
=\int\Big(\frac{\sin(8x)-\sin(2x)}{2}\Big)dx
\]

\[
=\frac{\cos(2x)}{4}-\frac{\cos(8x)}{16}+C
\]


%-------------------------------------------------------------------------------
\subsubsection{Example 2}
\label{sss:exam-5bdd}
%-------------------------------------------------------------------------------


Find\[
\int\sin(x)\sin(2x)dx
\] .


Using the identities


\[
\sin(x)\sin(2x)=\frac{\cos(-x)-\cos(3x)}{2}=\frac{\cos(x)-\cos(3x)}{2}
\]

Then


\[
\int\sin(x)\sin(2x)dx
\]

\[
=\frac{1}{2}\int\Big(\cos(x)-\cos(3x)\Big)dx
\]

\[
=\frac{\sin(x)}{2}-\frac{\sin(3x)}{6}+C
\]


%-------------------------------------------------------------------------------
\subsubsection{Resources}
\label{sss:reso-55b5}
%-------------------------------------------------------------------------------


\begin{itemize}
    \item \href{https://youtu.be/Sv-bSbAQeXc}{Trigonometric Integrals Involving Powers of Sine and Cosine - Part 1} by James Sousa
    \item \href{https://youtu.be/vPnuigP8I6I}{Trigonometric Integrals Involving Powers of Sine and Cosine - Part 2} by James Sousa
    \item \href{https://youtu.be/RGaDMqhOg8Y}{Trigonometric Integrals Involving Powers of Secant and Tangent - Part 1} by James Sousa
    \item \href{https://youtu.be/o8sIHlS17qc}{Trigonometric Integrals Involving Powers of Secant and Tangent - Part 1} by James Sousa
    \item \href{https://youtu.be/lTqnlihOC4o}{Trigonometric Integrals - Part 1 of 6} by patrickJMT
    \item \href{https://youtu.be/zyg9k1je7Fg}{Trigonometric Integrals - Part 2 of 6} by patrickJMT
    \item \href{https://youtu.be/BhJ4soojyAQ}{Trigonometric Integrals - Part 3 of 6} by patrickJMT
    \item \href{https://youtu.be/8CHHY-2Ctug}{Trigonometric Integrals - Part 4 of 6} by patrickJMT
    \item \href{https://youtu.be/QdNScjd5bno}{Trigonometric Integrals - Part 5 of 6} by patrickJMT
    \item \href{https://youtu.be/m3zG7c52QR4}{Trigonometric Integrals - Part 6 of 6} by patrickJMT
    \item \href{https://youtu.be/WYhyq_mTCZs}{Trigonometric integrals - sin\^mcos\^n, odd m} by Kriata King
    \item \href{https://youtu.be/RRDiT-djQPk}{Trigonometric integrals - sin\^mcos\^n, odd n} by Kriata King
    \item \href{https://youtu.be/rpbr2nH7lNY}{Trigonometric integrals - sin\^mcos\^n, m and n even} by Kriata King
    \item \href{https://youtu.be/CK2SfrzF_c4}{Integrals of trigonometric functions, tan\^msec\^n, even n} by Krista King
    \item \href{https://youtu.be/0TuZSSah5hc}{Integrals of trigonometric functions, tan\^msec\^n, odd m} by Krista King
\end{itemize}


%-------------------------------------------------------------------------------
\subsubsection{Licensing}
\label{sss:lice-7352}
%-------------------------------------------------------------------------------


Content obtained and/or adapted from:


\begin{itemize}
    \item \href{https://en.wikibooks.org/wiki/Calculus/Integration_techniques/Trigonometric_Integrals}{Trigonometric integrals, Wikibooks: Calculus/Integration techniques} under a CC BY-SA license
\end{itemize}

\subsection{Learning Objectives}

\begin{enumerate}
    \item Here are three student learning objectives (SLOs) related to the techniques for integrating functions involving powers of sine and cosine:
    \item \textbf{Integration of Powers of Sine and Cosine:}
    \item \textbf{Outcome:} Students will be able to evaluate integrals of the form $\int \cos^m(x) \sin^n(x) \, dx$ using appropriate trigonometric identities and substitution methods.
    \item \textbf{Details:} This includes identifying whether $m$ or $n$ is odd, applying substitutions such as $u = \sin(x)$ or $u = \cos(x)$, and using trigonometric identities to simplify and solve the integral. For instance, if $m$ is odd, students should substitute $u = \sin(x)$ and use the identity $\cos^2(x) = 1 - \sin^2(x)$. Conversely, if $n$ is odd, they should substitute $u = \cos(x)$ and use $\sin^2(x) = 1 - \cos^2(x)$.
    \item \textbf{Integration of Even Powers of Sine and Cosine:}
    \item \textbf{Outcome:} Students will be able to evaluate integrals of the form $\int \cos^m(x) \sin^n(x) \, dx$ when both $m$ and $n$ are even, using multiple angle identities.
    \item \textbf{Details:} This involves applying identities such as $\sin^2(x) = \frac{1 - \cos(2x)}{2}$ and $\cos^2(x) = \frac{1 + \cos(2x)}{2}$ to transform the integral into a form that can be integrated more easily. Students should be able to simplify the integrand using these identities and integrate using standard calculus techniques.
    \item \textbf{Integration of Powers of Tangent and Secant:}
    \item \textbf{Outcome:} Students will be able to evaluate integrals of the form $\int \tan^m(x) \sec^n(x) \, dx$ by applying substitution methods and using trigonometric identities.
    \item \textbf{Details:} This includes handling cases where $n$ is even, both $n$ and $m$ are odd, or $n$ is odd and $m$ is even. For instance, if $n$ is even, students should substitute $u = \tan(x)$ and use $\sec^2(x) = 1 + \tan^2(x)$. If both $n$ and $m$ are odd, they should substitute $u = \sec(x)$ and use $\tan^2(x) = \sec^2(x) - 1$. These SLOs focus on different integration techniques involving trigonometric functions and provide a structured approach to solving integrals based on the characteristics of the integrand.
\end{enumerate}

\subsection{Assessment Instrument}
\label{ssec:trig-fb0b}

Students use Pythagorean identities and half-angle formulas to reduce powers of $\sin$, $\cos$, $\tan$, and $\sec$ to integrable forms. The parity (odd vs.\ even) of each exponent determines the strategy.

\textbf{Example problems.}

\begin{enumerate}
    \item Evaluate $\int \sin^3 x\cos^2 x\,dx$ by saving one $\sin x$ factor and converting the remaining $\sin^2 x$ via a Pythagorean identity.
    \item Compute $\int \cos^4 x\,dx$ using the half-angle formula $\cos^2\theta = \tfrac{1+\cos 2\theta}{2}$.
    \item Find $\int \tan^3 x\sec^3 x\,dx$.
\end{enumerate}

%===============================================================================
\section{Trigonometric Substitution}
\label{sec:trigonometric-substitution}
%===============================================================================

The idea behind the trigonometric substitution is quite simple: to replace expressions involving square roots with expressions that involve standard trigonometric functions, but \textit{no square roots}. Integrals involving trigonometric functions are often easier to solve than integrals involving square roots.


Let us demonstrate this idea in practice. Consider the expression $\sqrt{1-x^2}$ . Probably the most basic trigonometric identity is $\sin^2(\theta)+\cos^2(\theta)=1$ for an arbitrary angle $\theta$ . If we replace $x$ in this expression by $\sin(\theta)$ , with the help of this trigonometric identity we see


\[
\sqrt{1-x^2}=\sqrt{1-\sin^2(\theta)}=\sqrt{\cos^2(\theta)}=\cos(\theta)
\]

Note that we could write $\theta=\arcsin(x)$ since we replaced $x^2$ with $\sin^2(\theta)$ .


We would like to mention that technically one should write the absolute value of $\cos(\theta)$ , in other words $\bigl|\cos(\theta)\bigr|$ as our final answer since $\sqrt{a^2}=|a|$ for all possible $a$ . But as long as we are careful about the domain of all possible $x$ and how $\cos(\theta)$ is used in the final computation, omitting the absolute value signs does not constitute a problem. However, we cannot directly interchange the simple expression $\cos(\theta)$ with the complicated $\sqrt{1-x^2}$ wherever it may appear, we must remember when integrating by substitution we need to take the derivative into account. That is we need to remember that $dx=\cos(\theta)d\theta$ , and to get a integral that only involves $\theta$ we need to also replace $dx$ by something in terms of $d\theta$ . Thus, if we see an integral of the form


\[
\int\sqrt{1-x^2}dx
\] we can rewrite it as


\[
\int\cos(\theta)\cos(\theta)d\theta=\int\cos^2(\theta)d\theta
\] Notice in the expression on the left that the first $\cos(\theta)$ comes from replacing the $\sqrt{1-x^2}$ and the $\cos(\theta)d\theta$ comes from substituting for the $dx$ .


Since $\cos^2(\theta)=\frac{1+\cos(2\theta)}{2}$ our original integral reduces to:


\[
\tfrac12\int d\theta+\tfrac12\int\cos(2\theta)d\theta
\]


These last two integrals are easily handled. For the first integral we get


\[
\tfrac12\int d\theta=\tfrac12\theta
\]


For the second integral we do a substitution, namely $u=2\theta,du=2d\theta$ to get:


\begin{align}
\tfrac{1}{2}\int\cos(2\theta)\,d\theta &= \tfrac{1}{2}\cdot\tfrac{1}{2}\int\cos(u)\,du \quad \text{(with } u = 2\theta \text{ and } du = 2\,d\theta\text{)} \\&= \tfrac{1}{4}\sin(u) \\&= \frac{\sin(2\theta)}{4}
\end{align}


Finally we see that:


\begin{align}
\int \cos^2(\theta)\,d\theta &= \int \frac{1 + \cos(2\theta)}{2}\,d\theta \\&= \frac{1}{2}\int (1 + \cos(2\theta))\,d\theta \\&= \frac{\theta}{2} + \frac{1}{2}\int \cos(2\theta)\,d\theta \\&= \frac{\theta}{2} + \frac{\sin(2\theta)}{4} + C
\end{align}


However, this is in terms of $\theta$ and not in terms of $x$ , so we must substitute back in order to rewrite the answer in terms of $x$ .


That is we worked out that:


\[
\sin(\theta)=x\qquad\cos(\theta)=\sqrt{1-x^2}\qquad\theta=\arcsin(x)
\]


So we arrive at our final answer


\[
\int\sqrt{1-x^2}dx=\frac{\arcsin(x)+x\sqrt{1-x^2}}{2}
\]


As you can see, even for a fairly harmless looking integral this technique can involve quite a lot of calculation. Often it is helpful to see if a simpler method will suffice before turning to trigonometric substitution. On the other hand, frequently in the case of integrands involving square roots, this is the most tractable way to solve the problem. We begin with giving some rules of thumb to help you decide which trigonometric substitutions might be helpful.


If the integrand contains a single factor of one of the forms $\sqrt{a^2-x^2}$, $\sqrt{a^2+x^2}$, or $\sqrt{x^2-a^2}$ we can try a trigonometric substitution.


> -   If the integrand contains $\sqrt{a^2-x^2}$ let $x=a\sin(\theta)$

>     and use the identity $1-\sin^2(\theta)=\cos^2(\theta)$ .

> -   If the integrand contains $\sqrt{a^2+x^2}$ let $x=a\tan(\theta)$

>     and use the identity $1+\tan^2(\theta)=\sec^2(\theta)$ .

> -   If the integrand contains $\sqrt{x^2-a^2}$ let $x=a\sec(\theta)$

>     and use the identity $\sec^2(\theta)-1=\tan^2(\theta)$ .


%-------------------------------------------------------------------------------
\subsubsection{Sine substitution}
\label{sss:sine-8011}
%-------------------------------------------------------------------------------


% [Image: This substitution is easily derived from a triangle, using the

Pythagorean

Theorem.]


Figure: This substitution is easily derived from a triangle, using the Pythagorean Theorem.


If the integrand contains a piece of the form $\sqrt{a^2-x^2}$ we use the substitution


\[
x=a\sin(\theta)\quad dx=a\cos(\theta)d\theta
\] This will transform the integrand to a trigonometric function. If the new integrand can't be integrated on sight then the tan-half-angle substitution described below will generally transform it into a more tractable algebraic integrand.


E.g., if the integrand is $\sqrt{1-x^2}$ ,


\begin{align}
\int\limits_0^1\sqrt{1-x^2}\,dx &= \int\limits_0^\frac{\pi}{2}\sqrt{1-\sin^2(\theta)}\cos(\theta)\,d\theta \\&= \int\limits_0^\frac{\pi}{2}\cos^2(\theta)\,d\theta \\&= \tfrac{1}{2}\int\limits_0^\frac{\pi}{2}\bigl(1+\cos(2\theta)\bigr)\,d\theta \\&= \frac{\pi}{4}
\end{align}


If the integrand is $\sqrt{\frac{1+x}{1-x}}$ , we can rewrite it as


\[
\sqrt{\frac{1+x}{1-x}}=\sqrt{\frac{1+x}{1+x}\cdot\frac{1+x}{1-x}}=\frac{1+x}{\sqrt{1-x^2}}
\]

Then we can make the substitution


\begin{align}
\int\limits_0^a\frac{1+x}{\sqrt{1-x^2}}\,dx 
&= \int\limits_0^\alpha\frac{1+\sin(\theta)}{\cos(\theta)}\cos(\theta)\,d\theta \quad\quad\quad 0 < a < 1 \\&= \int\limits_0^\alpha\bigl(1+\sin(\theta)\bigr)\,d\theta \quad\quad\quad \alpha = \text{arcsin}(a) \\&= \alpha + \left[-\cos(\theta)\right] _0^\alpha \\&= \alpha + 1 - \cos(\alpha) \\&= 1 + \text{arcsin}(a) - \sqrt{1-a^2}
\end{align}


%-------------------------------------------------------------------------------
\subsubsection{Tangent substitution}
\label{sss:tang-3508}
%-------------------------------------------------------------------------------


% [Image: This substitution is easily derived from a triangle, using the Pythagorean Theorem.]


Figure: This substitution is easily derived from a triangle, using the Pythagorean Theorem.


When the integrand contains a piece of the form $\sqrt{a^2+x^2}$ we use the substitution


\[
x=a\tan(\theta)\qquad\sqrt{x^2+a^2}=a\sec(\theta)\qquad dx=a\sec^2(\theta)d\theta
\]


E.g., if the integrand is $(x^2+a^2)^{-\frac32}$ then on making this substitution we find


\begin{align}
\int\limits_0^z(x^2+a^2)^{-\frac{3}{2}}\,dx 
&= a^{-2}\int\limits_0^\alpha \cos(\theta)\,d\theta \quad\quad\quad z>0 \\&= a^{-2}\left[\sin(\theta)\right] _0^\alpha \quad\quad\quad \alpha = \arctan\left(\tfrac{z}{a}\right) \\&= a^{-2}\sin(\alpha) \\&= a^{-2}\frac{\frac{z}{a}}{\sqrt{1+\frac{z^2}{a^2}}} \\&= \frac{z}{a^2\sqrt{a^2+z^2}}
\end{align}


If the integral is


\[
I=\int\limits_0^z\sqrt{x^2+a^2}\qquad z>0
\]


then on making this substitution we find


\begin{align}
I &= a^2 \int_0^\alpha \sec^3 \theta \, d\theta \quad\quad\quad \alpha = \tan^{-1} \left(\tfrac{z}{a}\right) \\&= a^2 \int_0^\alpha \sec \theta \, d\tan \theta \\&= a^2 \left[ \sec \theta \tan \theta \right] _0^\alpha - a^2 \int_0^\alpha \sec \theta \tan^2 \theta \, d\theta \\&= a^2 \sec \alpha \tan \alpha - a^2 \int_0^\alpha \sec^3 \theta \, d\theta  + a^2 \int_0^\alpha \sec \theta \, d\theta \\&= a^2 \sec \alpha \tan \alpha - I + a^2 \int_0^\alpha \sec \theta \, d\theta
\end{align}


After integrating by parts, and using trigonometric identities, we've ended up with an expression involving the original integral. In cases like this we must now rearrange the equation so that the original integral is on one side only


\begin{align}
I &= \tfrac{a^2}{2}\Bigg(\sec(\alpha)\tan(\alpha) + \int\limits_0^\alpha \sec(\theta)\,d\theta\Bigg) \\&= \tfrac{a^2}{2}\Bigg(\sec(\alpha)\tan(\alpha) + \bigg[\ln\Big(\sec(\theta) + \tan(\theta)\Big)\bigg] _0^\alpha\Bigg) \\&= \tfrac{a^2}{2}\bigg(\sec(\alpha)\tan(\alpha) + a^2\ln\Big(\sec(\alpha) + \tan(\alpha)\Big)\bigg) \\&= \tfrac{a^2}{2}\bigg(\tfrac{z}{a^2}\sqrt{a^2 + z^2} + \ln\left(\tfrac{z + \sqrt{a^2 + z^2}}{a}\right)\bigg) \\&= \tfrac{1}{2}\bigg(z\sqrt{z^2 + a^2} + a^2\ln\left(\tfrac{z + \sqrt{a^2 + z^2}}{a}\right)\bigg)
\end{align}


As we would expect from the integrand, this is approximately $\frac{z^2}{2}$ for large $z$ .


In some cases it is possible to do trigonometric substitution in cases when there is no $\sqrt{\ \ }$ appearing in the integral.


\textbf{Example}


\[
\int\frac{dx}{x^2+1}
\]


The denominator of this function is equal to

$\bigl(\sqrt{1+x^2}\bigr)^2$ . This suggests that we try to substitute

$x=\tan(u)$ and use the identity $1+\tan^2(u)=\sec^2(u)$ . With this substitution, we obtain that $dx=\sec^2(u)du$ and thus


\begin{align}
\int\frac{dx}{x^2+1} &= \int\frac{\sec^2(u)}{\tan^2(u)+1}\,du \\&= \int\frac{\sec^2(u)}{\sec^2(u)}\,du \\&= \int du \\&= u + C
\end{align}


Using the initial substitution $u=\arctan(x)$ gives


\[
\int\frac{dx}{x^2+1}=\arctan(x)+C
\]


%-------------------------------------------------------------------------------
\subsubsection{Secant substitution}
\label{sss:seca-a19a}
%-------------------------------------------------------------------------------


% [Image: This substitution is easily derived from a triangle, using the [Pythagorean Theorem]


Figure: This substitution is easily derived from a triangle, using the Pythagorean Theorem.


If the integrand contains a factor of the form

$\displaystyle\sqrt{x^2-a^2}$ we use the substitution


\[
x=a\sec(\theta)\quad dx=a\sec(\theta)\tan(\theta)d\theta\quad\sqrt{x^2-a^2}=a\tan(\theta)
\]


%-------------------------------------------------------------------------------
\subsubsection{Example 1}
\label{sss:exam-acc3}
%-------------------------------------------------------------------------------


Find $\int\limits_0^z \frac{\sqrt{x^2-1}}{x}dx$ .


\begin{align}
\int\limits_0^z \frac{\sqrt{x^2-1}}{x}\,dx
&= \int\limits_0^\alpha \frac{\tan(\theta)}{\sec(\theta)}\sec(\theta)\tan(\theta)\,d\theta \quad\quad\quad z > 1 \\&= \int\limits_0^\alpha \tan^2(\theta)\,d\theta \quad\quad\quad\quad\quad \alpha = \text{arcsec}(z) \\&= \bigg[\tan(\theta) - \theta\bigg] _0^\alpha \quad\quad\quad\quad \tan(\alpha) = \sqrt{\sec^2(\alpha) - 1} \\&= \tan(\alpha) - \alpha \quad\quad\quad\quad\quad \tan(\alpha) = \sqrt{z^2 - 1} \\&= \sqrt{z^2 - 1} - \text{arcsec}(z)
\end{align}


%-------------------------------------------------------------------------------
\subsubsection{Example 2}
\label{sss:exam-5bdd}
%-------------------------------------------------------------------------------


Find $\int\limits_1^z \frac{\sqrt{x^2-1}}{x^2}dx$ .


\begin{align}
\int\limits_1^z \frac{\sqrt{x^2-1}}{x^2}\,dx 
&= \int\limits_1^\alpha \frac{\tan(\theta)}{\sec^2(\theta)}\sec(\theta)\tan(\theta)\,d\theta \qquad z > 1 \\&= \int\limits_0^\alpha \frac{\sin^2(\theta)}{\cos(\theta)}\,d\theta \qquad\qquad\qquad \alpha = \text{arcsec}(z)
\end{align}


We can now integrate by parts


\begin{align}
\int\limits_1^z \frac{\sqrt{x^2-1}}{x^2}\,dx 
&= -\bigg[\tan(\theta)\cos(\theta)\bigg] _0^\alpha + \int\limits_0^\alpha \sec(\theta)\,d\theta \\&= -\sin(\alpha) + \bigg[\ln\bigl(\sec(\theta) + \tan(\theta)\bigr)\bigg] _0^\alpha \\&= \ln\bigl(\sec(\alpha) + \tan(\alpha)\bigr) - \sin(\alpha) \\&= \ln\bigl(z + \sqrt{z^2 - 1}\bigr) - \frac{\sqrt{z^2 - 1}}{z}
\end{align}

%-------------------------------------------------------------------------------
\subsubsection{Resources}
\label{sss:reso-55b5}
%-------------------------------------------------------------------------------

\begin{itemize}

\item \href{https://youtu.be/BVn8x7wCa3A}{Example 1: Integration Using Trigonometric Substitution}.  
Video by James Sousa.

\item \href{https://youtu.be/a-f63vZVuBQ}{Example 2: Integration Using Trigonometric Substitution}.  
Video by James Sousa.

\item \href{https://youtu.be/jZa4FM6l6pM}{Example 3: Integration Using Trigonometric Substitution}.  
Video by James Sousa.

\item \href{https://youtu.be/ObxsrPx44nw}{Example 4: Integration Using Trigonometric Substitution}.  
Video by James Sousa.

\item \href{https://youtu.be/OnT7Ub68FWg}{Example 5: Integration Using Trigonometric Substitution}.  
Video by James Sousa.

\item \href{https://youtu.be/3lC5AuCFK4c}{Trigonometric Substitution – Example 1}.  
Video by Patrick JMT.

\item \href{https://youtu.be/GVL4lHX6DgM}{Trigonometric Substitution – Example 2}.  
Video by Patrick JMT.

\item \href{https://youtu.be/yW6Odu0YHL0}{Trigonometric Substitution – Example 3 (Part 1)}.  
Video by Patrick JMT.

\item \href{https://youtu.be/4QNntd_KZbc}{Trigonometric Substitution – Example 4 (Rational Powers)}.  
Video by Patrick JMT.

\item \href{https://youtu.be/wuRNmNDLanc}{Trigonometric Substitution – More Examples}.  
Video by Patrick JMT.

\end{itemize}

%-------------------------------------------------------------------------------
\subsubsection{Licensing}
\label{sss:lice-7352}
%-------------------------------------------------------------------------------

Content obtained and/or adapted from:

\begin{itemize}

\item \href{https://en.wikibooks.org/wiki/Calculus/Integration_techniques/Trigonometric_Substitution}{Trigonometric Substitution, Wikibooks: Calculus/Integration Techniques}.  
Licensed under CC BY-SA.

\end{itemize}

\subsection{Learning Objectives}

\begin{enumerate}
    \item Here are three student learning objectives (SLOs) for understanding and applying trigonometric substitution in integrals:
    \item \textbf{Apply Sine Substitution:}
    \item \textbf{Outcome:} Students will be able to use the sine substitution method to simplify and solve integrals involving the square root of the form $\sqrt{a^2 - x^2}$.
    \item \textbf{Details:} This includes recognizing when to use the substitution $x = a \sin(\theta)$, transforming the integral into one involving trigonometric functions, and integrating using the resulting trigonometric identities. Students should be able to substitute back to obtain the final answer in terms of the original variable.
    \item \textbf{Utilize Tangent Substitution:}
    \item \textbf{Outcome:} Students will be able to apply the tangent substitution technique to integrals involving the square root of the form $\sqrt{x^2 + a^2}$.
    \item \textbf{Details:} This involves substituting $x = a \tan(\theta)$, simplifying the integral using the identity $\sqrt{x^2 + a^2} = a \sec(\theta)$, and integrating with respect to $\theta$. Students should also be able to convert the result back into the original variable if necessary.
    \item \textbf{Implement Secant Substitution:}
    \item \textbf{Outcome:} Students will be able to solve integrals involving the square root of the form $\sqrt{x^2 - a^2}$ using the secant substitution method.
    \item \textbf{Details:} This includes using the substitution $x = a \sec(\theta)$, transforming the integrand into a trigonometric function, and integrating using trigonometric identities. Students should be able to express the final result in terms of the original variable, demonstrating an understanding of the process and its application in various integral scenarios.
\end{enumerate}

\subsection{Assessment Instrument}
\label{ssec:trig-8a0f}

Assessment requires identifying the correct substitution ($x = a\sin\theta$, $a\tan\theta$, or $a\sec\theta$) based on the radical form, completing the square when necessary, and converting back to the original variable.

\textbf{Example problems.}

\begin{enumerate}
    \item Evaluate $\int\dfrac{dx}{\sqrt{9-x^2}}$ using $x = 3\sin\theta$.
    \item Compute $\int\dfrac{\sqrt{x^2-4}}{x}\,dx$ using $x = 2\sec\theta$.
    \item Find $\int\dfrac{dx}{(x^2+4)^2}$ using $x = 2\tan\theta$.
\end{enumerate}

%===============================================================================
\section{Partial Fractions}
\label{sec:partial-fractions}
%===============================================================================

Suppose we want to find $\int\frac{3x+1}{x^2+x}dx$ . One way to do this is to simplify the integrand by finding constants $A$ and $B$ so that


\[
\frac{3x+1}{x^2+x}=\frac{3x+1}{x(x+1)}=\frac{A}{x}+\frac{B}{x+1}
\] .

This can be done by cross multiplying the fraction which gives


\[
\frac{3x+1}{x(x+1)}=\frac{A(x+1)+Bx}{x(x+1)}
\] As both sides have the same denominator we must have


\[
3x+1=A(x+1)+Bx
\]

This is an equation for $x$ so it must hold whatever value $x$ is. If we put in $x=0$ we get $A=1$ and putting $x=-1$ gives $=-B=-2$ so $B=2$ . So we see that


\[
\frac{3x+1}{x^2+x}=\frac{1}{x}+\frac{2}{x+1}
\]


Returning to the original integral


\begin{align}
\int\frac{3x+1}{x^2+x}\,dx &= \int\frac{dx}{x} + \int\frac{2}{x+1}\,dx \\&= \int\frac{dx}{x} + 2\int\frac{dx}{x+1} \\&= \ln|x| + 2\ln\Big|x+1\Big| + C
\end{align}


Rewriting the integrand as a sum of simpler fractions has allowed us to reduce the initial integral to a sum of simpler integrals. In fact this method works to integrate any rational function.


%-------------------------------------------------------------------------------
\subsubsection{Method of Partial Fractions}
\label{sss:meth-ffe2}
%-------------------------------------------------------------------------------

To decompose the rational function $\frac{P(x)}{Q(x)}$:

\begin{itemize}
\item \textbf{Step 1:} Use long division (if necessary) to ensure that 
$\deg(P) < \deg(Q)$.
\item \textbf{Step 2:} Factor $Q(x)$ as far as possible.
\item \textbf{Step 3:} Write down the correct form for the partial fraction 
decomposition and solve for the constants.
\end{itemize}

\medskip

To factor $Q(x)$, express it as a product of linear factors $(ax+b)$ and 
irreducible quadratic factors $(ax^2+bx+c)$ with $b^2-4ac<0$.

\medskip

Some factors may be repeated. For example,
\begin{align}
Q(x) &= x^3 - 6x^2 + 9x \\
     &= x(x^2 - 6x + 9) \\
     &= x(x - 3)^2.
\end{align}

If instead
\begin{align}
Q(x) &= x^3 - 3x^2 + 2x \\
     &= x(x^2 - 3x + 2) \\
     &= x(x - 1)(x - 2),
\end{align}
then all factors are linear and distinct.

\medskip

In general, we express
\[
\frac{P(x)}{Q(x)} = 
\sum \frac{A}{(ax+b)^k} 
\quad \text{and} \quad
\frac{Ax+B}{(ax^2+bx+c)^k},
\]
depending on the factorization of $Q(x)$.

%-------------------------------------------------------------------------------
\subsubsection{$Q(x)$ is a Product of Distinct Linear Factors}
\label{sss:qxis-013e}
%-------------------------------------------------------------------------------

If
\[
Q(x) = (a_1x+b_1)(a_2x+b_2)\cdots(a_nx+b_n),
\]
with no repeated factors, then
\[
\frac{P(x)}{Q(x)} =
\frac{A_1}{a_1x+b_1} + \frac{A_2}{a_2x+b_2} + \cdots + \frac{A_n}{a_nx+b_n}.
\]

%-------------------------------------------------------------------------------
\subsubsection{Example Problem}
\label{sss:exam-accd}
%-------------------------------------------------------------------------------

Evaluate
\[
\int \frac{1+x^2}{(x+3)(x+5)(x+7)}\,dx.
\]

We write
\[
\frac{1+x^2}{(x+3)(x+5)(x+7)}
= \frac{A}{x+3} + \frac{B}{x+5} + \frac{C}{x+7}.
\]

Multiplying both sides:
\[
1+x^2 = A(x+5)(x+7) + B(x+3)(x+7) + C(x+3)(x+5).
\]

Substitute convenient values:

\[
x=-3:\quad 10 = 8A \Rightarrow A = \frac{5}{4}
\]

\[
x=-5:\quad 26 = -4B \Rightarrow B = -\frac{13}{2}
\]

\[
x=-7:\quad 50 = 8C \Rightarrow C = \frac{25}{4}
\]

Thus,
\[
\frac{1+x^2}{(x+3)(x+5)(x+7)}
= \frac{5}{4(x+3)} - \frac{13}{2(x+5)} + \frac{25}{4(x+7)}.
\]

Integrating:
\[
\int \frac{1+x^2}{(x+3)(x+5)(x+7)}\,dx
= \frac{5}{4}\ln|x+3| - \frac{13}{2}\ln|x+5| + \frac{25}{4}\ln|x+7| + C.
\]

%-------------------------------------------------------------------------------
\subsubsection{Exercises}
\label{sss:exer-pf-linear}
%-------------------------------------------------------------------------------

Evaluate the following using partial fractions:

\begin{enumerate}
\item $\displaystyle \int \frac{2x+11}{(x+6)(x+5)}\,dx$
\item $\displaystyle \int \frac{7x^2-5x+6}{(x-1)(x-3)(x-7)}\,dx$
\end{enumerate}

%-------------------------------------------------------------------------------
\subsubsection{$Q(x)$ Contains Repeated Linear Factors}
\label{sss:qxis-3eeb}
%-------------------------------------------------------------------------------

If $(ax+b)$ appears $k$ times in the factorization of $Q(x)$, then use
\[
\frac{A_1}{ax+b} + \frac{A_2}{(ax+b)^2} + \cdots + \frac{A_k}{(ax+b)^k}.
\]

%-------------------------------------------------------------------------------
\subsubsection{Example 2}
\label{sss:exam-5bdd}
%-------------------------------------------------------------------------------

Evaluate
\[
\int \frac{dx}{(x+1)(x+2)^2}.
\]

We write
\[
\frac{1}{(x+1)(x+2)^2}
= \frac{A}{x+1} + \frac{B}{x+2} + \frac{C}{(x+2)^2}.
\]

Multiplying:
\[
1 = A(x+2)^2 + B(x+1)(x+2) + C(x+1).
\]

Substitute values:

\[
x=-1:\quad 1 = A \Rightarrow A=1
\]

\[
x=-2:\quad 1 = -C \Rightarrow C=-1
\]

\[
x=0:\quad 1 = 4A + 2B + C = 4 + 2B -1 \Rightarrow B=-1
\]

Thus,
\[
\frac{1}{(x+1)(x+2)^2}
= \frac{1}{x+1} - \frac{1}{x+2} - \frac{1}{(x+2)^2}.
\]

Integrating:
\[
\int \frac{dx}{(x+1)(x+2)^2}
= \ln|x+1| - \ln|x+2| + \frac{1}{x+2} + C.
\]

Using logarithm properties:
\[
= \ln\left|\frac{x+1}{x+2}\right| + \frac{1}{x+2} + C.
\]


%-------------------------------------------------------------------------------
\subsubsection{Exercise}
\label{sss:exer-pf1}
%-------------------------------------------------------------------------------

\begin{enumerate}
\setcounter{enumi}{2}
\item Evaluate
\[
\int \frac{x^2 - x + 2}{x(x+2)^2}\,dx
\]
using the method of partial fractions.
\end{enumerate}

%-------------------------------------------------------------------------------
\subsubsection{$Q(x)$ Contains Non-Repeated Quadratic Factors}
\label{sss:qxco-886d}
%-------------------------------------------------------------------------------

If a quadratic factor $ax^2 + bx + c$ appears (and is not repeated), we use the form
\[
\frac{Ax + B}{ax^2 + bx + c}.
\]

%-------------------------------------------------------------------------------
\subsubsection{Exercises}
\label{sss:exer-pf2}
%-------------------------------------------------------------------------------

Evaluate the following using the method of partial fractions:

\begin{enumerate}
\setcounter{enumi}{3}

\item 
\[
\int \frac{2}{(x+2)(x^2+3)}\,dx
\]

\item 
\[
\int \frac{dx}{(x+2)(x^2+2)}
\]

\end{enumerate}

%-------------------------------------------------------------------------------
\subsubsection{$Q(x)$ Contains Repeated Quadratic Factors}
\label{sss:qxco-4e84}
%-------------------------------------------------------------------------------

If a quadratic factor $ax^2 + bx + c$ appears $k$ times, we use
\[
\frac{A_1x + B_1}{ax^2 + bx + c}
+ \frac{A_2x + B_2}{(ax^2 + bx + c)^2}
+ \cdots
+ \frac{A_kx + B_k}{(ax^2 + bx + c)^k}.
\]

%-------------------------------------------------------------------------------
\subsubsection{Exercise}
\label{sss:exer-pf3}
%-------------------------------------------------------------------------------

\begin{enumerate}
\setcounter{enumi}{5}

\item 
\[
\int \frac{dx}{(x-1)(x^2+1)^2}
\]

\end{enumerate}

%-------------------------------------------------------------------------------
\subsubsection{Exercise Solutions}
\label{sss:exer-eab6}
%-------------------------------------------------------------------------------

\begin{enumerate}

\item $\ln|x+6| + \ln|x+5| + C$

\item $\tfrac{2}{3}\ln|x-1| - \tfrac{27}{4}\ln|x-3| + \tfrac{157}{12}\ln|x-7| + C$

\item $\tfrac{1}{2}\ln|x(x+2)| + \frac{4}{x+2} + C$

\item $\ln\!\left(\sqrt[7]{\frac{(x+2)^2}{x^2+3}}\right)
+ \frac{4}{7\sqrt{3}} \arctan\!\left(\frac{x}{\sqrt{3}}\right) + C$

\item $\ln\!\left(\sqrt[12]{\frac{(x+2)^2}{x^2+2}}\right)
+ \frac{\sqrt{2}}{6} \arctan\!\left(\frac{x}{\sqrt{2}}\right) + C$

\item $\frac{1-x}{4(x^2+1)} + \tfrac{1}{8}\ln\!\left(\frac{(x-1)^2}{x^2+1}\right)
- \frac{1}{2}\arctan(x) + C$

\end{enumerate}

%-------------------------------------------------------------------------------
\subsubsection{Resources}
\label{sss:reso-55b5}
%-------------------------------------------------------------------------------

\begin{itemize}

\item \href{https://youtu.be/8MMALX7mUqQ}{Integration Using Partial Fraction Decomposition (Part 1)}.  
Video by James Sousa.

\item \href{https://youtu.be/BgVNLWLJmMs}{Integration Using Partial Fraction Decomposition (Part 2)}.  
Video by James Sousa.

\item \href{https://youtu.be/c2oLHtPA03U}{Partial Fraction Decomposition (Part 1: Linear Factors)}.  
Video by James Sousa.

\item \href{https://youtu.be/SQ4l8yDlBVE}{Partial Fraction Decomposition (Part 2)}.  
Video by James Sousa.

\item \href{https://youtu.be/HZTv4zCgEnA}{Partial Fraction Decomposition – Example 1}.  
Video by Patrick JMT.

\item \href{https://youtu.be/hnIs8uwAT3U}{Partial Fraction Decomposition – Example 2}.  
Video by Patrick JMT.

\item \href{https://youtu.be/lH9Mw65vAIA}{Partial Fraction Decomposition – Example 4}.  
Video by Patrick JMT.

\item \href{https://youtu.be/Slfrb8aEj08}{Partial Fraction Decomposition – Example 5}.  
Video by Patrick JMT.

\item \href{https://youtu.be/_YMG5mQTAOg}{Partial Fraction Decomposition – Example 6}.  
Video by Patrick JMT.

\item \href{https://youtu.be/7GcSOjLitF8}{Partial Fraction Decompositions}.  
Video by Patrick JMT.

\end{itemize}

%-------------------------------------------------------------------------------
\subsubsection{Licensing}
\label{sss:lice-7352}
%-------------------------------------------------------------------------------

Content obtained and/or adapted from:

\begin{itemize}

\item \href{https://en.wikibooks.org/wiki/Calculus/Integration_techniques/Partial_Fraction_Decomposition}{Partial Fraction Decomposition, Wikibooks: Calculus/Integration Techniques}.  
Licensed under CC BY-SA.

\end{itemize}

\subsection{Learning Objectives}

\begin{enumerate}
    \item Here are three Student Learning Objectives (SLOs) related to the method of partial fractions in integration:
    \item \textbf{Decompose Rational Functions:}
    \item \textbf{Outcome:} Students will be able to decompose a rational function into partial fractions.
    \item \textbf{Details:} This includes identifying the appropriate form of partial fractions based on the factorization of the denominator, whether it contains linear factors, quadratic factors, or repeated factors. Students will be able to find constants for each term in the decomposition and solve the resulting equations.
    \item \textbf{Integrate Decomposed Fractions:}
    \item \textbf{Outcome:} Students will be able to integrate rational functions by first decomposing them into partial fractions.
    \item \textbf{Details:} After decomposing a rational function, students should be able to integrate each term of the partial fraction decomposition separately. This involves applying integration techniques such as basic integration rules and logarithmic integration, and combining the results to obtain the final answer.
    \item \textbf{Apply Partial Fractions to Complex Integrals:}
    \item \textbf{Outcome:} Students will be able to apply the method of partial fractions to solve more complex integrals involving repeated and irreducible quadratic factors.
    \item \textbf{Details:} Students will be able to handle cases where the denominator of the rational function includes repeated linear factors or irreducible quadratic factors. They will use the appropriate forms for partial fractions and solve for the constants, integrating each term accordingly and combining the results into a final expression.
\end{enumerate}

\subsection{Assessment Instrument}
\label{ssec:part-8927}

Problems cover the full partial-fraction procedure: factoring the denominator, writing the decomposition form, solving for coefficients, and integrating each term. Repeated and irreducible quadratic factors are both included.

\textbf{Example problems.}

\begin{enumerate}
    \item Evaluate $\int\dfrac{3x+5}{(x+1)(x+2)}\,dx$ using partial fractions.
    \item Compute $\int\dfrac{2x^2+x-1}{x^3-x}\,dx$ after factoring the denominator.
    \item Find $\int\dfrac{x^2}{(x-1)(x^2+1)}\,dx$, handling the irreducible quadratic factor $x^2+1$ in the decomposition.
\end{enumerate}

%===============================================================================
\section{Improper Integrals}
\label{sec:improper-integrals}
%===============================================================================

The definition of a definite integral
\[
\int_a^b f(x)\,dx
\]
requires that the interval $[a,b]$ be finite and that $f$ be continuous 
on $[a,b]$. In this section, we extend the concept of integration to 
cases where one or both of these conditions fail.

\medskip

Such integrals are called \textbf{improper integrals}. These arise when:
\begin{itemize}
\item the limits of integration are infinite, or
\item the integrand has discontinuities on the interval.
\end{itemize}

%-------------------------------------------------------------------------------
\subsubsection{Improper Integrals with Infinite Limits}
\label{sss:impr-5c12}
%-------------------------------------------------------------------------------

Consider
\[
\int_1^\infty \frac{dx}{x^2}.
\]

We define this as
\[
\lim_{b\to\infty} \int_1^b \frac{dx}{x^2}
= \lim_{b\to\infty} \left(1 - \frac{1}{b}\right)
= 1.
\]

\medskip

\textbf{Definition (Infinite Limits).}

\begin{enumerate}
\item If $\int_a^b f(x)\,dx$ exists for all $b \ge a$, define
\[
\int_a^\infty f(x)\,dx = \lim_{b\to\infty} \int_a^b f(x)\,dx,
\]
provided the limit exists and is finite.

\item If $\int_a^b f(x)\,dx$ exists for all $a \le b$, define
\[
\int_{-\infty}^b f(x)\,dx = \lim_{a\to -\infty} \int_a^b f(x)\,dx.
\]

\item If both integrals converge, define
\[
\int_{-\infty}^\infty f(x)\,dx
= \int_{-\infty}^c f(x)\,dx + \int_c^\infty f(x)\,dx.
\]
\end{enumerate}

%-------------------------------------------------------------------------------
\subsubsection{Example: Convergent Integral}
\label{sss:exam-1179}
%-------------------------------------------------------------------------------

\[
\int_0^\infty e^{-x}\,dx
= \lim_{b\to\infty} \int_0^b e^{-x}\,dx
= \lim_{b\to\infty} \left[-e^{-x}\right]_0^b
= \lim_{b\to\infty} (-e^{-b} + 1)
= 1.
\]

%-------------------------------------------------------------------------------
\subsubsection{Example: Divergent Integral}
\label{sss:exam-7826}
%-------------------------------------------------------------------------------

\[
\int_1^\infty \frac{dx}{x}
= \lim_{b\to\infty} \int_1^b \frac{dx}{x}
= \lim_{b\to\infty} \ln b
= \infty.
\]

Hence, the integral diverges.

%-------------------------------------------------------------------------------
\subsubsection{Example: Improper Integral}
\label{sss:exam-d573}
%-------------------------------------------------------------------------------

Evaluate
\[
\int_0^\infty x^2 e^{-x}\,dx.
\]

Using integration by parts twice:
\begin{align}
\int_0^b x^2 e^{-x}\,dx
&= \left[-x^2 e^{-x}\right]_0^b + 2\int_0^b x e^{-x}\,dx \\
&= -b^2 e^{-b} + 2\left(\left[-x e^{-x}\right]_0^b + \int_0^b e^{-x}\,dx\right) \\
&= -b^2 e^{-b} + 2\left(-b e^{-b} + \left[-e^{-x}\right]_0^b\right) \\
&= -b^2 e^{-b} + 2\left(-b e^{-b} - e^{-b} + 1\right).
\end{align}

Taking the limit as $b \to \infty$:
\[
\int_0^\infty x^2 e^{-x}\,dx = 2.
\]

%-------------------------------------------------------------------------------
\subsubsection{Example: $p$-Integrals}
\label{sss:exam-6fd6}
%-------------------------------------------------------------------------------

Show
\[
\int_1^\infty \frac{dx}{x^p}
=
\begin{cases}
\frac{1}{p-1}, & p > 1, \\
\text{diverges}, & p \le 1.
\end{cases}
\]

For $p \ne 1$:
\begin{align}
\int_1^\infty x^{-p}\,dx
&= \lim_{b\to\infty} \left[\frac{x^{1-p}}{1-p}\right]_1^b \\
&= \lim_{b\to\infty} \frac{b^{1-p} - 1}{1-p}.
\end{align}

This converges iff $p>1$.

%-------------------------------------------------------------------------------
\subsubsection{Improper Integrals with Discontinuities}
\label{sss:impr-312b}
%-------------------------------------------------------------------------------

\textbf{Definition (Discontinuities).}

\begin{enumerate}
\item If $f$ is continuous on $[a,b)$ but discontinuous at $b$, define
\[
\int_a^b f(x)\,dx = \lim_{c\to b^-} \int_a^c f(x)\,dx.
\]

\item If $f$ is continuous on $(a,b]$ but discontinuous at $a$, define
\[
\int_a^b f(x)\,dx = \lim_{c\to a^+} \int_c^b f(x)\,dx.
\]

\item If $f$ has a discontinuity at $c \in (a,b)$, define
\[
\int_a^b f(x)\,dx
= \int_a^c f(x)\,dx + \int_c^b f(x)\,dx,
\]
provided both integrals converge.
\end{enumerate}

%-------------------------------------------------------------------------------
\subsubsection{Example 1}
\label{sss:exam-acc3}
%-------------------------------------------------------------------------------

Show that
\[
\int_0^1 \frac{dx}{x^p}
=
\begin{cases}
\dfrac{1}{1-p}, & \text{if } p<1, \\[4pt]
\text{diverges}, & \text{if } p\ge 1.
\end{cases}
\]

If $p\ne 1$, then
\begin{align}
\int_0^1 \frac{dx}{x^p}
&= \lim_{a\to 0^+}\int_a^1 x^{-p}\,dx \\
&= \lim_{a\to 0^+}\left(\frac{x^{1-p}}{1-p}\right)\Bigg|_a^1 \\
&= \frac{1}{1-p}\lim_{a\to 0^+}\bigl(1-a^{1-p}\bigr).
\end{align}

If $p<1$, then $1-p>0$, so $a^{1-p}\to 0$ as $a\to 0^+$, and hence
\[
\int_0^1 \frac{dx}{x^p} = \frac{1}{1-p}.
\]

If $p>1$, then $1-p<0$, so $a^{1-p}\to\infty$ as $a\to 0^+$, and the integral diverges.

Notice that we had to assume $p\ne 1$ to avoid dividing by $0$. So we do the case $p=1$ separately:
\begin{align}
\int_0^1 \frac{dx}{x}
&= \lim_{a\to 0^+}\int_a^1 \frac{dx}{x} \\
&= \lim_{a\to 0^+}\Bigl[\ln|x|\Bigr]_a^1 \\
&= \lim_{a\to 0^+}\bigl(0-\ln(a)\bigr) \\
&= \infty.
\end{align}

Thus, the integral diverges when $p=1$ as well.

%-------------------------------------------------------------------------------
\subsubsection{Example 2}
\label{sss:exam-5bdd}
%-------------------------------------------------------------------------------

The integral
\[
\int_{-1}^3 \frac{dx}{x-2}
\]
is improper because the integrand is not continuous at $x=2$. Had we not noticed this, we might have been tempted to apply the Fundamental Theorem of Calculus and conclude that it equals
\[
\ln|x-2|\Big|_{-1}^3 = \ln(1)-\ln(3) = -\ln(3),
\]
which is not correct.

In fact, the integral diverges, since
\begin{align}
\int_{-1}^3 \frac{dx}{x-2}
&= \lim_{b\to 2^-}\int_{-1}^b \frac{dx}{x-2}
 + \lim_{a\to 2^+}\int_a^3 \frac{dx}{x-2} \\
&= \lim_{b\to 2^-}\ln|x-2|\Big|_{-1}^b
 + \lim_{a\to 2^+}\ln|x-2|\Big|_a^3 \\
&= \lim_{b\to 2^-}\bigl(\ln(2-b)-\ln(3)\bigr)
 + \lim_{a\to 2^+}\bigl(\ln(1)-\ln(a-2)\bigr).
\end{align}

Now
\[
\lim_{b\to 2^-}\ln(2-b) = -\infty
\quad\text{and}\quad
\lim_{a\to 2^+}[-\ln(a-2)] = \infty,
\]
so the improper integral diverges.

\medskip

\textbf{Definition: Improper integrals with a finite number of discontinuities}

Suppose $f$ is continuous on $[a,b]$ except at points
\[
c_1<c_2<\cdots<c_n
\]
in $[a,b]$. We define
\[
\int_a^b f(x)\,dx
=
\int_a^{c_1} f(x)\,dx
+\int_{c_1}^{c_2} f(x)\,dx
+\cdots
+\int_{c_n}^b f(x)\,dx,
\]
provided that each improper integral on the right converges.

Notice that, by combining this definition with the definition for improper integrals with infinite endpoints, we can define the integral of a function with a finite number of discontinuities and one or more infinite endpoints.

%-------------------------------------------------------------------------------
\subsubsection{Comparison Test}
\label{sss:comp-edc1}
%-------------------------------------------------------------------------------

There are improper integrals that cannot easily be evaluated directly. However, it may still be possible to determine whether they converge or diverge by comparing them to integrals whose behavior is already known.

\medskip

\textbf{Theorem (Comparison Test).} Let $f$ and $g$ be continuous functions defined for all $x\ge a$.

\begin{enumerate}
\item Suppose
\[
g(x)\ge f(x)\ge 0
\quad\text{for all } x\ge a.
\]
If
\[
\int_a^\infty g(x)\,dx
\]
converges, then
\[
\int_a^\infty f(x)\,dx
\]
also converges.

\item Suppose
\[
f(x)\ge g(x)\ge 0
\quad\text{for all } x\ge a.
\]
If
\[
\int_a^\infty g(x)\,dx
\]
diverges, then
\[
\int_a^\infty f(x)\,dx
\]
also diverges.
\end{enumerate}

A similar theorem holds for improper integrals of the form
\[
\int_{-\infty}^b f(x)\,dx
\]
and for improper integrals with discontinuities.

%-------------------------------------------------------------------------------
\subsubsection{Example: Use of the Comparison Test to Show Convergence}
\label{sss:exam-0de4}
%-------------------------------------------------------------------------------

Show that
\[
\int_1^\infty \frac{\sin(x)+2}{x^2}\,dx
\]
converges.

For all $x$, we know that
\[
-1\le \sin(x)\le 1,
\]
so
\[
1\le \sin(x)+2\le 3.
\]
Therefore,
\[
0\le \frac{\sin(x)+2}{x^2}\le \frac{3}{x^2}.
\]

We have already seen that
\[
\int_1^\infty \frac{3}{x^2}\,dx
=
3\int_1^\infty \frac{dx}{x^2}
\]
converges. Hence, by the comparison test,
\[
\int_1^\infty \frac{\sin(x)+2}{x^2}\,dx
\]
also converges.

%-------------------------------------------------------------------------------
\subsubsection{Example: Use of the Comparison Test to Show Divergence}
\label{sss:exam-3ee7}
%-------------------------------------------------------------------------------

Show that
\[
\int_1^\infty \frac{\sin(x)+2}{x}\,dx
\]
diverges.

As before, we know that
\[
\sin(x)+2\ge 1
\quad\text{for all }x.
\]
Thus,
\[
\frac{\sin(x)+2}{x}\ge \frac{1}{x}\ge 0.
\]

Since
\[
\int_1^\infty \frac{dx}{x}
\]
diverges, the comparison test implies that
\[
\int_1^\infty \frac{\sin(x)+2}{x}\,dx
\]
also diverges.

%-------------------------------------------------------------------------------
\subsubsection{An Extension of the Comparison Theorem}
\label{sss:anex-6d45}
%-------------------------------------------------------------------------------

To apply the comparison theorem, it is not necessary that
\[
g(x)\ge f(x)\ge 0
\quad\text{for all }x\ge a.
\]
What is actually needed is that this inequality hold for all sufficiently large $x$; that is, there exists some number $c$ such that
\[
g(x)\ge f(x)\ge 0
\quad\text{for all }x\ge c.
\]

Indeed,
\[
\int_a^\infty f(x)\,dx
=
\int_a^c f(x)\,dx
+
\int_c^\infty f(x)\,dx.
\]
Thus, the convergence or divergence depends only on the tail integral, and the comparison test may be applied to
\[
\int_c^\infty f(x)\,dx.
\]

%-------------------------------------------------------------------------------
\subsubsection{Example}
\label{sss:exam-1a79}
%-------------------------------------------------------------------------------

Show that
\[
\int_1^\infty \sqrt{\frac{x^7}{e^{3x}}}\,dx
\]
converges.

We first rewrite the integrand:
\[
\sqrt{\frac{x^7}{e^{3x}}}
=
x^{7/2}e^{-3x/2}
=
x^{7/2}e^{-x/2}e^{-x}.
\]

Now,
\[
\lim_{x\to\infty} x^{7/2}e^{-x/2}=0,
\]
so for sufficiently large $x$ we have
\[
x^{7/2}e^{-x/2}\le 1.
\]
Therefore, for all sufficiently large $x$,
\[
x^{7/2}e^{-3x/2}
=
x^{7/2}e^{-x/2}e^{-x}
\le e^{-x}.
\]

Since
\[
\int_1^\infty e^{-x}\,dx
\]
converges, the comparison test implies that
\[
\int_1^\infty \sqrt{\frac{x^7}{e^{3x}}}\,dx
\]
also converges.


%-------------------------------------------------------------------------------
\subsubsection{Resources}
\label{sss:reso-55b5}
%-------------------------------------------------------------------------------

\begin{itemize}

\item \href{https://youtu.be/F19y9IiH-cM}{Improper Integrals}.  
Video by James Sousa.

\item \href{https://youtu.be/qWMJpYajOww}{Example 1: Improper Integral (Infinite Interval $(-\infty,\infty)$)}.  
Video by James Sousa.

\item \href{https://youtu.be/q-l1WbnjKhw}{Example 2: Improper Integral (Infinite Interval $(-\infty, c)$)}.  
Video by James Sousa.

\item \href{https://youtu.be/8cc720Y_UWU}{Example 3: Improper Integral (Infinite Interval $(-\infty,\infty)$)}.  
Video by James Sousa.

\item \href{https://youtu.be/85-HNJyuyrU}{Improper Integral – Basic Idea and Example}.  
Video by Patrick JMT.

\item \href{https://youtu.be/KvC4XyayEC0}{Improper Integral – More Complicated Example}.  
Video by Patrick JMT.

\item \href{https://youtu.be/f6cGotvktxs}{Improper Integral – Infinite Limits}.  
Video by Patrick JMT.

\item \href{https://youtu.be/GZ-vy0bApWc}{What Makes an Integral Improper?}.  
Video by Krista King.

\item \href{https://youtu.be/DUowajlKIKA}{Improper Integrals}.  
Video by Krista King.

\item \href{https://youtu.be/266JzwOkSJs}{Improper Integral Example 2 (Part 1)}.  
Video by Krista King.

\item \href{https://youtu.be/6fVgBQMvVJY}{Improper Integral Example 2 (Part 2)}.  
Video by Krista King.

\end{itemize}

%-------------------------------------------------------------------------------
\subsubsection{Licensing}
\label{sss:lice-7352}
%-------------------------------------------------------------------------------

Content obtained and/or adapted from:

\begin{itemize}

\item \href{https://en.wikibooks.org/wiki/Calculus/Improper_Integrals}{Improper Integrals, Wikibooks: Calculus}.  
Licensed under CC BY-SA.

\end{itemize}


\subsection{Learning Objectives}

\begin{enumerate}
    \item Here are three Student Learning Objectives (SLOs) related to the definition and evaluation of improper integrals, as outlined in the provided text:
    \item \textbf{Understand the Definition of Improper Integrals with Infinite Limits:}
    \item \textbf{Outcome:} Students will be able to explain the definition of improper integrals with infinite limits of integration and determine whether such integrals converge or diverge.
    \item \textbf{Details:} This involves understanding how to evaluate integrals where the limits of integration are infinite by taking limits and applying convergence criteria. Students should be able to apply these concepts to examples such as \[
\int\limits_1^\infty \frac{dx}{x^2}
\] and \[
\int\limits_1^\infty \frac{dx}{x}
\] to determine their convergence or divergence.
    \item \textbf{Evaluate Improper Integrals with Finite Discontinuities:}
    \item \textbf{Outcome:} Students will be able to evaluate improper integrals that have finite discontinuities by applying the appropriate limits and combining integrals across discontinuities.
    \item \textbf{Details:} This includes calculating integrals where the function has discontinuities within the interval, such as \[
\int\limits _{-1}^3 \frac{dx}{x-2}
\], and ensuring that each part of the integral converges individually. Students should be able to use the definition for improper integrals with finite discontinuities to evaluate these cases correctly.
    \item \textbf{Apply the Comparison Test for Convergence and Divergence:}
    \item \textbf{Outcome:} Students will be able to use the comparison test to determine the convergence or divergence of improper integrals by comparing them with known convergent or divergent integrals.
    \item \textbf{Details:} This includes using inequalities to compare integrands, such as in the examples of \[
\int\limits_1^\infty \frac{\sin(x)+2}{x^2}dx
\] and \[
\int\limits_1^\infty \frac{\sin(x)+2}{x}dx
\]. Students should be able to apply the comparison test effectively to infer the behavior of more complex integrals based on simpler, known cases.
\end{enumerate}

\subsection{Assessment Instrument}
\label{ssec:impr-a882}

Students replace infinite limits (or discontinuities) with a parameter, evaluate the resulting proper integral, and take the limit. Convergence must be stated with justification; divergence must be proved rather than assumed.

\textbf{Example problems.}

\begin{enumerate}
    \item Determine whether $\int_1^{\infty}\dfrac{1}{x^2}\,dx$ converges or diverges, and if it converges, find its value.
    \item Evaluate $\int_0^1\dfrac{1}{\sqrt{x}}\,dx$, handling the discontinuity at $x=0$.
    \item Determine the convergence of $\int_{-\infty}^{\infty}x\,e^{-x^2}\,dx$ by splitting at $0$ and evaluating each part separately.
\end{enumerate}

%===============================================================================
\section{Separable Equations}
\label{sec:separable-equations}
%===============================================================================

This wasn't listed in the Wiki under MAT 1224.

%===============================================================================
\section{Exponential Growth and Decay.}
\label{sec:expo-fa23}
%===============================================================================

This wasn't listed in the Wiki under MAT 1224.

%===============================================================================
\section{The logistic equation}
\label{sec:thel-e63c}
%===============================================================================

This wasn't listed in the Wiki under MAT 1224.

%===============================================================================
\section{First order linear equations}
\label{sec:firs-9bf7}
%===============================================================================

This wasn't listed in the Wiki under MAT 1224.

%===============================================================================
\section{Sequences}
\label{sec:sequences}
%===============================================================================


%-------------------------------------------------------------------------------
\subsubsection{Finite Sequences}
\label{sss:fini-2cff}
%-------------------------------------------------------------------------------


Definition of a Sequence


A \textbf{sequence} is an ordered collection of terms in which repetition is allowed. The number of terms in a sequence is called the \textbf{length} of the sequence.


Sequences are often denoted by brackets like

$\left \\{2,3,3,4,4 \right \\}$. Furthermore if we have a sequence $a$ such that $a = \left \\{1,2,3,4,5 \right \\}$ then $a_1 = 1, a_2=2, a_3=3...$.


The subscript must be a non-negative integer. Also notice that $n$ starts from one and counts up.


We can describe the terms in this sequence with a formula $a_n=n$ for all non-negative integers $n < 6$. So under this definition $a_6$ is not defined, and indeed $a_6$ is not in the sequence.


%-------------------------------------------------------------------------------
\subsubsection{Infinite sequences}
\label{sss:infi-58fd}
%-------------------------------------------------------------------------------


Definition of an infinite sequence


An \textbf{infinite sequence} is a sequence with an infinite number of elements.


Infinite sequences have infinite terms. For such a sequence, we can again give a formula for any term in the sequence. For our previous sequence $a$, we can say $a_n = n$ for all non-negative integers $n$. This sequence could also be denoted as $\left \\{0,1,2,3,4,5,... \right \\}$ where the period of ellipses implies that this sequence is infinite.


%-------------------------------------------------------------------------------
\subsubsection{Discrete Functions}
\label{sss:disc-f4b3}
%-------------------------------------------------------------------------------


Earlier, we defined the members of the infinite sequence $a$ as

$a_n = n$ for all non-negative integers $n$. This is known as a \textbf{discrete function}, \textbf{discrete definition}, or \textbf{explicit definition}. A discrete function is any function whose domain is not the set of all real or imaginary numbers, but is instead a smaller, countable set like the set of all integers or the set of all rational numbers. Note that a set differs from a sequence, but that is beyond the scope of this discussion.


Discrete functions only take ``countable'', discrete domains. The set of all integers is countable, because there are not infinitely many values between two values in the set; there is no extra value between 2 and 1, as 1.5 is not an integer and is not contained in the set. Also note that given a discrete function or explicit definition, as long as the domain is discrete, the range must also be discrete. This means that if the input of a discrete function is countable, the output must also be countable.


%-------------------------------------------------------------------------------
\subsubsection{Example 1}
\label{sss:exam-acc3}
%-------------------------------------------------------------------------------


$q_n = n + 1$


$q = \left \\{2, 3, 4, 5, 6... \right \\}$


This is known as an arithmetic sequence. These will be discussed later.


%-------------------------------------------------------------------------------
\subsubsection{Example 2}
\label{sss:exam-5bdd}
%-------------------------------------------------------------------------------


$c_n = \cos(n-1)$


$c = \left \\{1, 0.5403..., -0.4161..., -0.9899..., 0.2836...,...\right \\}$


This result may be interesting: a sequence does not need to be a collection of integers, indeed it can be any collection, as long as it is countable. Here, we are simply taking the cosine of all integers, and any discrete function must have both a discrete domain \textit{and} range.


%-------------------------------------------------------------------------------
\subsubsection{Recursive Functions}
\label{sss:recu-4e24}
%-------------------------------------------------------------------------------


\textbf{Recursive functions}, \textbf{recursive formulas}, or **recursive

definitions** are formulas in which $a _{n}$ is defined in terms of $a _{n-1}$. Knowing any term in a recursively defined sequence requires you to know all the terms before it, which means you must know the first term, sometimes denoted $a_0$ or $a_1$. The first term must be defined in order to have a proper recursive sequence; it cannot be assumed that the first term is 1.


Sometimes, one can have a sequence that is \textit{necessarily} defined by a recursive function. For instance, the recursively defined sequence $u _{n+1} = \cos(u_n), a_1 = 1$. This sequence cannot be expressed any other ``easy'' way and in this kind of situation it is best to use the recursive definition.


%-------------------------------------------------------------------------------
\subsubsection{Example 1}
\label{sss:exam-acc3}
%-------------------------------------------------------------------------------


The sequence


$p _{n+1} = p _{n} + 1, p_1 = 2$


$p = \left \\{2, 3, 4, 5... \right \\}$


is the same arithmetic sequence mentioned earlier. However, this time it uses a recursive definition which is essentially the same.


%-------------------------------------------------------------------------------
\subsubsection{Example 2}
\label{sss:exam-5bdd}
%-------------------------------------------------------------------------------


This is the sequence of cosine mentioned earlier:


$u _{n+1} = \cos(u_n), a_1 = 1$


$u = \left \\{0.5403..., -0.9111..., 0.6128..., 0.8180..., ... \right \\}$


%-------------------------------------------------------------------------------
\subsubsection{Example 3}
\label{sss:exam-ad87}
%-------------------------------------------------------------------------------


$s_n = 3\times s _{n-1}$


Notice that this time, instead of saying $s _{n+1} = ...$, we defined $s_n$ in terms of $s _{n-1}$. This definition is still valid.


%-------------------------------------------------------------------------------
\subsubsection{Resources}
\label{sss:reso-55b5}
%-------------------------------------------------------------------------------


%-------------------------------------------------------------------------------
\subsubsection{Videos}
\label{sss:vide-cce3}
%-------------------------------------------------------------------------------


\href{https://youtu.be/p-rc5mTDt9E}{Introduction to Sequences} by James Sousa


\href{https://youtu.be/Kxh7yJC9Jr0}{What is a Sequence? Basic Sequence Info}

by patrickJMT


\href{https://youtu.be/XdoqvWxTVEw}{Calculating the first terms of the

sequence} by Krista King


\href{https://youtu.be/0aCpT3dalh8}{Ex: Find the General Formula For a Sequence in Fraction Form

(Arithmetic/Geometric)} by James Sousa


\href{https://youtu.be/zoP3UnulRcA}{Geometric Sequences: A Formula for the' n - th ' Term , Part

1} by patrickJMT


\href{https://youtu.be/INjRH9h3oxA}{Finding the nth Term of the Sequence} by

Krista King


\href{https://youtu.be/hc64LUtPjP0}{Limits of a Sequence} by James Sousa


\href{https://youtu.be/9K1xx6wfN-U}{Sequences - Examples showing convergence or

divergence} by patrickJMT


\href{https://youtu.be/0piu2uP4zr4}{Finding the Limit of a Sequence, 3 more

examples} by patrickJMT


\href{https://youtu.be/BknXGSoCNB0}{Does the sequence converge or diverge?}

by Krista King


\href{https://youtu.be/bvPxNhs-yrI}{Intro to Monotonic and Bounded Sequences, Ex

1} by patrickJMT


\href{https://youtu.be/npboNfgQff8}{Increasing, decreasing and not monotonic

sequences} by Krista King


\href{https://youtu.be/tHy3TXmZpF0}{Monotonic Sequences and Bounded

Sequences} by The Organic Chemistry Tutor


\href{https://youtu.be/RjsyEWDEQe0}{Ex: Finding Terms in a Sequence Given a Recursive

Formula} by James Sousa


\href{https://youtu.be/IEyojX9jzv0}{Recursive Sequences} by patrickJMT


\href{https://youtu.be/IFHZQ6MaG6w}{Recursive Formulas For Sequences} by The

Organic Chemistry Tutor


%-------------------------------------------------------------------------------
\subsubsection{Licensing}
\label{sss:lice-7352}
%-------------------------------------------------------------------------------


Content obtained and/or adapted from:


\begin{itemize}
    \item \href{https://en.wikibooks.org/wiki/Calculus/Definition_of_a_Sequence}{Definition of a sequence, Wikibooks:
\end{itemize}
    Calculus}

    under a CC BY-SA license

\subsection{Learning Objectives}

\begin{enumerate}
    \item Here are three Student Learning Objectives (SLOs) for the topic of finite sequences, infinite sequences, and discrete functions:
    \item \textbf{Understand and Define Finite Sequences:}
    \item \textbf{Outcome:} Students will be able to define and recognize finite sequences, including their length and notation.
    \item \textbf{Details:} This involves understanding that a finite sequence is an ordered list of terms with a specific length. Students should be able to identify and list the terms of a sequence given its definition and understand the concept of indexing, where subscripts start from one and count up.
    \item \textbf{Analyze and Describe Infinite Sequences:}
    \item \textbf{Outcome:} Students will be able to describe infinite sequences and identify the general formula for any term in such sequences.
    \item \textbf{Details:} This includes recognizing that an infinite sequence has an endless number of terms and being able to express its general term formula. Students should be able to differentiate between finite and infinite sequences and understand how the concept of infinity affects sequence notation and definitions.
    \item \textbf{Interpret and Use Discrete Functions:}
    \item \textbf{Outcome:} Students will be able to interpret discrete functions and apply them to sequences with countable domains and ranges.
    \item \textbf{Details:} This involves understanding that discrete functions are defined over countable sets such as integers and can be used to generate sequences. Students should be able to distinguish discrete functions from continuous functions and recognize that both the domain and range of a discrete function are countable. These SLOs should help guide the learning process and assessment for students studying sequences and discrete functions.
\end{enumerate}

\subsection{Assessment Instrument}
\label{ssec:sequ-6e01}

Assessment covers the definition of convergence, standard limit techniques (Squeeze Theorem, L\'H\^{o}pital, ratio comparison), and the Monotone Convergence Theorem. Students must prove convergence or divergence explicitly.

\textbf{Example problems.}

\begin{enumerate}
    \item Determine whether $a_n = \dfrac{3n^2+1}{2n^2-n}$ converges, and if so, find its limit.
    \item Use the Squeeze Theorem to show that $a_n = \dfrac{\cos n}{n}$ converges to $0$.
    \item Show that $a_n = \dfrac{n}{e^n}$ converges to $0$ using L\'H\^{o}pital\'s Rule on the corresponding function.
\end{enumerate}

%===============================================================================
\section{Series}
\label{sec:series}
%===============================================================================

A \textit{series} is the sum of a sequence of terms. An \textit{infinite series} is

the sum of an infinite number of terms (the actual sum of the series

need not be infinite, as we will see below).


While there is usually some pattern in the terms being added,

technically a series can be the sum of \textit{any} sequence of terms:


\[
1+2+\pi+\cos\tfrac{\pi}{12}+e+\ldots
\] The problem with the above

series is, it is not at all clear what the next term of the series would

be, nor any subsequent ones. To say anything useful about the series,

there needs to be a clear pattern. Two such patterns seen in sequences

that also appear in the study of series are called \textit{arithmetic} and

\textit{geometric}.


An \textit{arithmetic series} is the sum of a sequence of terms having a

\textit{common difference} (i.e., the difference between consecutive terms is

always the same). For example,


\[
1+4+7+10+13+\dots
\] is an arithmetic series with common difference 3,

since $a_2-a_1=3$, $a_3-a_2=3$, and so forth. Unlike the first series

above, this series has a clear pattern. It is obvious that the next term

of the series is 16, and then 19, and so on.


A \textit{geometric series} is the sum of a sequence of terms having a *common

ratio* (i.e., the ratio between consecutive terms is always the same).

For example,


\[
1+\frac{1}{2}+\frac{1}{4}+\frac{1}{8}+\frac{1}{16}+\ldots
\] is a

geometric series with common ratio $\tfrac{1}{2}$, since

$a_2/a_1=\tfrac{1}{2}$, $a_3/a_2=\tfrac{1}{2}$, and so forth. As before,

there is a definite pattern, and the next term of the series is clearly

$\tfrac{1}{32}$.


%-------------------------------------------------------------------------------
\subsubsection{Summation notation}
\label{sss:summ-2e90}
%-------------------------------------------------------------------------------


\textit{Summation notation} provides a compact way of writing an infinite

series. The arithmetic series above can be written


\[
\sum _{n=1}^{\infty} [1+3(n-1)]
\] or, more simply


\[
\sum _{n=1}^{\infty} (3n-2).
\] The symbol on the left is a Sigma (the

Greek letter for capital ``S'', standing for ``Sum''); it is also known

as a \textit{summation symbol}. Along with the expressions above and below the

Sigma, it is read:


$\quad$ ``The sum, as $n$ goes from 1 to infinity, of\...''


The algebraic expression on the right is the \textit{n}th term of the series,

and is often generically called "$a_n$, as is also done with

sequences. This means that the sum could have been written a bit more

verbosely (and perhaps more

\textit{pedantically} as


\[
\sum _{n=1}^{\infty} a_n \text{, where } a_n=3n-2,
\] a form that more

explicitly reflects the idea that a series is the sum of a sequence of

terms. (The expression $a_n=3n-2$ by itself defines the sequence of

terms, but note that it is a bit ambiguous since it's not clear without

additional context whether the sequence starts at $n=1$ or $n=0$, or

indeed any other value. Incidentally, the first summation formula above,

which might have seemed a bit strange at first, came from the formula

for the \textit{n}th term of an arithmetic sequence, which you probably first

learned in algebra class: $a_n=a_1+d(n-1)$.)


To verify that this summation formula represents the same sum as in the

previous section, we simply evaluate the algebraic expression ($a_n$)

for the first value of $n$ (in this case, 1), then the next value (2),

and so forth:


\[
\sum _{n=1}^{\infty} (3n-2) = [3(1)-2]+[3(2)-2]+[3(3)-2]+\ldots = 1+4+7+\ldots
\]


As for the geometric series, it should be easy to verify the

representation


\[
\sum _{n=1}^{\infty} \frac{1}{2^{n-1}}
\] or, equivalently


\[
\sum _{n=1}^{\infty} \left(\frac{1}{2}\right)^{n-1}.
\]


%-------------------------------------------------------------------------------
\subsubsection{Non-uniqueness of summation notation}
\label{sss:nonu-9309}
%-------------------------------------------------------------------------------


Note that the representation of a series in summation notation is not

unique, even ignoring simple algebraic rewrites. All of the sums seen so

far have started at $n=1$, as is usually the convention adopted for

sequences, but we could have chosen to start the sums at any other value

of $n$. In fact, it is quite common to start infinite series at $n=0$.

If we follow that convention, the arithmetic series would be written


\[
\sum _{n=0}^{\infty} (3n+1)
\] and the geometric series would be


\[
\sum _{n=0}^{\infty} \left(\frac{1}{2}\right)^n.
\] Again, these

expressions should be easy to verify.


Furthermore, as with functions, the variable used in the sum (called the

\textit{index variable}) is completely arbitrary. Thus, the following sums also

represent the same series (respectively) as those above:


\[
\sum _{k=0}^{\infty} (3k+1)
\]


\[
\sum _{i=0}^{\infty} \left(\frac{1}{2}\right)^i
\]


As a result of this fact, any series can be rewritten by simply changing

the starting value of the series and compensating for this change in the

formula for the \textit{n}th term. To be specific, consider the first series

above in this section. Using the substitution $n=m-2$, we get the series


\[
\sum _{m-2=0}^{\infty} (3[m-2]+1) = \sum _{m=2}^\infty (3m-5).
\] But

then the variable $m$ can simply be changed back to $n$ to get


\[
\sum _{n=2}^{\infty} (3n-5).
\] Notice how this new series can be seen

as resulting from ``bumping up'' the starting value of the original

series by 2 and ``bumping down'' each $n$ in the $a_n$ formula by the

same amount (that is, replacing each $n$ by $n-2$ and then simplifying).

This kind of change to a series is sometimes necessary to verify a given

series formula (say, when looking up the answer to an exercise in a

textbook).


%-------------------------------------------------------------------------------
\subsubsection{Series versus sequence}
\label{sss:seri-1648}
%-------------------------------------------------------------------------------


As mentioned \href{above}, there is an important

difference between an infinite series, which is the sum of an infinite

number of terms, and the \textit{sequence} formed by those same terms.


At the risk of being repetitive, our arithmetic series is


\[
\sum _{n=1}^{\infty} (3n-2) = 1+4+7+10+\ldots,
\] but the corresponding

sequence of terms is


\[
\\{a_n\\} _{n=1}^{\infty} = \\{3n-2\\} _{n=1}^{\infty} = 1,4,7,10,\ldots.
\]

It might seem silly at this point to make such a big deal out of this

distinction, but it will be very important going forward not to confuse

these two concepts.


%-------------------------------------------------------------------------------
\subsubsection{Sequence of partial sums}
\label{sss:sequ-97ba}
%-------------------------------------------------------------------------------


One special kind of sequence that is very useful to consider when

studying series is the \textit{sequence of partial sums}, defined in the

following way (assuming a starting value of 1 for the index variable)

for any positive integer $n$:


> \[
s_n=\sum _{k=1}^{n}a_k=a_1+a_2+a_3+\ldots+a_n
\]


For our arithmetic series, the sequence of partial sums is:


\[
s_1 =a_1=1
\]


\[
s_2 =a_1+a_2=1+4=5
\]


\[
s_3 =a_1+a_2+a_3=1+4+7=12
\]


\[
s_4 =a_1+a_2+a_3+a_4=1+4+7+10=22
\]


\[
\ldots
\] Thus


\[
\left\\{s_n\right\\}=1,5,12,22,\ldots.
\] Note how this sequence of

partial sums is very different from the sequence of terms discussed in

the last section!


Just as the series itself can be written more compactly in summation

notation by finding the pattern in its terms, the sequence of partial sums for a series can (sometimes) also be written compactly as a

function of $n$.


It should be easy to check that the following expression accurately represents the (first few terms of the) sequence of partial sums shown above:


\[
s_n =\frac{3n^2-n}{2}
\] To see where this expression came from, we

first need to review some properties of sums that you probably remember

from arithmetic and algebra, but might not be familiar with in summation

notation.


%-------------------------------------------------------------------------------
\subsubsection{Some Properties of Finite Sums}
\label{sss:some-e12c}
%-------------------------------------------------------------------------------

\textbf{Basic Properties}

\begin{itemize}

\item \textbf{Constant multiple rule:}
\[
\sum_{k=1}^{n} c\,a_k = c \sum_{k=1}^{n} a_k
\]

\item \textbf{Sum and difference rules:}
\[
\sum_{k=1}^{n} a_k + \sum_{k=1}^{n} b_k
= \sum_{k=1}^{n} (a_k + b_k)
\]
\[
\sum_{k=1}^{n} a_k - \sum_{k=1}^{n} b_k
= \sum_{k=1}^{n} (a_k - b_k)
\]

\item \textbf{Linearity:}
\[
\sum_{k=1}^{n} (c\,a_k + d\,b_k)
= c \sum_{k=1}^{n} a_k + d \sum_{k=1}^{n} b_k
\]

\end{itemize}

\textbf{Special Sum Formulas}

\[
\sum_{k=1}^{n} 1 = n
\]
\[
\sum_{k=1}^{n} k = \frac{n(n+1)}{2}
\]
\[
\sum_{k=1}^{n} r^k = \frac{r(1 - r^n)}{1 - r}, \quad r \ne 1
\]

The first formula is immediate, since it represents the sum of $n$ copies of 1. 
The second and third formulas are standard results from algebra and precalculus.

\paragraph{Example: Arithmetic Sum}

\begin{align}
\sum_{k=1}^{n} (3k - 2)
&= \sum_{k=1}^{n} 3k - \sum_{k=1}^{n} 2 \\
&= 3 \sum_{k=1}^{n} k - 2 \sum_{k=1}^{n} 1 \\
&= 3\left(\frac{n(n+1)}{2}\right) - 2n \\
&= \frac{3n^2 + 3n}{2} - \frac{4n}{2} \\
s_n &= \frac{3n^2 - n}{2}
\end{align}

\paragraph{Example: Geometric Sum}

\begin{align}
\sum_{k=1}^{n} \left(\tfrac{1}{2}\right)^{k-1}
&= \sum_{k=1}^{n} \left(\tfrac{1}{2}\right)^k \cdot 2 \\
&= 2 \sum_{k=1}^{n} \left(\tfrac{1}{2}\right)^k \\
&= 2 \left( \frac{\frac{1}{2}\left(1 - \left(\tfrac{1}{2}\right)^n\right)}{1 - \tfrac{1}{2}} \right) \\
s_n &= 2\left[1 - \left(\tfrac{1}{2}\right)^n\right]
\end{align}


%-------------------------------------------------------------------------------
\subsubsection{Finding a series from its partial sums}
\label{sss:find-5833}
%-------------------------------------------------------------------------------


Using the fact that


\[
s_n=a_1+a_2+\ldots+a _{n-1}+a_n=s _{n-1}+a_n,
\] we see that


> \[
a_n=s_n-s _{n-1}.
\]


This provides a way of ``recovering'' the original series from its sequence of partial sums.


For our arithmetic series:


\[
a_n = s_n-s _{n-1}
\]

\[
=\left[\frac{3n^2-n}{2}\right]-\left[\frac{3(n-1)^2-(n-1)}{2}\right]
\]

\[
=\left[\frac{3n^2-n}{2}\right]-\left[\frac{3(n^2-2n+1)-(n-1)}{2}\right]
\]

\[
=\left[\frac{3n^2-n}{2}\right]-\left[\frac{3n^2-6n+3-n+1}{2}\right]
\]

\[
=\frac{3n^2-n-3n^2+6n-3+n-1}{2}
\]

\[
=\frac{6n-4}{2}
\]

\[
a_n=3n-2
\]


%-------------------------------------------------------------------------------
\subsubsection{Sum of an infinite series}
\label{sss:sumo-5512}
%-------------------------------------------------------------------------------


Now we can finally formally define the \textit{sum of an infinite series} as the limit of its sequence of partial sums:


> \[
\sum _{n=1}^{\infty} a_n=\lim _{n\to\infty}s_n
\]


If the limit converges to a real number, say $s$, then the infinite series is said to \textit{converge to the sum} $s$, or to \textit{be convergent with sum} $s$. If the limit diverges (including the cases where the limit is infinity or negative infinity), then the series is said to \textit{diverge} or to \textit{be divergent} in the same way, and its sum is said to \textit{not exist} (or to be infinity or negative infinity, as appropriate). Note that one does not describe the series itself as ``not existing'' when it diverges, only its sum.


Consider again the arithmetic and geometric series we have been discussing up to this point. It is obvious by simply looking at the original arithmetic series


\[
1+4+7+10+13+\dots
\] that it does not have a finite sum. The terms being added are themselves getting larger and larger without bound, so the sum is getting larger and larger without bound.


The sequence of partial sums given above formalizes this idea. In particular, because


\[
\sum _{n=1}^{\infty} (3n-2) = \lim _{n\to\infty} \frac{3n^2-n}{2} = \infty,
\]

the arithmetic series diverges to infinity.


Now consider the original geometric series:


\[
1+\frac{1}{2}+\frac{1}{4}+\frac{1}{8}+\frac{1}{16}+\ldots
\] The terms here are getting smaller and smaller, and indeed are approaching zero. Since adding zero to something doesn't change its value, it seems reasonable to suspect that the sum might be a fixed, finite number. What does the sequence of partial sums reveal?


\[
\sum _{n=1}^{\infty} (\tfrac{1}{2})^{n-1} = \lim _{n\to\infty} 2[1-(\tfrac{1}{2})^n] = 2(1-0)=2
\]

So the geometric series does, in fact, converge to the finite sum 2.


It seems obvious that any series whose terms ``blow up'' to infinity (like our arithmetic series) will diverge, but does every series whose terms shrink to zero (like our geometric series) converge to a finite sum? It turns out the answer to that question is no.


%-------------------------------------------------------------------------------
\subsubsection{Harmonic Series}
\label{sss:harm-8fb1}
%-------------------------------------------------------------------------------

There is a very important series whose terms approach zero, yet the series 
diverges because its sequence of \textit{partial sums} diverges to infinity. 
It is called the \textit{harmonic series}:
\[
\sum_{n=1}^{\infty} \frac{1}{n}.
\]

Clearly, the terms approach zero:
\[
\lim_{n \to \infty} \frac{1}{n} = 0.
\]

To analyze the partial sums, consider the $2^n$th partial sum and group 
terms as follows:
\[
\begin{aligned}
\sum_{k=1}^{2^n} \frac{1}{k}
&= 1 
+ \underbrace{\frac{1}{2}}_{\text{group 1}} 
+ \underbrace{\left(\frac{1}{3} + \frac{1}{4}\right)}_{\text{group 2}} 
+ \underbrace{\left(\frac{1}{5} + \cdots + \frac{1}{8}\right)}_{\text{group 3}} 
+ \cdots 
+ \underbrace{\sum_{k=2^{n-1}+1}^{2^n} \frac{1}{k}}_{\text{group }n}.
\end{aligned}
\]

Each group can be bounded below:
\[
\begin{aligned}
\sum_{k=1}^{2^n} \frac{1}{k}
&> 1 
+ \frac{1}{2}
+ \frac{2}{4}
+ \frac{4}{8}
+ \cdots
+ \frac{2^{n-1}}{2^n} \\
&= 1 + \frac{1}{2} + \frac{1}{2} + \cdots + \frac{1}{2}.
\end{aligned}
\]

There are $n$ terms equal to $\tfrac{1}{2}$ in this expression, so
\[
\sum_{k=1}^{2^n} \frac{1}{k} > 1 + \frac{n}{2}.
\]

As $n \to \infty$, the right-hand side diverges, so the sequence of 
partial sums diverges. Therefore, the harmonic series diverges.

\medskip

Even though its terms approach zero, the harmonic series diverges. 
Many other series share this property.

%-------------------------------------------------------------------------------
\subsubsection{``Eventually''}
\label{sss:even-a4ef}
%-------------------------------------------------------------------------------

When determining whether a series converges or diverges, it is often useful 
to ``ignore'' a finite number of initial terms.

For example, the geometric series
\[
1 + \frac{1}{2} + \frac{1}{4} + \frac{1}{8} + \cdots
\]
converges to $2$, whereas the series
\[
-1 + 3 + 1 + \frac{1}{2} + \frac{1}{4} + \frac{1}{8} + \cdots
\]
is not geometric at the beginning, but becomes geometric from the third 
term onward. We say that this series is \textit{eventually geometric}.

Since convergence depends only on the behavior of the tail of the series, 
this series also converges. Indeed,
\[
-1 + 3 + \sum_{n=1}^{\infty} \frac{1}{2^{n-1}} = (-1 + 3) + 2 = 4.
\]

Similarly, the series
\[
\frac{1}{4} + \frac{1}{8} + \frac{1}{16} + \cdots
\]
is simply the original geometric series with the first two terms removed. Thus,
\[
\sum_{n=3}^{\infty} \frac{1}{2^{n-1}}
= \left(\sum_{n=1}^{\infty} \frac{1}{2^{n-1}}\right)
- \left(\sum_{n=1}^{2} \frac{1}{2^{n-1}}\right)
= 2 - \left(1 + \tfrac{1}{2}\right)
= \tfrac{1}{2}.
\]

\medskip

\textbf{Key Insight:}  
Whether a series converges does not depend on its initial terms, but the 
\textit{value of the sum} does.


%-------------------------------------------------------------------------------
\subsubsection{Series convergence and divergence}
\label{sss:seri-f0f7}
%-------------------------------------------------------------------------------


Now that we've covered enough background about series, we can start to
systematically investigate the convergence and divergence of many
different kinds of infinite series. In some cases we will be able to
find a formula for the partial sums, and thus the sum of the series, but
in most cases we will not.


%-------------------------------------------------------------------------------
\subsubsection{Necessary condition for convergence}
\label{sss:nece-6314}
%-------------------------------------------------------------------------------


Note that we have already considered one convergent series whose terms shrink to zero (the geometric series), one divergent series whose terms shrink to zero (the harmonic series), and one divergent series whose terms do not shrink to zero (the arithmetic series). What about the remaining case: a convergent series whose terms do not shrink to zero?


Turns out, that's not possible.


> \textbf{Theorem: A necessary condition for series convergence}

> If the infinite series $\sum _{n=c}^\infty a_n$ converges (where $c$ is any integer), then $\lim _{n\to\infty} a_n = 0$.


So, every convergent series has a sequence of terms that converges to zero. But, as we saw with the harmonic series, just because a series has terms that shrink to zero, that doesn't mean that the series converges. Therefore, having terms that go to zero is a \textit{necessary} condition that a convergent series must satisfy, but not a \textit{sufficient} condition that allows us to \textit{conclude} that a given series converges. (Put more bluntly: \textit{Just because the terms go to zero doesn't mean the series converges!})


So what \textit{is} a sufficient condition for series convergence? Well, there are many such conditions, several of which will be discussed later. All of them, however, will be saying essentially the same thing: that the terms of the series are going to zero ``fast enough'' that it allows the series to converge.


%-------------------------------------------------------------------------------
\subsubsection{Sufficient condition for divergence}
\label{sss:suff-4468}
%-------------------------------------------------------------------------------


Recall from basic logic that the statement ``If A then B'' is equivalent to ``If not B then not A''. Applying this idea to the previous theorem gives a way of determining that an infinite series diverges (the so-called \textit{divergence test}, which we will see again).


> \textbf{Theorem: A sufficient condition for series divergence}

> If it is not true that $\lim _{n\to\infty} a_n = 0$ (including the case in which the limit does not exist), then the infinite series $\sum _{n=c}^\infty a_n$ (where $c$ is any integer) diverges.

> \textbf{Proof}

> This is an immediate consequence of the last theorem. The argument is a simple proof by contradiction: Let $\lim _{n\to\infty} a_n$ either not equal 0 or not exist. Assume that $\sum _{n=c}^\infty a_n$ converges (for some integer $c$). Then by the theorem above, $\lim _{n\to\infty} a_n$ must equal 0. This is clearly a contradiction. Thus $\sum _{n=c}^\infty a_n$ must diverge.


While this gives a very useful way of checking for divergence, it is not a very ``powerful'' method, since: as we've already said: there are many divergent infinite series whose terms do actually go to zero.


%-------------------------------------------------------------------------------
\subsubsection{Properties of infinite series}
\label{sss:prop-58bd}
%-------------------------------------------------------------------------------


There are several properties of infinite series that are direct

consequences of the corresponding properties of finite sums:


Constants can be factored out of (or multiplied into) infinite series:


\begin{itemize}
    \item For any real number $c$, $\textstyle \sum _{k=1}^{\infty} c\,a_k$
\end{itemize}
    converges if $\textstyle \sum _{k=1}^{\infty} a_k$ converges, in which case $\newline \textstyle \sum _{k=1}^{\infty} c\,a_k = c\sum _{k=1}^{\infty}a_k.\newline$

\begin{itemize}
    \item For any non-zero real number $c$,
\end{itemize}
    $\textstyle \sum _{k=1}^{\infty} c\,a_k$ diverges if

    $\textstyle \sum _{k=1}^{\infty} a_k$ diverges.


Note that $\textstyle \sum _{k=1}^{\infty} c\,a_k$ \textit{converges} if

$\textstyle \sum _{k=1}^{\infty} a_k$ diverges but $c=0$ (because it is a

series whose terms are all zero).


Sums and differences of convergent infinite series are convergent:


\begin{itemize}
    \item If $\textstyle \sum _{k=1}^{\infty} a_k$ and
\end{itemize}
    $\textstyle \sum _{k=1}^{\infty} b_k$ both converge, then

    $\textstyle \sum _{k=1}^{\infty} (a_k+b_k)$ also converges and $\newline \textstyle \sum _{k=1}^{\infty} (a_k+b_k) = \sum _{k=1}^{\infty} a_k + \sum _{k=1}^{\infty} b_k.\newline$

\begin{itemize}
    \item If $\textstyle \sum _{k=1}^{\infty} a_k$ and
\end{itemize}
    $\textstyle \sum _{k=1}^{\infty} b_k$ both converge, then

    $\textstyle \sum _{k=1}^{\infty} (a_k-b_k)$ also converges and

    $\newline \textstyle \sum _{k=1}^{\infty} (a_k-b_k) = \sum _{k=1}^{\infty} a_k - \sum _{k=1}^{\infty} b_k. \newline$


Sums and differences of one convergent and one divergent series are

divergent:


\begin{itemize}
    \item If $\textstyle \sum _{k=1}^{\infty} a_k$ converges and
\end{itemize}
    $\textstyle \sum _{k=1}^{\infty} b_k$ diverges, then

    $\textstyle \sum _{k=1}^{\infty} (a_k+b_k)$ and

    $\textstyle \sum _{k=1}^{\infty} (a_k-b_k)$ both diverge.

\begin{itemize}
    \item If $\textstyle \sum _{k=1}^{\infty} a_k$ diverges and
\end{itemize}
    $\textstyle \sum _{k=1}^{\infty} b_k$ converges, then

    $\textstyle \sum _{k=1}^{\infty} (a_k+b_k)$ and

    $\textstyle \sum _{k=1}^{\infty} (a_k-b_k)$ both diverge.


Sums and differences of two divergent series may be convergent or

divergent (thus one cannot decide whether they converge or diverge

without applying a specific convergence or divergence test):


\begin{itemize}
    \item If $\textstyle \sum _{k=1}^{\infty} a_k$ and
\end{itemize}
    $\textstyle \sum _{k=1}^{\infty} b_k$ both diverge, then

    $\textstyle \sum _{k=1}^{\infty} (a_k+b_k)$ and

    $\textstyle \sum _{k=1}^{\infty} (a_k-b_k)$ may either (both)

    converge or (both) diverge.


More generally, linear combinations of convergent infinite series are

convergent:


\begin{itemize}
    \item If $\textstyle \sum _{k=1}^{\infty} a_k$ and
\end{itemize}
    $\textstyle \sum _{k=1}^{\infty} b_k$ both converge and $c$ and $d$

    are any real numbers, then

    $\textstyle \sum _{k=1}^{\infty} (c\,a_k+d\,b_k)$ also converges and

    $\textstyle \sum _{k=1}^{\infty} (c\,a_k+d\,b_k) = c \sum _{k=1}^{\infty} a_k + d \sum _{k=1}^{\infty} b_k.$


Linear combinations of a mix of convergent and divergent series are

(usually) divergent:


\begin{itemize}
    \item If $\textstyle \sum _{k=1}^{\infty} a_k$ converges and
\end{itemize}
    $\textstyle \sum _{k=1}^{\infty} b_k$ diverges, and $c$ and $d$ are

    any real numbers with $d \ne 0$, then

    $\textstyle \sum _{k=1}^{\infty} (c\,a_k+d\,b_k)$ diverges. (If

    $d=0$, then $\textstyle \sum _{k=1}^{\infty} (c\,a_k+d\,b_k)$

    converges.)


However, arbitrary linear combinations of divergent series are (usually)

inconclusive with respect to their convergence without further testing:


\begin{itemize}
    \item If $\textstyle \sum _{k=1}^{\infty} a_k$ and $\textstyle \sum _{k=1}^{\infty} b_k$ both diverge and $c$ and $d$ are non-zero real numbers, then $\textstyle \sum _{k=1}^{\infty} (c\,a_k+d\,b_k)$ may either converge or diverge. (If exactly one of the constants $c$ and $d$ is zero, then the latter series diverges. If both $c=0$ and $d=0$, then the latter series converges.)
\end{itemize}


%-------------------------------------------------------------------------------
\subsubsection{Resources}
\label{sss:reso-55b5}
%-------------------------------------------------------------------------------


\href{https://youtu.be/jPx2S4mSZWc}{Introduction to Infinite Series} by James Sousa


\href{https://youtu.be/LPFSQoozBL0}{Calculating the first terms in a series of partial sums} by Krista King


\href{https://youtu.be/wpB7nk5J0AA}{Sum of the series of partial sums} by Krista King


\href{https://youtu.be/-wvF8OQSMx8}{Convergence \& Divergence - Geometric Series, Telescoping Series, Harmonic Series, Divergence Test} by The Organic Chemistry Tutor


\href{https://youtu.be/6LfoMmCckFY}{Telescoping Series} by James Sousa


\href{https://youtu.be/cyoiIBs7kIg}{Finding a Formula for a Partial Sum of a Telescoping Series} by patrickJMT


\href{https://youtu.be/xmyLUFv2mDM}{Telescoping Series ,Showing Divergence Using Partial Sums} by patrickJMT


\href{https://youtu.be/H8HVYiEqzBo}{Telescoping Series , Finding the Sum, Example 1} by patrickJMT


\href{https://youtu.be/pFcJZnxqzNc}{Infinite Geometric Series} by James Sousa


\href{https://youtu.be/xjmy5hkZccY}{Geometric Series and the Test for Divergence - Part 1} by patrickJMT


\href{https://youtu.be/C8piSCOdo1Y}{Geometric Series and the Test for Divergence - Part 2} by patrickJMT


\href{https://youtu.be/LhO7i3RFri0}{Convergence of a geometric series} by Krista King


\href{https://youtu.be/1eInMsp20cI}{Values for which the geometric series converges} by Kritsa King


\href{https://youtu.be/-5kIBPR2Npk}{Geometric Series and Geometric Sequences - Basic Introduction} by The Organic Chemistry Tutor


\href{https://youtu.be/jxRqRLMliPc}{Finding The Sum of an Infinite Geometric Series} by The Organic Chemistry Tutor


%-------------------------------------------------------------------------------
\subsubsection{Licensing}
\label{sss:lice-7352}
%-------------------------------------------------------------------------------


Content obtained and/or adapted from:


\begin{itemize}
    \item \href{https://en.wikibooks.org/wiki/User:Dcljr/Series}{Series, Wikibooks}
\end{itemize}
    under a CC BY-SA license

\subsection{Learning Objectives}

\begin{enumerate}
    \item \textbf{Identify Series Types:}
    \item \textbf{Outcome:} Students will be able to identify and classify series as arithmetic or geometric based on the pattern of their terms.
    \item \textbf{Details:} This involves recognizing whether a series has a constant difference (arithmetic) or a constant ratio (geometric) between consecutive terms. Students will also be able to provide examples of each type.
    \item \textbf{Apply Summation Notation:}
    \item \textbf{Outcome:} Students will be able to express series using summation notation.
    \item \textbf{Details:} This includes converting a given series into its summation notation form, understanding the notation components (such as the Sigma symbol and index variable), and interpreting series written in different notations. Students should be able to rewrite series starting at different index values and understand the implications of these changes.
    \item \textbf{Analyze Convergence and Divergence:}
    \item \textbf{Outcome:} Students will be able to determine whether a given series converges or diverges and compute its sum if it converges.
    \item \textbf{Details:} This involves evaluating series by examining the sequence of partial sums or using known convergence tests. Students should be able to apply these concepts to arithmetic and geometric series, as well as recognize cases where series diverge, such as the harmonic series.
\end{enumerate}

\subsection{Assessment Instrument}
\label{ssec:seri-2a98}

Students classify series as geometric, telescoping, or general, apply the correct convergence test, and compute sums when possible. The divergence test (necessary condition) is applied before any deeper test.

\textbf{Example problems.}

\begin{enumerate}
    \item Determine whether $\sum_{n=1}^{\infty}\dfrac{3}{4^n}$ converges. If so, find its exact sum.
    \item Write out partial sums and find $\sum_{n=1}^{\infty}\dfrac{1}{n(n+1)}$ as a telescoping series.
    \item Show that $\sum_{n=1}^{\infty}\dfrac{n}{n+1}$ diverges using the divergence test.
\end{enumerate}

%===============================================================================
\section{The divergence and integral tests}
\label{sec:thed-219c}
%===============================================================================


%-------------------------------------------------------------------------------
\subsubsection{Theorem (p-series)}
\label{sss:theo-5ce0}
%-------------------------------------------------------------------------------

The \textbf{$p$-series}
\[
\sum_{n=1}^{\infty} \frac{1}{n^p}
\]
converges if $p > 1$ and diverges if $0 < p \leq 1$.

\textbf{Sketch of Proof:}

If $p = 1$, the series becomes the harmonic series
\[
\sum_{n=1}^{\infty} \frac{1}{n},
\]
which diverges.

If $p > 1$, one can show (e.g., via the integral test) that
\[
\sum_{n=1}^{\infty} \frac{1}{n^p}
\]
converges.

If $0 < p < 1$, then for all $n \geq 1$,
\[
\frac{1}{n^p} \geq \frac{1}{n}.
\]
Since the harmonic series diverges, the comparison test implies that
\[
\sum_{n=1}^{\infty} \frac{1}{n^p}
\]
also diverges.

%-------------------------------------------------------------------------------
\subsubsection{Limit Test}
\label{sss:limi-19f1}
%-------------------------------------------------------------------------------

The limit test determines whether a series can possibly converge.

\paragraph{Limit Test for Convergence}

If a series
\[
\sum_{n=1}^{\infty} s_n
\]
satisfies
\[
\lim_{n \to \infty} s_n \neq 0,
\]
then the series diverges.

If
\[
\lim_{n \to \infty} s_n = 0,
\]
the test is inconclusive.

\textbf{Explanation:}

If the terms $s_n$ do not approach zero, then for large $n$ they remain 
close to some nonzero value $L$, so the partial sums behave like those 
of an arithmetic sequence and cannot converge.

\textbf{Important Note:}

This is a test for \textbf{divergence only}. If $\lim_{n \to \infty} s_n = 0$, 
the series may either converge or diverge, and further tests are required.

%-------------------------------------------------------------------------------
\subsubsection{Example 1}
\label{sss:exam-acc3}
%-------------------------------------------------------------------------------

Determine whether the series
\[
\sum_{n=1}^{\infty} \frac{n+1}{n}
\]
diverges or whether the limit test is inconclusive.

%-------------------------------------------------------------------------------
\subsubsection{Solution}
\label{sss:solu-acc3}
%-------------------------------------------------------------------------------

We compute
\[
\lim_{n \to \infty} \frac{n+1}{n} = 1.
\]

Since the limit is not zero, the series diverges by the limit test.

%-------------------------------------------------------------------------------
\subsubsection{Example 2}
\label{sss:exam-5bdd}
%-------------------------------------------------------------------------------

Determine whether the series
\[
\sum_{n=1}^{\infty} \frac{n^2 + 5n + 6}{3n^2 + 1}
\]
diverges or whether the limit test is inconclusive.

%-------------------------------------------------------------------------------
\subsubsection{Solution}
\label{sss:solu-5bdd}
%-------------------------------------------------------------------------------

We compute
\[
\lim_{n \to \infty} \frac{n^2 + 5n + 6}{3n^2 + 1} = \frac{1}{3}.
\]

Since the limit is not zero, the series diverges by the limit test.

%-------------------------------------------------------------------------------
\subsubsection{Example 3}
\label{sss:exam-ad87}
%-------------------------------------------------------------------------------

Determine whether the series
\[
\sum_{n=1}^{\infty} \frac{n^2 + 5n + 6}{n^3}
\]
diverges or whether the limit test is inconclusive.

%-------------------------------------------------------------------------------
\subsubsection{Solution}
\label{sss:solu-ad87}
%-------------------------------------------------------------------------------

We compute
\[
\lim_{n \to \infty} \frac{n^2 + 5n + 6}{n^3}
= \lim_{n \to \infty} \left(\frac{1}{n} + \frac{5}{n^2} + \frac{6}{n^3}\right)
= 0.
\]

Since the limit is zero, the limit test is inconclusive. Further analysis 
is required to determine whether the series converges or diverges.


%-------------------------------------------------------------------------------
\subsubsection{Solution}
\label{sss:solu-5502}
%-------------------------------------------------------------------------------


Because


$\lim _{n \rightarrow \infty}{\frac{n^2+5n+6}{n^3}}= 0$


the limit \textit{is} zero and so the test is inconclusive. Further analysis is needed to determine whether or not the series converges or diverges.


%-------------------------------------------------------------------------------
\subsubsection{Integral Test}
\label{sss:inte-e4c4}
%-------------------------------------------------------------------------------


The next test for convergence for infinite series is the \textbf{integral test}. The integral test utilizes the fact that an integral is essentially an Riemann Sum: which is itself an infinite sum: over an infinite interval which is useful because integration is relatively straight forward and familiar.


The test is as follows:


> \textbf{Integral Test for Convergence and Divergence}

> If a series $S= \sum _{n=j}^{\infty}{s_n}$, and if the discrete function $s(n)$ is continuous and decreasing on the interval $[j, \infty)$ then if

>

> 1.  $\int _{j}^{\infty}{s(n)dn}$ is convergent, so is $S$

> 2.  $\int _{j}^{\infty}{s(n)dn}$ is divergent, so is $S$


Let us look at why this this test works. As an example, we'll use the \textbf{harmonic series} $\sum _{n=1}^{\infty}{\frac{1}{n}}$. The harmonic series is well known series that actually happens to be divergent. This series actually passes the limit test, however it does not pass the integral test. If we want to approximate the integral, we can use rectangles like in a Riemann Sum:


% [Image: Right-handed Riemann Sum]


Figure: Right-handed Riemann Sum


Note that this right-handed method will always under approximate the integral (assuming that the function is decreasing on our selected interval). This implies that if the right-handed sum is equal to the actual infinite series, then the integral itself must be greater than the sum. This can help show convergence, because if the integral from the starting point to infinity is convergent, then by the comparison test the original function on that interval must also be convergent. And so we can see that the integral test is actually a ``special case'' of the comparison test.


But what about divergence? This case is also satisfied: if we use a left-handed approximation instead of a right handed one, we see that we again attain the original series, however there is an important difference:


% [Image: Left-handed Riemann Sum]


Figure: Left-handed Riemann Sum


The key difference, in this case, is that the integral becomes an under approximation for the series, and we can use the new ``series'' of the integral to show divergence with the comparison test.


This test is useful but is unfortunately only useful on functions that can be integrated \textit{and} are decreasing in size. The latter may seem like a trivial and unnecessary addition, however consider how this test works; it relies on the fact that the integral of a function decreasing on an interval will always yield an under/over approximation of a series; if the function is not decreasing everywhere on the interval, the integral will not necessarily yield an under/over approximation every time.


%-------------------------------------------------------------------------------
\subsubsection{Example 1}
\label{sss:exam-acc3}
%-------------------------------------------------------------------------------


Use the integral test to determine if the following series is convergent or divergent:


$\sum _{n=1}^{\infty}{\frac{1}{n^2}}$


%-------------------------------------------------------------------------------
\subsubsection{Solution}
\label{sss:solu-5502}
%-------------------------------------------------------------------------------


The improper integral from $1$ returns

$\lim _{x \rightarrow \infty} \frac{-1}{n} - \frac{-1}{(1)}$. The limit of the first term is zero, and so because the integral converges the series does too.


%-------------------------------------------------------------------------------
\subsubsection{Example 2}
\label{sss:exam-5bdd}
%-------------------------------------------------------------------------------


Determine if the series $\sum _{x=1}^{\infty}{\frac{x^2+1}{x}}$ converges or diverges.


%-------------------------------------------------------------------------------
\subsubsection{Solution}
\label{sss:solu-5502}
%-------------------------------------------------------------------------------


This series does not pass the first requirement that the series is decreasing over the desired interval; $\lim _{x \rightarrow \infty}{\frac{x^2+1}{x}}=\infty$. Applying the integral test would still show the series' convergence.


%-------------------------------------------------------------------------------
\subsubsection{Example 3}
\label{sss:exam-ad87}
%-------------------------------------------------------------------------------


Determine if the series $\sum _{n=1}^{\infty}{\frac{1}{(n-3)^2+1}}$ converges or diverges.


%-------------------------------------------------------------------------------
\subsubsection{Solution}
\label{sss:solu-5502}
%-------------------------------------------------------------------------------


However, this series is not decreasing everywhere on the interval. However, it has a relative maximum at $n=1$, after which it does decrease forever. And so we can write this series as $\sum _{n=1}^{2}{\frac{1}{(n-3)^2+1}} + \sum _{n=3}^{\infty}{\frac{1}{(n-3)^2+1}}$. Integrating the function yields the improper integral $\tan^{-1}(n-3)$ from $3$ to infinity, which converges, and so the series converges too.


%-------------------------------------------------------------------------------
\subsubsection{Resources}
\label{sss:reso-55b5}
%-------------------------------------------------------------------------------


%-------------------------------------------------------------------------------
\subsubsection{The Divergence Test}
\label{sss:thed-5a78}
%-------------------------------------------------------------------------------


\begin{itemize}
    \item \href{https://www.youtube.com/watch?v=4vGr104Qpj4}{The Nth Term Divergence
\end{itemize}
    Test} Video by James Sousa, Math is Power 4U

\begin{itemize}
    \item \href{https://www.youtube.com/watch?v=isIs1-Rdpv0}{The Nth Term Divergence Test
\end{itemize}
    (Rational)} Video by James Sousa, Math is Power 4U

\begin{itemize}
    \item \href{https://www.youtube.com/watch?v=5ciNNv2eL6k}{The Nth Term Divergence Test
\end{itemize}
    (Geometric)} Video by James Sousa, Math is Power 4U

\begin{itemize}
    \item \href{https://youtu.be/DcduHCJuiWk}{Test for Divergence for Series, Two
\end{itemize}
    Examples} Video by Patrick JMT

\begin{itemize}
    \item \href{https://youtu.be/e50wqj2Y-18}{The nth Term Test, Divergence Test, Zero
\end{itemize}
    Test} Video by Krista King

\begin{itemize}
    \item \href{https://youtu.be/nC_IU1IzzS4}{Divergence Test For Series} Video by The Organic Chemistry Tutor
\end{itemize}


%-------------------------------------------------------------------------------
\subsubsection{The Integral Test}
\label{sss:thei-442b}
%-------------------------------------------------------------------------------


\begin{itemize}
    \item \href{https://www.youtube.com/watch?v=qQlnPrpwKY8}{The Integral Test}
\end{itemize}
    Video by James Sousa, Math is Power 4U

\begin{itemize}
    \item \href{https://www.youtube.com/watch?v=Lla2xi6W8cQ}{Infinite Series - The Integral
\end{itemize}
    Test} Video by James Sousa, Math is Power 4U

\begin{itemize}
    \item \href{https://www.youtube.com/watch?v=tXOMw8sVnKY}{Integral Test Example Rational
\end{itemize}
    Function} Video by James Sousa, Math is Power 4U

\begin{itemize}
    \item \href{https://www.youtube.com/watch?v=GeM-KGovJi0}{Integral Test Example Rational
\end{itemize}
    Function} Video by James Sousa, Math is Power 4U

\begin{itemize}
    \item \href{https://www.youtube.com/watch?v=YqGuL3kFp04}{Integral Test Example Radical
\end{itemize}
    Function} Video by James Sousa, Math is Power 4U

\begin{itemize}
    \item \href{https://www.youtube.com/watch?v=Jxu-ggQU-Bc}{Integral Test Example Exponential
\end{itemize}
    Function} Video by James Sousa, Math is Power 4U

\begin{itemize}
    \item \href{https://www.youtube.com/watch?v=xTA0eb6USl0}{Integral Test Example Involving
\end{itemize}
    Arctangent} Video by James Sousa, Math is Power 4U

\begin{itemize}
    \item \href{https://www.youtube.com/watch?v=1riNZY-8ATM}{Integral Test Example Involving Integration by
\end{itemize}
    Parts} Video by James Sousa, Math is Power 4U





\begin{itemize}
    \item \href{https://www.youtube.com/watch?v=gvZ7UIaC50M}{Integral Test - Basic
\end{itemize}
    Idea} Video by Patrick JMT

\begin{itemize}
    \item \href{https://www.youtube.com/watch?v=XHUxLLu_c_0}{Integral Test for Series Example
\end{itemize}
    1} Video by Patrick JMT

\begin{itemize}
    \item \href{https://www.youtube.com/watch?v=f0hGrW6_weo}{Integral Test for Series Example
\end{itemize}
    2} Video by Patrick JMT

\begin{itemize}
    \item \href{https://www.youtube.com/watch?v=TZAmpPZ4DH4}{Integral Test to Evaluate Series Ex
\end{itemize}
    1} Video by Patrick JMT

\begin{itemize}
    \item \href{https://www.youtube.com/watch?v=MvfCLjfz-cc}{Integral Test to Evaluate Series Ex
\end{itemize}
    2} Video by Patrick JMT

\begin{itemize}
    \item \href{https://www.youtube.com/watch?v=xHaV9Ambgqc}{Integral Test to Evaluate Series Ex
\end{itemize}
    3} Video by Patrick JMT

\begin{itemize}
    \item \href{https://www.youtube.com/watch?v=_KDtZ5wRq8Q}{Integral Test to Evaluate Series Ex
\end{itemize}
    4} Video by Patrick JMT





\begin{itemize}
    \item \href{https://www.youtube.com/watch?v=qiwGXzbans4}{How to Use the Integral
\end{itemize}
    Test} Video by Krista King





\begin{itemize}
    \item \href{https://www.youtube.com/watch?v=f9SJz4-UaQQ}{Integral Test for Series} Video by The Organic Chemistry Tutor
\end{itemize}


%-------------------------------------------------------------------------------
\subsubsection{The p-Series Test}
\label{sss:thep-f2b0}
%-------------------------------------------------------------------------------


\begin{itemize}
    \item \href{https://en.wikibooks.org/wiki/Real_Analysis/Series}{Series},
\end{itemize}
    Wikibooks: Real Analysis

\begin{itemize}
    \item \href{https://www.youtube.com/watch?v=gB43EA9659c}{The p-Series Test}
\end{itemize}
    Video by James Sousa, Math is Power 4U

\begin{itemize}
    \item \href{https://www.youtube.com/watch?v=L3eqCqzBH1I}{Infinite Series - The p-Series
\end{itemize}
    Test} Video by James Sousa, Math is Power 4U

\begin{itemize}
    \item \href{https://www.youtube.com/watch?v=YOy72v7IwuM}{The p-Series Test Example 1} Video by James
\end{itemize}
    Sousa, Math is Power 4U

\begin{itemize}
    \item \href{https://www.youtube.com/watch?v=Ehq3QwGVidU}{The p-Series Test Example 2} Video by James Sousa, Math is Power 4U
\end{itemize}




\begin{itemize}
    \item \href{https://www.youtube.com/watch?v=w9fC8RaklhA}{The p-Series Test}
\end{itemize}
    Video by Patrick JMT





\begin{itemize}
    \item \href{https://www.youtube.com/watch?v=koZR5uEBkdQ}{The p-Series Test}
\end{itemize}
    Video by Krista King





\begin{itemize}
    \item \href{https://www.youtube.com/watch?v=x621rtIwb_4}{The p-Series Test}
\end{itemize}
    Video by The Organic Chemistry Tutor


%-------------------------------------------------------------------------------
\subsubsection{The Remainder Estimate for the Integral Test}
\label{sss:ther-4888}
%-------------------------------------------------------------------------------


\begin{itemize}
    \item \href{https://www.youtube.com/watch?v=sVQdVA-icGA}{Remainder Estimate for the Integral
\end{itemize}
    Test} Video by Patrick JMT





\begin{itemize}
    \item \href{https://www.youtube.com/watch?v=93-zn9jY-78}{Estimating the Error/Remainder of a
\end{itemize}
    Series} Video by Krista King





\begin{itemize}
    \item \href{https://www.youtube.com/watch?v=vaU5DDcTfZk}{Remainder Estimate for the Integral
\end{itemize}
    Test} Video by The Organic Chemistry Tutor


%-------------------------------------------------------------------------------
\subsubsection{Licensing}
\label{sss:lice-7352}
%-------------------------------------------------------------------------------


Content obtained and/or adapted from:


\begin{itemize}
    \item \href{https://en.wikibooks.org/wiki/Real_Analysis/Series}{Series, Wikibooks: Real
\end{itemize}
    Analysis} under

    a CC BY-SA license

\begin{itemize}
    \item \href{https://en.wikibooks.org/wiki/Calculus/Integral_Test_for_Convergence}{Integral Test for Convergence, Wikibooks:
\end{itemize}
    Calculus}

    under a CC BY-SA license

\begin{itemize}
    \item \href{https://en.wikibooks.org/wiki/Calculus/Limit_Test_for_Convergence}{Limit Test for Convergence, Wikibooks:
\end{itemize}
    Calculus}

    under a CC BY-SA license

\subsection{Learning Objectives}

\begin{enumerate}
    \item Student Learning Objectives (SLOs) for Theorems and Tests on Series
    \item \textbf{Understand the p-Series Test:}
    \item \textbf{Outcome:} Students will be able to apply the p-series test to determine the convergence or divergence of a series.
    \item \textbf{Details:} This involves recognizing the form of a p-series, understanding the conditions under which the series converges (when $p > 1$) or diverges (when $0 < p \le 1$), and justifying these conclusions through appropriate comparison or integral tests.
    \item \textbf{Apply the Limit Test for Divergence:}
    \item \textbf{Outcome:} Students will be able to use the limit test to identify divergent series by evaluating the limit of the general term.
    \item \textbf{Details:} This includes calculating the limit of the general term of a series, determining if the limit is non-zero, and concluding the divergence of the series if the limit is not zero. Students will also recognize when the limit test is inconclusive and requires further analysis.
    \item \textbf{Utilize the Integral Test for Series Convergence:}
    \item \textbf{Outcome:} Students will be able to apply the integral test to determine the convergence or divergence of a series.
    \item \textbf{Details:} This involves checking if the function associated with the series is continuous, positive, and decreasing, and then evaluating the corresponding improper integral to draw conclusions about the series. Students will need to perform integration and compare results to assess series convergence.
\end{enumerate}

\subsection{Assessment Instrument}
\label{ssec:thed-18c5}

Problems require applying the Divergence Test as a first check, then the Integral Test where the series terms correspond to a positive, decreasing, continuous function. Students must verify the hypotheses of the Integral Test before applying it.

\textbf{Example problems.}

\begin{enumerate}
    \item Apply the Divergence Test to $\sum_{n=1}^{\infty}\dfrac{n^2}{2n^2+1}$ and state your conclusion.
    \item Use the Integral Test to determine whether $\sum_{n=1}^{\infty}\dfrac{1}{n^2}$ converges. Verify the hypotheses before integrating.
    \item Determine convergence of $\sum_{n=1}^{\infty}\dfrac{1}{n\ln n}$ using the Integral Test.
\end{enumerate}

%===============================================================================
\section{Comparison tests}
\label{sec:comparison-tests}
%===============================================================================


%-------------------------------------------------------------------------------
\subsubsection{Comparison Test}
\label{sss:comp-edc1}
%-------------------------------------------------------------------------------


The first real determiner of convergence is the \textbf{comparison test}. This test is very basic and intuitive.


> \textbf{Comparison for Convergence and Divergence}

> If two series, $S= \sum _{n=j}^{\infty}{s_n}$ and $ Z= \sum _{n=j}^{\infty}{z_n}$, 

$\newline$ and if $0 \leq z_n \leq s_n$ in the interval $[j, \infty)$, then if

>

> 1.  $Z$ is divergent, so is $S$

> 2.  $S$ is convergent, so is $Z$


First, a few words about this test. Notice that this test applies even if the two series' summands are \textit{equal}. This is because if summands are the same, this means that the series must also be the same, and so if one of them converges or diverges by the equality property they must both converge or diverge. However, if the starting point $j$ is different from series to series, then they \textit{will not converge to the same value}, that is to say $Z \neq S$, but this test will still apply.


The test itself follows from the fact that if we know that $S$ converges to some finite number, and we know that $z_n$ is less than (or equal to) $s_n$ for all $n$ then it follows that $Z$ should also converge to some finite number greater than zero. i.e., if there is a sum $1+2+3+4$ and a sum $2+3+4+5$ then we know that the first sum will be smaller because it has smaller numbers; the only thing smaller than a finite number is another finite number. The same is true for the divergence portion of this test. If $Z$ diverges and $z_n$ is less than or equal to $s_n$ for all $n$ then $S$ will diverge for essentially the same reason: the summand is bigger, and the sum of a set of numbers greater than the sum of numbers that is infinite must also be infinite as there is no finite number larger than an infinite number.


One last key note is that all the terms $z_n$ and $s_n$ must be larger than zero in order for this test to be conclusive. The series $\sum _{n=1}^{\infty}{\frac{(-1)^n}{n}}$ cannot be tested with the comparison test because it is alternating and half the terms are less than zero.


%-------------------------------------------------------------------------------
\subsubsection{Example 1}
\label{sss:exam-acc3}
%-------------------------------------------------------------------------------


Use the divergent and monotonic harmonic series

$\sum _{n=1}^{\infty}{\frac{1}{n}}$ to determine if the following series

are divergent or if the test is inconclusive.


1.  $\sum _{n=1}^{\infty}{\frac{1}{n+1}}$

2.  $\sum _{n=2}^{\infty}{\frac{1}{n-1}}$

3.  $\sum _{n=1}^{\infty}{\frac{3}{n}}$

4.  $\sum _{n=1}^{\infty}{\frac{1}{\sqrt{n}}}$

5.  $\sum _{n=1}^{\infty}{\frac{1}{n^2}}$


%-------------------------------------------------------------------------------
\subsubsection{Solutions}
\label{sss:solu-8692}
%-------------------------------------------------------------------------------


1.  Notice that this sum can be rewritten as

    $\sum _{n=2}^{\infty}{\frac{1}{n}}$, making it have the same summand as the harmonic series which is divergent; therefore, this series is divergent.

2.  This series is similar: it can be rewritten as

    $\sum _{n=1}^{\infty}{\frac{1}{n}}$ which is the harmonic series and so it is divergent.

3.  For each $n$, this series is larger because $3$ divided by any integer is larger than $1$ divided by any integer. This can also be seen as $3 \times \sum _{n=1}^{\infty}{\frac{1}{n}}$, which is essentially $3 \times \infty$ and so this series is divergent.

4.  For every $n$ in this series, the summand of

    $\sum _{n=1}^{\infty}{\frac{1}{\sqrt{n}}}$ is larger than the summand     of the harmonic series and so this series is divergent. This can be     seen by simply plotting the graph of $f(n)=\frac{1}{\sqrt{n}}$. Something interesting to note is that when $n < 1$, the summand of     the harmonic series is actually larger.

5.  Via plotting/plugging in values of $n$, we see that for every $n$ in

    the series, the summand of the harmonic series is larger and so the

    test fails and is inconclusive.


%-------------------------------------------------------------------------------
\subsubsection{Example 2}
\label{sss:exam-5bdd}
%-------------------------------------------------------------------------------

Use the convergent geometric series
\[
\sum_{n=1}^{\infty} \frac{1}{2^n}
\]
to determine whether the following series converge or whether the test is inconclusive.

\begin{enumerate}
\item $\displaystyle \sum_{n=1}^{\infty} \frac{1}{n^2}$
\item $\displaystyle \sum_{n=1}^{\infty} e^{-n}$
\item $\displaystyle \sum_{n=1}^{\infty} \frac{\sin^2(n)}{2^n}$
\item $\displaystyle \sum_{n=1}^{\infty} \frac{2}{2^n}$
\item $\displaystyle \sum_{n=1}^{\infty} \frac{1}{1.5^n}$
\item $\displaystyle \sum_{n=1}^{\infty} \frac{(-1)^n}{2^n}$
\end{enumerate}

%-------------------------------------------------------------------------------
\subsubsection{Solutions}
\label{sss:solu-8692}
%-------------------------------------------------------------------------------

\begin{enumerate}

\item The comparison with $\sum \frac{1}{2^n}$ is not appropriate here. 
However, $\sum \frac{1}{n^2}$ is a $p$-series with $p=2>1$, so it converges.

\item Since $e^{-n} = (1/e)^n$ and $1/e < 1$, this is a geometric series. 
Hence it converges.

\item Since $0 \leq \sin^2(n) \leq 1$ for all $n$, we have
\[
0 \leq \frac{\sin^2(n)}{2^n} \leq \frac{1}{2^n}.
\]
By the comparison test, the series converges.

\item This is
\[
\sum_{n=1}^{\infty} \frac{2}{2^n}
= 2 \sum_{n=1}^{\infty} \frac{1}{2^n},
\]
which converges because it is a constant multiple of a convergent geometric series.

\item Since $1.5 < 2$, we have
\[
\frac{1}{1.5^n} > \frac{1}{2^n}.
\]
Thus, comparison with $\sum \frac{1}{2^n}$ is inconclusive. However, this is a geometric series with ratio $1/1.5 < 1$, so it converges.

\item The terms are not nonnegative, so the comparison test does not apply. 
However, since
\[
\sum_{n=1}^{\infty} \left|\frac{(-1)^n}{2^n}\right| = \sum_{n=1}^{\infty} \frac{1}{2^n}
\]
converges, the series converges absolutely.

\end{enumerate}


%-------------------------------------------------------------------------------
\subsubsection{Resources}
\label{sss:reso-55b5}
%-------------------------------------------------------------------------------


%-------------------------------------------------------------------------------
\subsubsection{The Direct Comparison Test}
\label{sss:thed-585f}
%-------------------------------------------------------------------------------


\begin{itemize}
    \item \href{https://www.youtube.com/watch?v=2J9vhLkA1Xo}{The Direct Comparison
\end{itemize}
    Test} Video by James Sousa, Math is Power 4U

\begin{itemize}
    \item \href{https://www.youtube.com/watch?v=aJGSLBU4WuY}{Infinite Series - The Direct Comparison
\end{itemize}
    Test} Video by James Sousa, Math is Power 4U

\begin{itemize}
    \item \href{https://www.youtube.com/watch?v=X6zp_ybeTAQ}{Ex - Direct Comparison Test
\end{itemize}
    (Convergent)} Video by James Sousa, Math is Power 4U

\begin{itemize}
    \item \href{https://www.youtube.com/watch?v=_R85cPUCUKc}{Ex - Direct Comparison Test
\end{itemize}
    (Divergent)} Video by James Sousa, Math is Power 4U

\begin{itemize}
    \item \href{https://www.youtube.com/watch?v=0TZEfrMb3q0}{Ex - Direct Comparison Test
\end{itemize}
    (Inconclusive)} Video

    by James Sousa, Math is Power 4U





\begin{itemize}
    \item \href{https://www.youtube.com/watch?v=O-Quebw2e-8}{Direct Comparison Test - Another Example
\end{itemize}
    1} Video by Patrick JMT

\begin{itemize}
    \item \href{https://www.youtube.com/watch?v=ZTGeBm6xmy8}{Direct Comparison Test - Another Example
\end{itemize}
    2} Video by Patrick JMT

\begin{itemize}
    \item \href{https://www.youtube.com/watch?v=hpF3wqiPHQU}{Direct Comparison Test - Another Example
\end{itemize}
    3} Video by Patrick JMT

\begin{itemize}
    \item \href{https://www.youtube.com/watch?v=EZf5U6YaXXc}{Direct Comparison Test - Another Example
\end{itemize}
    4} Video by Patrick JMT

\begin{itemize}
    \item \href{https://www.youtube.com/watch?v=zrGC3nawUJA}{Direct Comparison Test - Another Example
\end{itemize}
    5} Video by Patrick JMT





\begin{itemize}
    \item \href{https://www.youtube.com/watch?v=ctJGrBzGoss}{Direct Comparison
\end{itemize}
    Test} Video by Krista King





\begin{itemize}
    \item \href{https://www.youtube.com/watch?v=oZtAgihok5s}{Direct Comparison
\end{itemize}
    Test} Video by The Organic Chemistry Tutor


%-------------------------------------------------------------------------------
\subsubsection{The Limit Comparison Test}
\label{sss:thel-318b}
%-------------------------------------------------------------------------------


\begin{itemize}
    \item \href{https://www.youtube.com/watch?v=LQ_ktobTEp0}{The Limit Comparison
\end{itemize}
    Test} Video by James Sousa, Math is Power 4U

\begin{itemize}
    \item \href{https://www.youtube.com/watch?v=uYC64L_UEJs}{Ex - Limit Comparison Test
\end{itemize}
    (Convergent)} Video by James Sousa, Math is Power 4U

\begin{itemize}
    \item \href{https://www.youtube.com/watch?v=5G8BPDhQsW8}{Ex - Limit Comparison Test
\end{itemize}
    (Convergent)} Video by James Sousa, Math is Power 4U

\begin{itemize}
    \item \href{https://www.youtube.com/watch?v=NDaxM6bJL90}{Ex - Limit Comparison Test
\end{itemize}
    (Convergent)} Video by James Sousa, Math is Power 4U

\begin{itemize}
    \item \href{https://www.youtube.com/watch?v=Q77TLOQ9Y9k}{Ex - Limit Comparison Test
\end{itemize}
    (Divergent)} Video by James Sousa, Math is Power 4U

\begin{itemize}
    \item \href{https://www.youtube.com/watch?v=cXhwDCJ2HiE}{Ex - Limit Comparison Test
\end{itemize}
    (Divergent)} Video by James Sousa, Math is Power 4U

\begin{itemize}
    \item \href{https://www.youtube.com/watch?v=mzrmkveJ8nk}{Ex - Limit Comparison Test
\end{itemize}
    (Divergent)} Video by James Sousa, Math is Power 4U





\begin{itemize}
    \item \href{https://www.youtube.com/watch?v=Hg5RwzLAqRs}{Limit Comparison Test - Another Example
\end{itemize}
    1} Video by Patrick JMT

\begin{itemize}
    \item \href{https://www.youtube.com/watch?v=gNSWa3VDx9k}{Limit Comparison Test - Another Example
\end{itemize}
    2} Video by Patrick JMT

\begin{itemize}
    \item \href{https://www.youtube.com/watch?v=dX_bwYp-lVc}{Limit Comparison Test - Another Example
\end{itemize}
    3} Video by Patrick JMT

\begin{itemize}
    \item \href{https://www.youtube.com/watch?v=PxtE4reKltY}{Limit Comparison Test - Another Example
\end{itemize}
    5} Video by Patrick JMT

\begin{itemize}
    \item \href{https://www.youtube.com/watch?v=-uU5QbQALcM}{Limit Comparison Test - Another Example
\end{itemize}
    6} Video by Patrick JMT

\begin{itemize}
    \item \href{https://www.youtube.com/watch?v=QAZZgGDnnCg}{Limit Comparison Test - Another Example
\end{itemize}
    7} Video by Patrick JMT

\begin{itemize}
    \item \href{https://www.youtube.com/watch?v=PIr6NMbfTCM}{Limit Comparison Test - Another Example
\end{itemize}
    8} Video by Patrick JMT





\begin{itemize}
    \item \href{https://www.youtube.com/watch?v=PPTLtFhAp6c}{Limit Comparison Test}
\end{itemize}
    Video by Krista King





\begin{itemize}
    \item \href{https://www.youtube.com/watch?v=LBxYQ0TJxYM}{Limit Comparison Test}
\end{itemize}
    Video by The Organic Chemistry Tutor


%-------------------------------------------------------------------------------
\subsubsection{Licensing}
\label{sss:lice-7352}
%-------------------------------------------------------------------------------


Content obtained and/or adapted from:


\begin{itemize}
    \item \href{https://en.wikibooks.org/wiki/Calculus/Comparison_Test_for_Convergence}{Comparison Test for Convergence, Wikibooks:
\end{itemize}
    Calculus}

    under a CC BY-SA license

\subsection{Learning Objectives}

\begin{enumerate}
    \item Here are three student learning objectives (SLOs) related to the Comparison Test for series:
    \item \textbf{Apply the Comparison Test to Determine Convergence or Divergence:}
    \item \textbf{Outcome:} Students will be able to apply the Comparison Test to determine if a given series converges or diverges by comparing it to a known series.
    \item \textbf{Details:} This includes identifying appropriate series for comparison, ensuring that the conditions of the test are met (i.e., the comparison series must have positive terms), and correctly concluding the behavior of the given series based on the behavior of the comparison series.
    \item \textbf{Analyze and Justify Results Using Known Series:}
    \item \textbf{Outcome:} Students will be able to justify the convergence or divergence of a series by analyzing how the terms of the series relate to those of a known convergent or divergent series.
    \item \textbf{Details:} This involves using known series, such as the harmonic series or geometric series, to either show that the series in question is bounded above or below by a series with known behavior or to establish divergence through comparison with a divergent series.
    \item \textbf{Handle Special Cases and Inconclusive Results:}
    \item \textbf{Outcome:} Students will be able to recognize and handle cases where the Comparison Test is inconclusive or cannot be applied directly.
    \item \textbf{Details:} This includes understanding situations where the terms of the series to be tested are not all positive, where the series does not fit the criteria for comparison, or where additional adjustments or alternative tests may be needed to determine convergence or divergence.
\end{enumerate}

\subsection{Assessment Instrument}
\label{ssec:comp-6dff}

Students apply the Direct Comparison Test and the Limit Comparison Test, selecting an appropriate benchmark series (geometric or $p$-series) and verifying the required inequality or limit condition.

\textbf{Example problems.}

\begin{enumerate}
    \item Use the Direct Comparison Test to determine whether $\sum_{n=1}^{\infty}\dfrac{1}{n^2+1}$ converges.
    \item Apply the Limit Comparison Test to $\sum_{n=1}^{\infty}\dfrac{3n^2-1}{n^4+n+2}$ with an appropriate $p$-series.
    \item Determine convergence of $\sum_{n=1}^{\infty}\dfrac{\ln n}{n^2}$ using a suitable comparison.
\end{enumerate}

%===============================================================================
\section{Alternating series tests}
\label{sec:alte-c075}
%===============================================================================


%-------------------------------------------------------------------------------
\subsubsection{Alternating Series}
\label{sss:alte-4145}
%-------------------------------------------------------------------------------

An \textit{alternating series} is any series whose terms alternate in sign; 
that is, consecutive terms have opposite signs.

Equivalently, an alternating series can be written in the form
\[
\sum_{n=c}^{\infty} (-1)^{n+d} \, b_n,
\]
for some fixed integers $c$ and $d$, where $\{b_n\}$ is a sequence of 
positive terms. If the $n$th term of the series is $a_n$, then 
$b_n = |a_n|$.

In many calculus textbooks, alternating series are defined more restrictively as
\[
\sum_{n=1}^{\infty} (-1)^{n+1} b_n,
\]
or
\[
\sum_{n=1}^{\infty} (-1)^n b_n
\quad \text{or} \quad
\sum_{n=1}^{\infty} (-1)^{n-1} b_n,
\]
with $b_n > 0$. These are all special cases of the general form above.

\medskip

While it is generally difficult to find closed formulas for partial sums of 
alternating series, we can determine convergence using the following result.

\paragraph{Alternating Series Test}

An alternating series
\[
\sum_{n=1}^{\infty} (-1)^n b_n
\]
converges if:
\begin{enumerate}
\item $\displaystyle \lim_{n \to \infty} b_n = 0$,
\item $\{b_n\}$ is eventually decreasing.
\end{enumerate}

That is, the magnitudes of the terms must decrease to zero.

\medskip

\textbf{Important:} If either condition fails, the test is inconclusive. 
The series may still converge or diverge, and other tests are required.

\medskip

Note that any geometric series with ratio $-1 < r < 0$ is an alternating 
series and converges absolutely.

%-------------------------------------------------------------------------------
\subsubsection{Resources}
\label{sss:reso-alte}
%-------------------------------------------------------------------------------

\begin{itemize}

\item \href{https://www.youtube.com/watch?v=-K9Qt6YUIrI}{Alternating Series}. 
Video by Patrick JMT.

\item \href{https://www.youtube.com/watch?v=oJEhVFTos5o}{Alternating Series – Basic Example}. 
Video by Patrick JMT.

\item \href{https://www.youtube.com/watch?v=e4_dE0YiKlA}{Two Alternating Series Examples}. 
Video by Patrick JMT.

\item \href{https://www.youtube.com/watch?v=zgTRNpKcenE}{More Alternating Series Examples}. 
Video by Patrick JMT.

\item \href{https://www.youtube.com/watch?v=xL30cRRXWsI}{Alternating Series – Another Example 1}. 
Video by Patrick JMT.

\item \href{https://www.youtube.com/watch?v=SFs2ys2gpoE}{Alternating Series – Another Example 2}. 
Video by Patrick JMT.

\item \href{https://www.youtube.com/watch?v=XDC53pvATcs}{Alternating Series – Another Example 3}. 
Video by Patrick JMT.

\item \href{https://www.youtube.com/watch?v=aOiZvfFAMW8}{Alternating Series – Another Example 4}. 
Video by Patrick JMT.

\item \href{https://www.youtube.com/watch?v=iYNgFPFKWy0}{The Alternating Series Test}. 
Video by Krista King.

\item \href{https://www.youtube.com/watch?v=DRO1kPT4iS8}{The Alternating Series Test}. 
Video by The Organic Chemistry Tutor.

\end{itemize}

%-------------------------------------------------------------------------------
\subsubsection{The Alternating Series Estimation Theorem}
\label{sss:thea-9bbb}
%-------------------------------------------------------------------------------

\begin{itemize}

\item \href{https://www.youtube.com/watch?v=oZ3PKvQKffE}{Find the Error in Using a Partial Sum to Approximate the Sum of an Alternating Series}. 
Video by James Sousa, Math is Power 4U.

\item \href{https://www.youtube.com/watch?v=_hWa3INOo-Y}{Number of Terms Needed in a Partial Sum to Approximate the Sum of an Alternating Series with Given Error}. 
Video by James Sousa, Math is Power 4U.

\item \href{https://www.youtube.com/watch?v=ilRkVcAxdCo}{Alternating Series Estimation Theorem}. 
Video by Patrick JMT.

\item \href{https://www.youtube.com/watch?v=BNEULxileT8}{Alternating Series Error Estimation}. 
Video by Patrick JMT.

\item \href{https://www.youtube.com/watch?v=H8lHvQashkE}{Alternating Series Estimation Theorem}. 
Video by Krista King.

\item \href{https://www.youtube.com/watch?v=FkUrAgBzAZo}{Alternating Series Estimation Theorem}. 
Video by The Organic Chemistry Tutor.

\end{itemize}

%-------------------------------------------------------------------------------
\subsubsection{Absolute Convergence and Conditional Convergence}
\label{sss:abso-c62f}
%-------------------------------------------------------------------------------

\begin{itemize}

\item \href{https://www.youtube.com/watch?v=rEyeRKd3TWo}{Absolutely and Conditionally Convergent Series}. 
Video by James Sousa, Math is Power 4U.

\item \href{https://www.youtube.com/watch?v=NHkxgvXVH7g}{Ex 1: Determine if a Series is Absolutely Convergent, Conditionally Convergent, or Divergent}. 
Video by James Sousa, Math is Power 4U.

\item \href{https://www.youtube.com/watch?v=awNWNGkD6Vk}{Ex 2: Determine if a Series is Absolutely Convergent, Conditionally Convergent, or Divergent}. 
Video by James Sousa, Math is Power 4U.

\item \href{https://www.youtube.com/watch?v=kP9sMcirMqw}{Ex 3: Determine if a Series is Absolutely Convergent, Conditionally Convergent, or Divergent}. 
Video by James Sousa, Math is Power 4U.

\item \href{https://www.youtube.com/watch?v=xGQ-sR2KLwc}{Ex 4: Determine if a Series is Absolutely Convergent, Conditionally Convergent, or Divergent}. 
Video by James Sousa, Math is Power 4U.

\item \href{https://www.youtube.com/watch?v=6hOeqjoQvNw}{Absolute Convergence, Conditional Convergence, and Divergence}. 
Video by Patrick JMT.

\item \href{https://www.youtube.com/watch?v=UvWgxCSypFc}{Absolute Convergence, Conditional Convergence – Another Example 1}. 
Video by Patrick JMT.

\item \href{https://www.youtube.com/watch?v=XQLa66wZYI4}{Absolute Convergence, Conditional Convergence – Another Example 2}. 
Video by Patrick JMT.

\item \href{https://www.youtube.com/watch?v=_OYmHjNQV1A}{Absolute Convergence, Conditional Convergence – Another Example 3}. 
Video by Patrick JMT.

\item \href{https://www.youtube.com/watch?v=8h9Eo148IF4}{Absolute and Conditional Convergence}. 
Video by Krista King.

\item \href{https://www.youtube.com/watch?v=FPK6LO1iiXc}{Absolute Convergence, Conditional Convergence, and Divergence}. 
Video by The Organic Chemistry Tutor.

\end{itemize}

%-------------------------------------------------------------------------------
\subsubsection{Licensing}
\label{sss:lice-7352}
%-------------------------------------------------------------------------------

Content obtained and/or adapted from:

\begin{itemize}

\item \href{https://en.wikibooks.org/wiki/User:Dcljr/Series}{Series, Wikibooks}. 
Licensed under CC BY-SA.

\end{itemize}

\subsection{Learning Objectives}

\begin{enumerate}
    \item Here are three student learning objectives (SLOs) related to alternating series:
    \item \textbf{Determine Convergence of Alternating Series:}
    \item \textbf{Outcome:} Students will be able to determine the convergence of an alternating series using the Alternating Series Test (Leibniz's Test).
    \item \textbf{Details:} This includes verifying that the terms of the series decrease in magnitude to zero and understanding the criteria for an alternating series to converge.
    \item \textbf{Differentiate Between Absolute and Conditional Convergence:}
    \item \textbf{Outcome:} Students will be able to differentiate between absolute and conditional convergence of an alternating series.
    \item \textbf{Details:} This includes understanding the definitions of absolute and conditional convergence and applying these concepts to determine whether a given alternating series converges absolutely or conditionally.
    \item \textbf{Analyze Convergence of Specific Alternating Series:}
    \item \textbf{Outcome:} Students will be able to analyze the convergence of specific alternating series through examples.
    \item \textbf{Details:} This includes working through examples such as the alternating harmonic series and geometric series with negative ratios, and applying the Alternating Series Test to determine their convergence properties.
\end{enumerate}

\subsection{Assessment Instrument}
\label{ssec:alte-69b6}

Assessment covers the Alternating Series Test (Leibniz criterion), the alternating series estimation theorem for error bounds, and the distinction between absolute and conditional convergence.

\textbf{Example problems.}

\begin{enumerate}
    \item Use the Alternating Series Test to determine whether $\sum_{n=1}^{\infty}\dfrac{(-1)^{n+1}}{n}$ converges.
    \item For the alternating series $\sum_{n=1}^{\infty}\dfrac{(-1)^n}{n^2}$, how many terms are needed to approximate the sum with error less than $0.01$?
    \item Determine whether $\sum_{n=1}^{\infty}\dfrac{(-1)^n}{\sqrt{n}}$ converges absolutely, conditionally, or diverges.
\end{enumerate}

%===============================================================================
\section{Ratio and Root tests}
\label{sec:rati-81b0}
%===============================================================================


%-------------------------------------------------------------------------------
\subsubsection{Ratio Test}
\label{sss:rati-a950}
%-------------------------------------------------------------------------------


In mathematics, the \textbf{ratio test} is a test (or ``criterion'') for the convergence of a series


\[
\sum _{n=1}^\infty a_n,
\]


where each term is a real or complex number and $a_n$ is nonzero when $n$ is large. The test was first published by Jean le Rond d'Alembert and is sometimes known as \textbf{d'Alembert's ratio test} or as the \textbf{Cauchy ratio test}.


%-------------------------------------------------------------------------------
\subsubsection{The Test}
\label{sss:thet-762c}
%-------------------------------------------------------------------------------


% [Image: Decision diagram for the ratio test]


Figure: Decision diagram for the ratio test


The usual form of the test makes use of the limit


\[
L = \lim _{n\to\infty}\left|\frac{a _{n+1}}{a_n}\right|.
\]


The ratio test states that:


\begin{itemize}
    \item if $L < 1$ then the series converges absolutely;
    \item if $L > 1$ then the series is divergent;
    \item if $L = 1$ or the limit fails to exist, then the test is inconclusive, because there exist both convergent and divergent series that satisfy this case.
\end{itemize}

It is possible to make the ratio test applicable to certain cases where the limit \textit{L} fails to exist, if limit superior and limit inferior are used. The test criteria can also be refined so that the test is sometimes conclusive even when $L = 1$. More specifically, let


\[
R = \lim\sup \left|\frac{a _{n+1}}{a_n}\right|
\]


\[
r = \lim\inf \left|\frac{a _{n+1}}{a_n}\right|
\].


Then the ratio test states that:


\begin{itemize}
    \item if $R < 1$, the series converges absolutely;
    \item if $r > 1$, the series diverges;
    \item if $\left|\frac{a _{n+1}}{a_n}\right|\ge 1$ for all large \textit{n} (regardless of the value of \textit{r}), the series also diverges; this is because $|a_n|$ is nonzero and increasing and hence $a_n$ does not approach zero;
    \item the test is otherwise inconclusive.
\end{itemize}

If the limit $L$ in (1) exists, we must have $L = R = r$. So the original ratio test is a weaker version of the refined one.


%-------------------------------------------------------------------------------
\subsubsection{Examples}
\label{sss:exam-bfeb}
%-------------------------------------------------------------------------------


%-------------------------------------------------------------------------------
\subsubsection{Convergent because $L < 1$}
\label{sss:conv-9602}
%-------------------------------------------------------------------------------


Consider the series


\[
\sum _{n=1}^\infty\frac{n}{e^n}
\]


Applying the ratio test, one computes the limit


\[
L = \lim _{n\to\infty} \left| \frac{a _{n+1}}{a_n} \right| = \lim _{n\to\infty} \left| \frac{\frac{n+1}{e^{n+1}}}{\frac{n}{e^n}}\right| = \frac{1}{e} < 1.
\]


Since this limit is less than 1, the series converges.


%-------------------------------------------------------------------------------
\subsubsection{Divergent because \textit{L}
\label{sss:dive-77ed}
%-------------------------------------------------------------------------------
 \> 1}


Consider the series


\[
\sum _{n=1}^\infty\frac{e^n}{n}.
\]


Putting this into the ratio test:


\[
L = \lim _{n\to\infty} \left| \frac{a _{n+1}}{a_n} \right| = \lim _{n\to\infty} \left| \frac{\frac{e^{n+1}}{n+1}}{\frac{e^n}{n}} \right|

= e > 1.
\]


Thus the series diverges.


%-------------------------------------------------------------------------------
\subsubsection{Inconclusive because $L = 1$}
\label{sss:inco-32dd}
%-------------------------------------------------------------------------------


Consider the three series


\[
\sum _{n=1}^\infty 1,
\]


\[
\sum _{n=1}^\infty \frac{1}{n^2},
\]


\[
\sum _{n=1}^\infty \frac{(-1)^{n+1}}{n}.
\]


The first series (1 + 1 + 1 + 1 + $\cdots$) diverges, the second one (the one central to the Basel problem) converges absolutely and the third one (the alternating harmonic series) converges conditionally. However, the term-by-term magnitude ratios $\left|\frac{a _{n+1}}{a_n}\right|$ of the three series are respectively $1,$ $\frac{n^2}{(n+1)^2}$ and $\frac{n}{n+1}$. So, in all three cases, one has that the limit $\lim _{n\to\infty}\left|\frac{a _{n+1}}{a_n}\right|$ is equal to 1. This illustrates that when \textit{L} = 1, the series may converge or diverge, and hence the original ratio test is inconclusive. In such cases, more refined tests are required to determine convergence or divergence.


%-------------------------------------------------------------------------------
\subsubsection{Proof}
\label{sss:proo-028d}
%-------------------------------------------------------------------------------


% [Image: Sequence Convergence]


Figure: In this example, the ratio of adjacent terms in the blue sequence converges to $L=1/2$. We choose $r = (L+1)/2 = 3/4$. Then the blue sequence is dominated by the red sequence $rk$ for all $n \geq 2$. The red sequence converges, so the blue sequence does as well.


Below is a proof of the validity of the original ratio test.


Suppose that

$L = \lim _{n\to\infty} \left| \frac{a _{n+1}}{a _{n}}\right| < 1$. We can then show that the series converges absolutely by showing that its terms will eventually become less than those of a certain convergent geometric series. To do this, consider a real number $r$ \textit{such that} $L < r < 1$. This implies that $|a _{n+1}| < r |a _{n}|$ for sufficiently large \textit{n}; say, for all \textit{n} greater than \textit{N}. Hence $|a _{n+i}| < r^i|a _{n}|$ for each \textit{n} \> \textit{N} and \textit{i} \> 0, and so


\[
\sum _{i=N+1}^{\infty}|a_i| = \sum _{i=1}^{\infty} \left |a _{N+i} \right | < \sum _{i=1}^{\infty}r^{i}|a _{N}| = |a _{N}| \sum _{i=1}^{\infty} r^{i} = |a _{N}|\frac{r}{1 - r} < \infty.
\]


That is, the series converges absolutely.


On the other hand, if \textit{L} \> 1, then $|a _{n+1}| > |a _{n}|$ for sufficiently large \textit{n}, so that the limit of the summands is non-zero. Hence the series diverges.


%-------------------------------------------------------------------------------
\subsubsection{Root Test}
\label{sss:root-9d7b}
%-------------------------------------------------------------------------------


In mathematics, the \textbf{root test} is a criterion for the convergence (a convergence test) of an infinite series. It depends on the quantity


\[
\limsup _{n\rightarrow\infty}\sqrt[n]{|a_n|},
\] where $a_n$ are the terms of the series, and states that the series converges absolutely if this quantity is less than one, but diverges if it is greater than one. It is particularly useful in connection with power series.


%-------------------------------------------------------------------------------
\subsubsection{Root test explanation}
\label{sss:root-bf3f}
%-------------------------------------------------------------------------------


% [Image: Decision diagram for the root test]


Figure: Decision diagram for the root test


The root test was developed first by Augustin-Louis Cauchy who published it in his textbook Cours d'analyse (1821). Thus, it is sometimes known as the \textbf{Cauchy root test} or \textbf{Cauchy's radical test}. For a series


\[
\sum _{n=1}^\infty a_n.
\]


the root test uses the number


\[
C = \limsup _{n\rightarrow\infty}\sqrt[n]{|a_n|},
\]


where ``lim sup'' denotes the limit superior, possibly 8+. Note that if


\[
\lim _{n\rightarrow\infty}\sqrt[n]{|a_n|},
\]


converges then it equals \textit{C} and may be used in the root test instead.


The root test states that:


\begin{itemize}
    \item if $C < 1$ then the series converges absolutely,
    \item if $C > 1$ then the series diverges,
    \item if $C = 1$ and the limit approaches strictly from above then the series diverges,
    \item otherwise the test is inconclusive (the series may diverge, converge absolutely or converge conditionally).
\end{itemize}

There are some series for which \textit{C} = 1 and the series converges, e.g. $\textstyle \sum 1/{n^2}$, and there are others for which \textit{C} = 1 and the series diverges, e.g. $\textstyle\sum 1/n$.


%-------------------------------------------------------------------------------
\subsubsection{Application to power series}
\label{sss:appl-3c8a}
%-------------------------------------------------------------------------------


This test can be used with a power series


\[
f(z) = \sum _{n=0}^\infty c_n (z-p)^n
\]


where the coefficients $c_n$, and the center $p$ are complex numbers and the argument $z$ is a complex variable.


The terms of this series would then be given by $a_n = c_n (z - p) ^n$. One then applies the root test to the $a_n$ as above.


Note that sometimes a series like this is called a power series ``around \textit{p}'', because the radius of convergence is the radius \textit{R} of the largest interval or disc centred at \textit{p} such that the series will converge for all points \textit{z} strictly in the interior (convergence on the boundary of the interval or disc generally has to be checked separately). A corollary of the root test applied to such a power series is the Cauchy--Hadamard theorem: the radius of convergence is exactly $1/\limsup _{n \rightarrow \infty}{\sqrt[n]{|c_n|}},$ taking care that we really mean $\infty$ if the denominator is 0.


%-------------------------------------------------------------------------------
\subsubsection{Proof}
\label{sss:proo-028d}
%-------------------------------------------------------------------------------


The proof of the convergence of a series $\sum a_n$ is an application of the comparison test. If for all \textit{n} = \textit{N} (\textit{N} some fixed natural number) we have $\sqrt[n]{|a_n|} \le k < 1,$ then $|a_n| \le k^n < 1$. Since the geometric series $\sum _{n=N}^\infty k^n$ converges so does $\sum _{n=N}^\infty |a_n|$ by the comparison test. Hence $\sum a_n$ converges absolutely.


If $\sqrt[n]{|a_n|} > 1$ for infinitely many $n$, then $a_n$ fails to converge to 0, hence the series is divergent.


\textbf{Proof of corollary}: For a power series $\sum a_n = \sum c_n (z - p)^n$, we see by the above that the series converges if there exists an \textit{N} such that for all $n \geq N$ we have


\[
\sqrt[n]{|a_n|} = \sqrt[n]{|c_n(z - p)^n|} < 1,
\]


equivalent to


\[
\sqrt[n]{|c_n|}\cdot|z - p| < 1
\]


for all $n \geq N$, which implies that in order for the series to converge we must have $|z - p| < 1/\sqrt[n]{|c_n|}$ for all sufficiently large $n$. This is equivalent to saying


\[
|z - p| < 1/\limsup _{n \rightarrow \infty}{\sqrt[n]{|c_n|}},
\]


so $R \le 1/\limsup _{n \rightarrow \infty}{\sqrt[n]{|c_n|}}.$ Now the only other place where convergence is possible is when


\[
\sqrt[n]{|a_n|} = \sqrt[n]{|c_n(z - p)^n|} = 1,
\]


(since points > 1 will diverge) and this will not change the radius of convergence since these are just the points lying on the boundary of the interval or disc, so


\[
R = 1/\limsup _{n \rightarrow \infty}{\sqrt[n]{|c_n|}}.
\]


%-------------------------------------------------------------------------------
\subsubsection{Examples}
\label{sss:exam-bfeb}
%-------------------------------------------------------------------------------

\textit{Example 1:}

\[
\sum_{n=1}^{\infty} \frac{2^n}{n^9}
\]

Applying the root test and using $\lim_{n \to \infty} n^{1/n} = 1$,
\[
C = \lim_{n \to \infty} \sqrt[n]{\left|\frac{2^n}{n^9}\right|}
= \lim_{n \to \infty} \frac{\sqrt[n]{2^n}}{\sqrt[n]{n^9}}
= \lim_{n \to \infty} \frac{2}{(n^{1/n})^9}
= 2.
\]

Since $C = 2 > 1$, the series diverges.

\medskip

\textit{Example 2:}

\[
1 + 1 + \tfrac{1}{2} + \tfrac{1}{2} + \tfrac{1}{4} + \tfrac{1}{4} + \cdots
\]

Observe that this is not a standard geometric series, but it behaves like one. 
Define the general term:
\[
a_n \sim \left(\tfrac{1}{2}\right)^{\lfloor n/2 \rfloor}.
\]

Applying the root test,
\[
\limsup_{n \to \infty} \sqrt[n]{|a_n|} = \sqrt{\tfrac{1}{2}} < 1.
\]

Hence, the series converges.

%-------------------------------------------------------------------------------
\subsubsection{Resources}
\label{sss:reso-ratio-root}
%-------------------------------------------------------------------------------

\textbf{Ratio Test}

\begin{itemize}

\item \href{https://www.youtube.com/watch?v=iy8mhbZTY7g}{Using the Ratio Test to Determine if a Series Converges \#1}. 
Video by Patrick JMT.

\item \href{https://www.youtube.com/watch?v=gay5CEXnkhA}{Using the Ratio Test to Determine if a Series Converges \#2}. 
Video by Patrick JMT.

\item \href{https://www.youtube.com/watch?v=AwJ0P8B25tY}{Using the Ratio Test to Determine if a Series Converges \#3}. 
Video by Patrick JMT.

\item \href{https://www.youtube.com/watch?v=nFsoA_dArn4}{The Ratio Test – Another Example \#1}. 
Video by Patrick JMT.

\item \href{https://www.youtube.com/watch?v=KjAkB_2XHaQ}{The Ratio Test – Another Example \#2}. 
Video by Patrick JMT.

\item \href{https://www.youtube.com/watch?v=4H57ndz-FRA}{The Ratio Test – Another Example \#3}. 
Video by Patrick JMT.

\item \href{https://www.youtube.com/watch?v=U5a3Ls4ZhyA}{The Ratio Test – Another Example \#4}. 
Video by Patrick JMT.

\item \href{https://www.youtube.com/watch?v=WeFse6u19fA}{The Ratio Test}. 
Video by Krista King.

\item \href{https://www.youtube.com/watch?v=uMLBA6CHR1E}{The Ratio Test with Factorials}. 
Video by Krista King.

\item \href{https://www.youtube.com/watch?v=kkfILYpkJS8}{The Ratio Test}. 
Video by The Organic Chemistry Tutor.

\end{itemize}

\textbf{Root Test}

\begin{itemize}

\item \href{https://www.youtube.com/watch?v=ZE5Pj7s9PWg}{The Root Test}. 
Video by Math is Power 4U.

\item \href{https://www.youtube.com/watch?v=eHzBsFxBzpA}{The Root Test I}. 
Video by Math is Power 4U.

\item \href{https://www.youtube.com/watch?v=eMMtyG7hIyI}{The Root Test II}. 
Video by Math is Power 4U.

\item \href{https://www.youtube.com/watch?v=yVva1LwaE9c}{Ex 1: The Root Test}. 
Video by Math is Power 4U.

\item \href{https://www.youtube.com/watch?v=oKUfBJF5XTI}{Ex 2: The Root Test}. 
Video by Math is Power 4U.

\item \href{https://www.youtube.com/watch?v=Xxzs4sa_zZg}{Ex 3: The Root Test}. 
Video by Math is Power 4U.

\item \href{https://www.youtube.com/watch?v=5HLqCSEgaFI}{Ex 4: The Root Test}. 
Video by Math is Power 4U.

\item \href{https://www.youtube.com/watch?v=JEDrZfwsW1M}{Ex 5: The Root Test}. 
Video by Math is Power 4U.

\item \href{https://www.youtube.com/watch?v=vDdDLfIya0I}{The Root Test for Series}. 
Video by Patrick JMT.

\item \href{https://www.youtube.com/watch?v=nTfJCbqHyxI}{The Root Test – Another Example 1}. 
Video by Patrick JMT.

\item \href{https://www.youtube.com/watch?v=s2ox6IuKiX0}{The Root Test – Another Example 2}. 
Video by Patrick JMT.

\item \href{https://www.youtube.com/watch?v=wo4c4NYIiIo}{The Root Test – Another Example 3}. 
Video by Patrick JMT.

\item \href{https://www.youtube.com/watch?v=pSvQaiqTajk}{The Root Test}. 
Video by Krista King.

\item \href{https://www.youtube.com/watch?v=ahf0eXOll1M}{The Root Test}. 
Video by The Organic Chemistry Tutor.

\end{itemize}

%-------------------------------------------------------------------------------
\subsubsection{References}
\label{sss:refe-44a2}
%-------------------------------------------------------------------------------

\begin{enumerate}
\item Rudin, W. (1976). \textit{Principles of Mathematical Analysis} (3rd ed.). McGraw-Hill.
\item Apostol, T. (1974). \textit{Mathematical Analysis} (2nd ed.). Addison-Wesley.
\item Bottazzini, U. (1986). \textit{The Higher Calculus}. Springer.
\item Briggs, W., \& Cochran, L. (2011). \textit{Calculus: Early Transcendentals}.
\item Knopp, K. (1956). \textit{Infinite Sequences and Series}. Dover.
\item Whittaker, E. T., \& Watson, G. N. (1963). \textit{A Course in Modern Analysis}.
\end{enumerate}

%-------------------------------------------------------------------------------
\subsubsection{Licensing}
\label{sss:lice-7352}
%-------------------------------------------------------------------------------

Content obtained and/or adapted from:

\begin{itemize}
\item \href{https://en.wikipedia.org/wiki/Root_test}{Root test, Wikipedia}. Licensed under CC BY-SA.
\item \href{https://en.wikipedia.org/wiki/Ratio_test}{Ratio test, Wikipedia}. Licensed under CC BY-SA.
\end{itemize}

\subsection{Assessment Instrument}
\label{ssec:rati-8609}

Problems are chosen so that the ratio or root test gives a definitive answer ($L \ne 1$). Students must compute the limit carefully, state the conclusion, and identify when the test is inconclusive.

\textbf{Example problems.}

\begin{enumerate}
    \item Apply the Ratio Test to $\sum_{n=1}^{\infty}\dfrac{n!}{n^n}$ and determine convergence or divergence.
    \item Use the Root Test on $\sum_{n=1}^{\infty}\left(\dfrac{2n}{n+1}\right)^n$.
    \item Apply the Ratio Test to $\sum_{n=1}^{\infty}\dfrac{(-1)^n 3^n}{n\cdot 2^n}$ and discuss the result.
\end{enumerate}

%===============================================================================
\section{Power Series}
\label{sec:power-series}
%===============================================================================

The study of \textbf{power series} is aimed at investigating series which can

approximate some function over a certain interval.


%-------------------------------------------------------------------------------
\subsubsection{Motivations}
\label{sss:moti-f8c2}
%-------------------------------------------------------------------------------


Elementary calculus (differentiation) is used to obtain information on a

line which touches a curve at one point (i.e. a tangent). This is done

by calculating the gradient, or slope of the curve, at a single point.

However, this does not provide us with reliable information on the

curve's actual \textit{value} at given points in a wider interval. This is

where the concept of power series becomes useful.


%-------------------------------------------------------------------------------
\subsubsection{An example}
\label{sss:anex-9765}
%-------------------------------------------------------------------------------


Consider the curve of $y=\cos(x)$ , about the point $x=0$ . A naïve

approximation would be the line $y=1$ . However, for a more accurate

approximation, observe that $\cos(x)$ looks like an inverted parabola

around $x=0$ - therefore, we might think about which parabola could

approximate the shape of $\cos(x)$ near this point. This curve might

well come to mind:


\[
y=1-\frac{x^2}{2}
\] In fact, this is the best estimate for $\cos(x)$

which uses polynomials of degree 2 (i.e. a highest term of $x^2$) - but

how do we know this is true? This is the study of power series: finding

optimal approximations to functions using polynomials.


%-------------------------------------------------------------------------------
\subsubsection{Definition}
\label{sss:defi-3061}
%-------------------------------------------------------------------------------


A \textit{power series} (in one variable) is a infinite series of the form


\[
f(x)=a_0(x-c)^0+a_1(x-c)^1+a_2(x-c)^2+\cdots
\] (where $c$ is a

constant) or, equivalently,


\[
f(x)=\sum _{n=0}^\infty a_n(x-c)^n
\]


%-------------------------------------------------------------------------------
\subsubsection{Radius of convergence}
\label{sss:radi-7672}
%-------------------------------------------------------------------------------


When using a power series as an alternative method of calculating a

function's value, the equation


\[
f(x)=\sum _{n=0}^\infty a_n(x-c)^n
\] can only be used to study $f(x)$

where the power series converges - this may happen for a finite range,

or for all real numbers.


The size of the interval (around its center) in which the power series

converges to the function is known as the \textit{radius of convergence}.


%-------------------------------------------------------------------------------
\subsubsection{An Example}
\label{sss:anex-9765}
%-------------------------------------------------------------------------------

A fundamental example of a power series is the geometric series:
\[
\frac{1}{1 - x} = \sum_{n=0}^{\infty} x^n,
\]
which converges for
\[
|x| < 1.
\]

This provides a way to represent a function as an infinite polynomial.  
Similarly, we have the related series:
\[
\frac{1}{1 + x} = \sum_{n=0}^{\infty} (-1)^n x^n
= 1 - x + x^2 - x^3 + \cdots,
\]
which also converges for $|x| < 1$.

\medskip

These examples illustrate how functions can be expressed as power series 
within their \textbf{interval of convergence}.

%-------------------------------------------------------------------------------
\subsubsection{Resources}
\label{sss:reso-power}
%-------------------------------------------------------------------------------

\textbf{Radius and Interval of Convergence}

\begin{itemize}

\item \href{https://www.youtube.com/watch?v=01LzAU__J-0}{Interval and Radius of Convergence for a Series Ex 1}. 
Video by Patrick JMT.

\item \href{https://www.youtube.com/watch?v=QQBtF7_KJrE}{Interval and Radius of Convergence for a Series Ex 2}. 
Video by Patrick JMT.

\item \href{https://www.youtube.com/watch?v=-gQ2E33V77o}{Interval and Radius of Convergence for a Series Ex 3}. 
Video by Patrick JMT.

\item \href{https://www.youtube.com/watch?v=AZeAZsK0sOw}{Interval and Radius of Convergence for a Series Ex 4}. 
Video by Patrick JMT.

\item \href{https://www.youtube.com/watch?v=uhNQ4X8hXoY}{Interval and Radius of Convergence for a Series Ex 5}. 
Video by Patrick JMT.

\item \href{https://www.youtube.com/watch?v=hLaJU2BLS2Y}{Interval and Radius of Convergence for a Series Ex 6}. 
Video by Patrick JMT.

\item \href{https://www.youtube.com/watch?v=ze54foaEUEQ}{Interval and Radius of Convergence for a Series Ex 7}. 
Video by Patrick JMT.

\item \href{https://www.youtube.com/watch?v=h4CNG3iAxKc}{Interval and Radius of Convergence for a Series Ex 8}. 
Video by Patrick JMT.

\item \href{https://www.youtube.com/watch?v=tcIKlExw-zA}{Interval and Radius of Convergence for a Series Ex 9}. 
Video by Patrick JMT.

\item \href{https://www.youtube.com/watch?v=MM3BXtVu9eM}{Radius of Convergence for a Power Series}. 
Video by Patrick JMT.

\item \href{https://www.youtube.com/watch?v=r8bK9j4vKco}{Radius of Convergence}. 
Video by Krista King.

\item \href{https://www.youtube.com/watch?v=deGN83_kzpc}{Interval of Convergence Ex 1}. 
Video by Krista King.

\item \href{https://www.youtube.com/watch?v=RWTri1aM7rM}{Interval of Convergence Ex 2}. 
Video by Krista King.

\item \href{https://www.youtube.com/watch?v=EGni2-m5yxM}{Finding the Radius and Interval of Convergence}. 
Video by The Organic Chemistry Tutor.

\end{itemize}

\textbf{Representing Functions as Power Series}

\begin{itemize}

\item \href{https://www.youtube.com/watch?v=Sw7PcBTgE0A}{Representing a Function as a Geometric Power Series (Part 1)}. 
Video by James Sousa.

\item \href{https://www.youtube.com/watch?v=DyJetBjvT4M}{Representing a Function as a Geometric Power Series (Part 2)}. 
Video by James Sousa.

\item \href{https://www.youtube.com/watch?v=CIDaNHGcQ08}{Ex 1: Find a Power Series to Represent a Rational Function}. 
Video by James Sousa.

\item \href{https://www.youtube.com/watch?v=Ps06EgRs3tI}{Ex 2: Find a Power Series to Represent a Rational Function}. 
Video by James Sousa.

\item \href{https://www.youtube.com/watch?v=ALxAnRKG_1c}{Ex 3: Find a Power Series to Represent a Rational Function}. 
Video by James Sousa.

\item \href{https://www.youtube.com/watch?v=XWGPjZK0Yzw}{Power Series Representation of Functions}. 
Video by Patrick JMT.

\item \href{https://www.youtube.com/watch?v=wrDo-4bwB-8}{Power Series Representation, Radius and Interval of Convergence}. 
Video by Krista King.

\item \href{https://www.youtube.com/watch?v=54yldObvvwY}{Power Series Representation of Functions}. 
Video by The Organic Chemistry Tutor.

\end{itemize}

%-------------------------------------------------------------------------------
\subsubsection{Licensing}
\label{sss:lice-7352}
%-------------------------------------------------------------------------------

Content obtained and/or adapted from:

\begin{itemize}
\item \href{https://en.wikibooks.org/wiki/Calculus/Power_series}{Power series, Wikibooks: Calculus}. Licensed under CC BY-SA.
\end{itemize}


\subsection{Assessment Instrument}
\label{ssec:powe-8f90}

Assessment requires finding the radius of convergence using the Ratio Test, testing endpoints separately, and stating the interval of convergence with correct bracket/parenthesis notation.

\textbf{Example problems.}

\begin{enumerate}
    \item Find the radius and interval of convergence for $\sum_{n=0}^{\infty}\dfrac{(x-2)^n}{3^n}$.
    \item Determine the interval of convergence for $\sum_{n=1}^{\infty}\dfrac{(-1)^n x^n}{n}$, testing the endpoints carefully.
    \item Find all $x$ for which $\sum_{n=0}^{\infty}n!\,x^n$ converges.
\end{enumerate}

%===============================================================================
\section{Properties of power series}
\label{sec:prop-066f}
%===============================================================================


%-------------------------------------------------------------------------------
\subsubsection{Addition and subtraction}
\label{sss:addi-546e}
%-------------------------------------------------------------------------------


When two functions \textit{f} and \textit{g} are decomposed into power series around the same center \textit{c}, the power series of the sum or difference of the functions can be obtained by termwise addition and subtraction. That is, if


$\quad$ $f(x) = \sum _{n=0}^\infty a_n (x - c)^n$ and $g(x) = \sum _{n=0}^\infty b_n (x - c)^n$


then


$\quad$ $f(x) \pm g(x) = \sum _{n=0}^\infty (a_n \pm b_n) (x - c)^n.$


It is not true that if two power series $\sum _{n=0}^\infty a_n x^n$ and $\sum _{n=0}^\infty b_n x^n$ have the same radius of convergence, then $\sum _{n=0}^\infty \left(a_n + b_n\right) x^n$ also has this radius of convergence. If $a_n = (-1)^n$ and $b_n = (-1)^{n+1} \left(1 - \frac{1}{3^n}\right)$, then both series have the same radius of convergence of 1, but the series $\sum _{n=0}^\infty \left(a_n + b_n\right) x^n = \sum _{n=0}^\infty \frac{(-1)^n}{3^n} x^n$ has a radius of convergence of 3.



%-------------------------------------------------------------------------------
\subsubsection{Multiplication and Division}
\label{sss:mult-c79b}
%-------------------------------------------------------------------------------

Let
\[
f(x) = \sum_{n=0}^{\infty} a_n (x-c)^n,
\qquad
g(x) = \sum_{n=0}^{\infty} b_n (x-c)^n.
\]

The product is given by:
\begin{align}
f(x)g(x)
&= \left(\sum_{n=0}^{\infty} a_n (x-c)^n\right)
   \left(\sum_{n=0}^{\infty} b_n (x-c)^n\right) \\
&= \sum_{i=0}^{\infty} \sum_{j=0}^{\infty} a_i b_j (x-c)^{i+j} \\
&= \sum_{n=0}^{\infty} \left( \sum_{i=0}^{n} a_i b_{n-i} \right) (x-c)^n.
\end{align}

The sequence
\[
m_n = \sum_{i=0}^{n} a_i b_{n-i}
\]
is called the \textbf{convolution} of the sequences $\{a_n\}$ and $\{b_n\}$.

\medskip

For division, suppose
\[
\frac{f(x)}{g(x)} = \sum_{n=0}^{\infty} d_n (x-c)^n,
\qquad \text{with } b_0 \neq 0.
\]

Then
\[
f(x) = g(x)\sum_{n=0}^{\infty} d_n (x-c)^n,
\]
so by comparing coefficients we obtain a recursive formula:
\[
d_0 = \frac{a_0}{b_0},
\]
\[
d_n = \frac{1}{b_0}\left(a_n - \sum_{i=1}^{n} b_i d_{n-i}\right), \quad n \ge 1.
\]

\medskip

Alternatively, $d_n$ can be expressed using determinants:
\[
d_n = \frac{1}{b_0^{\,n+1}}
\begin{vmatrix}
a_n      & b_1      & b_2      & \cdots & b_n \\
a_{n-1}  & b_0      & b_1      & \cdots & b_{n-1} \\
a_{n-2}  & 0        & b_0      & \cdots & b_{n-2} \\
\vdots   & \vdots   & \vdots   & \ddots & \vdots \\
a_0      & 0        & 0        & \cdots & b_0
\end{vmatrix}.
\]

%-------------------------------------------------------------------------------
\subsubsection{Differentiation and Integration}
\label{sss:diff-fcae}
%-------------------------------------------------------------------------------

If
\[
f(x) = \sum_{n=0}^{\infty} a_n (x-c)^n,
\]
then $f$ is differentiable inside its interval of convergence, and:

\[
f'(x) = \sum_{n=1}^{\infty} a_n n (x-c)^{n-1}
= \sum_{n=0}^{\infty} a_{n+1}(n+1)(x-c)^n,
\]

\[
\int f(x)\,dx
= \sum_{n=0}^{\infty} \frac{a_n}{n+1}(x-c)^{n+1} + K
= \sum_{n=1}^{\infty} \frac{a_{n-1}}{n}(x-c)^n + K.
\]

Both resulting series have the \textbf{same radius of convergence} as the original series.

%-------------------------------------------------------------------------------
\subsubsection{Resources}
\label{sss:reso-power-diff}
%-------------------------------------------------------------------------------

\textbf{Differentiation and Integration of Power Series}

\begin{itemize}
\item \href{https://www.youtube.com/watch?v=zCQuEN7TFNE}{Differentiation and Integration Using Power Series} by James Sousa.
\item \href{https://www.youtube.com/watch?v=Ap47feo6xUM}{Estimate a Definite Integral Using Power Series} by James Sousa.
\item \href{https://www.youtube.com/watch?v=qPl9nr8my2Q}{Differentiating and Integrating Power Series} by Patrick JMT.
\item \href{https://www.youtube.com/watch?v=q12rXAMfEXw}{Integrating a Power Series} by Patrick JMT.
\item \href{https://www.youtube.com/watch?v=mGAE8CWH45M}{Integrating a Power Series, Example 2} by Patrick JMT.
\item \href{https://www.youtube.com/watch?v=2dQ4UvCIBV8}{Power Series Differentiation} by Krista King.
\item \href{https://www.youtube.com/watch?v=5p-EH8vP_bU}{Expressing the Integral as a Power Series} by Krista King.
\item \href{https://www.youtube.com/watch?v=GHvPArQIXSg}{Using Power Series to Estimate a Definite Integral} by Krista King.
\item \href{https://www.youtube.com/watch?v=nD6hai32ykQ}{Power Series - Differentiation and Integration} by The Organic Chemistry Tutor.
\end{itemize}

%-------------------------------------------------------------------------------
\subsubsection{Power Series Representation of Functions Using Differentiation and Integration}
\label{sss:powe-4b59}
%-------------------------------------------------------------------------------

\begin{itemize}
\item \href{https://www.youtube.com/watch?v=bvNPmQRY8Ts}{Find a Power Series to Represent $\arctan(x)$ Using Integration} by Math is Power 4U.
\item \href{https://www.youtube.com/watch?v=ujB-5zQWZAY}{Find a Power Series to Represent a Rational Function Using Differentiation} by Math is Power 4U.
\item \href{https://www.youtube.com/watch?v=JCF3jgmsJuY}{Finding a Power Series by Differentiation} by Patrick JMT.
\item \href{https://www.youtube.com/watch?v=fLJsUpZnt1M}{Finding Power Series by Differentiation - Three Examples} by Patrick JMT.
\item \href{https://www.youtube.com/watch?v=L9xrmXtmMcA}{Power Series Representation Using Integration} by Patrick JMT.
\item \href{https://www.youtube.com/watch?v=PmHaNjDBh_c}{Power Series Representation by Differentiation} by The Organic Chemistry Tutor.
\item \href{https://www.youtube.com/watch?v=0HyM3nM87mk}{Power Series Representation by Integration} by The Organic Chemistry Tutor.
\end{itemize}

%-------------------------------------------------------------------------------
\subsubsection{Licensing}
\label{sss:lice-7352}
%-------------------------------------------------------------------------------


Content obtained and/or adapted from:


\begin{itemize}
    \item \href{https://en.wikipedia.org/wiki/Power_series}{Power series, Wikipedia} under a CC BY-SA license
\end{itemize}

\subsection{Learning Objectives}

\begin{enumerate}
    \item Addition and Subtraction
    \item \textbf{Decompose and Recompose Power Series for Sum of Functions:}
    \item \textbf{Outcome:} Students will be able to decompose functions \( f \) and \( g \) into their respective power series around the same center \( c \) and recompose the power series to represent the sum of the functions.
    \item \textbf{Details:} This involves understanding the power series representation of functions, performing termwise addition, and verifying the resulting power series. For example, given \( f(x) = \sum_{n=0}^\infty a_n (x - c)^n \) and \( g(x) = \sum_{n=0}^\infty b_n (x - c)^n \), students will demonstrate \( f(x) + g(x) = \sum_{n=0}^\infty (a_n + b_n) (x - c)^n \).
    \item \textbf{Decompose and Recompose Power Series for Difference of Functions:}
    \item \textbf{Outcome:} Students will be able to decompose functions \( f \) and \( g \) into their respective power series around the same center \( c \) and recompose the power series to represent the difference of the functions.
    \item \textbf{Details:} This involves understanding the power series representation of functions, performing termwise subtraction, and verifying the resulting power series. For example, given \( f(x) = \sum_{n=0}^\infty a_n (x - c)^n \) and \( g(x) = \sum_{n=0}^\infty b_n (x - c)^n \), students will demonstrate \( f(x) - g(x) = \sum_{n=0}^\infty (a_n - b_n) (x - c)^n \).
    \item \textbf{Analyze Radius of Convergence for Combined Power Series:}
    \item \textbf{Outcome:} Students will be able to analyze the radius of convergence for the combined power series of the sum or difference of functions \( f \) and \( g \) and determine if it changes.
    \item \textbf{Details:} This includes examining cases where the radii of convergence of the original power series are the same, and understanding why the radius of convergence of the sum or difference may differ. For example, with \( a_n = (-1)^n \) and \( b_n = (-1)^{n+1} \left(1 - \frac{1}{3^n}\right) \), students will show that although the series for \( f \) and \( g \) have a radius of convergence of 1, the series for \( f(x) + g(x) \) has a radius of convergence of 3. Multiplication and Division
    \item \textbf{Calculate the Power Series for Product of Functions:}
    \item \textbf{Outcome:} Students will be able to calculate the power series for the product of two functions \( f \) and \( g \) given their individual power series representations.
    \item \textbf{Details:} This involves understanding the convolution of sequences and performing the necessary calculations. Given \( f(x) = \sum_{n=0}^\infty a_n (x - c)^n \) and \( g(x) = \sum_{n=0}^\infty b_n (x - c)^n \), students will compute \( f(x)g(x) = \sum_{n=0}^\infty \left( \sum_{i=0}^n a_i b_{n-i} \right) (x - c)^n \).
    \item \textbf{Solve for Coefficients in Power Series Division:}
    \item \textbf{Outcome:} Students will be able to solve for the coefficients of the power series representation of the quotient of two functions \( f \) and \( g \).
    \item \textbf{Details:} This involves using the recursive relationship derived from the product of the series. Given \( \frac{f(x)}{g(x)} = \sum_{n=0}^\infty d_n (x - c)^n \), students will solve for \( d_n \) using the provided determinant formula.
    \item \textbf{Apply Convolution in Power Series Multiplication:}
    \item \textbf{Outcome:} Students will be able to apply the concept of convolution to find the power series representation of the product of two functions.
    \item \textbf{Details:} This includes performing the convolution of the sequences of coefficients from the power series of the functions. For instance, students will compute \( m_n = \sum_{i=0}^n a_i b_{n-i} \) to find the coefficients for the power series of \( f(x)g(x) \). Differentiation and Integration
    \item \textbf{Differentiate Power Series Term-by-Term:}
    \item \textbf{Outcome:} Students will be able to differentiate a power series term-by-term to find the power series representation of the derivative of a function.
    \item \textbf{Details:} This includes applying the differentiation
\end{enumerate}

\subsection{Assessment Instrument}
\label{ssec:prop-3ce7}

Students differentiate and integrate known power series term-by-term to derive new series, verify that the interval of convergence is preserved (except possibly at endpoints), and apply the results to evaluate limits or integrals.

\textbf{Example problems.}

\begin{enumerate}
    \item Starting from $\dfrac{1}{1-x} = \sum_{n=0}^{\infty}x^n$, differentiate term-by-term to obtain a power series for $\dfrac{1}{(1-x)^2}$ and state its interval of convergence.
    \item Integrate $\dfrac{1}{1+x^2} = \sum_{n=0}^{\infty}(-1)^n x^{2n}$ term-by-term to derive the Maclaurin series for $\arctan x$.
    \item Use the power series for $e^x$ to express $\int_0^1 e^{-x^2}\,dx$ as an infinite series and estimate it to within $0.001$.
\end{enumerate}

%===============================================================================
\section{Taylor and Maclaurin series}
\label{sec:tayl-313c}
%===============================================================================


%-------------------------------------------------------------------------------
\subsubsection{Taylor Series}
\label{sss:tayl-325c}
%-------------------------------------------------------------------------------


\textbf{Definition: Taylor series}
A function $f(x)$ is said to be \textbf{analytic} if it can be represented by the an infinite power series
$\newline$ $\sum _{n=0}^\infty c_n(x-a)^n$
$\newline$ The \textbf{Taylor expansion} or \textbf{Taylor series} representation of a function, then, is
$\newline$ $\sum _{n=0}^\infty\frac{f^{(n)}(a)}{n!}(x-a)^n$

Figure: sin(x) and Taylor approximations, polynomials of degree 1 (red), 3 (lime), 5 (blue), 7 (gold), 9 (cyan), 11 (magnetta), and 13 (grey).


Here, $n!$ is the factorial of $n$ and $f^{(n)}(a)$ denotes the $n$th derivative of $f$ at the point $a$ . If this series converges for every $x$ in the interval $(a-r,a+r)$ and the sum is equal to $f(x)$ , then the function $f(x)$ is called \textbf{analytic}. To check whether the series converges towards $f(x)$, one normally uses estimates for the remainder term of Taylor's theorem. A function is analytic if and only if a power series converges to the function; the coefficients in that power series are then necessarily the ones given in the above Taylor series formula.


If $a=0$ , the series is also called a \textbf{Maclaurin series}.


The importance of such a power series representation is threefold. First, differentiation and integration of power series can be performed term by term and is hence particularly easy. Second, an analytic function can be uniquely extended to a holomorphic function defined on an open disk in the complex plane, which makes the whole machinery of complex analysis available. Third, the (truncated) series can be used to approximate values of the function near the point of expansion.


% [Image: Around 0, the function is very flat]


Figure: The function $f(z) = e^{-1/z^2}$ is not analytic; the Taylor series is 0, but the function is not.


Note that there are examples of infinitely often differentiable functions $f(x)$ whose Taylor series converge, but are \textit{not} equal to $f(x)$ . For instance, for the function defined piecewise by saying that $f(x)=\begin{cases}0 & x=0 \\ e^{-\frac{1}{x^2}} & x\ne 0 \end{cases}$ , all the derivatives are 0 at $x=0$ , so the Taylor series of $f(x)$ is 0, and its radius of convergence is infinite, even though the function most definitely is not 0. This particular pathology does not afflict complex-valued functions of a complex variable. Notice that $f(z)=e^{-\frac{1}{z^2}}$ does not approach 0 as $z$ approaches 0 along the imaginary axis.


Some functions cannot be written as Taylor series because they have a singularity; in these cases, one can often still achieve a series expansion if one allows also negative powers of the variable $x$; see Laurent series. For example, $f(z)=e^{-\frac{1}{z^2}}$ \textit{can} be written as a Laurent series.


The Parker-Sockacki theorem is a recent advance in finding Taylor series which are solutions to differential equations. This theorem is an expansion on the Picard iteration.


%-------------------------------------------------------------------------------
\subsubsection{Derivation}
\label{sss:deri-b600}
%-------------------------------------------------------------------------------

Suppose we want to represent a function as an infinite power series; that is, 
as a polynomial with infinitely many terms:
\[
f(x) = c_0 + c_1(x-a) + c_2(x-a)^2 + c_3(x-a)^3 + \cdots
= \sum_{n=0}^{\infty} c_n (x-a)^n,
\]
where $a$ is the \textbf{center} of the series and $c_0, c_1, c_2, \dots$ are coefficients.

\medskip

To determine the coefficients, we evaluate the function at $x=a$:
\[
f(a) = c_0.
\]

Next, differentiate the series:
\[
f'(x) = c_1 + 2c_2(x-a) + 3c_3(x-a)^2 + \cdots
= \sum_{n=1}^{\infty} n c_n (x-a)^{n-1}.
\]

Evaluating at $x=a$ gives:
\[
f'(a) = c_1.
\]

Differentiating again:
\[
f''(x) = 2c_2 + 2\cdot 3\, c_3 (x-a) + \cdots
= \sum_{n=2}^{\infty} n(n-1)c_n (x-a)^{n-2},
\]
so
\[
f''(a) = 2c_2.
\]

Continuing this process, the $n$th derivative satisfies:
\[
f^{(n)}(x) = n! \, c_n + \text{terms involving } (x-a),
\]
and hence
\[
f^{(n)}(a) = n! \, c_n.
\]

Solving for $c_n$:
\[
c_n = \frac{f^{(n)}(a)}{n!}.
\]

\medskip

Substituting back into the series gives the \textbf{Taylor series}:
\[
f(x) = \sum_{n=0}^{\infty} \frac{f^{(n)}(a)}{n!}(x-a)^n.
\]

When $a=0$, this becomes the \textbf{Maclaurin series}:
\[
f(x) = \sum_{n=0}^{\infty} \frac{f^{(n)}(0)}{n!} x^n.
\]

%-------------------------------------------------------------------------------
\paragraph{Example: Maclaurin Series for $\sin(x)$}
%-------------------------------------------------------------------------------

Let $f(x) = \sin(x)$ and $a=0$.

We compute derivatives:
\[
\sin(x), \quad \cos(x), \quad -\sin(x), \quad -\cos(x), \quad \sin(x), \dots
\]

Evaluating at $x=0$:
\[
\sin(0)=0,\quad \cos(0)=1,\quad -\sin(0)=0,\quad -\cos(0)=-1,\dots
\]

Substituting into the Taylor formula:
\[
\sin(x) = 0 + \frac{1}{1!}x + 0 - \frac{1}{3!}x^3 + 0 + \frac{1}{5!}x^5 - \cdots
\]

Ignoring zero terms, we obtain:
\[
\sin(x) = x - \frac{x^3}{3!} + \frac{x^5}{5!} - \frac{x^7}{7!} + \cdots
\]

In summation form:
\[
\sin(x) = \sum_{n=1}^{\infty} \frac{(-1)^{n-1}}{(2n-1)!} x^{2n-1}.
\]

%-------------------------------------------------------------------------------
\subsubsection{List of Taylor Series}
\label{sss:list-782a}
%-------------------------------------------------------------------------------

Several important Taylor series expansions follow. All these expansions are also valid for complex arguments $x$ (within their respective domains of convergence).

\medskip

\textbf{Exponential function and natural logarithm:}
\[
e^x = \sum_{n=0}^{\infty} \frac{x^n}{n!}
\quad \text{for all } x,
\]
\[
\ln(1+x) = \sum_{n=1}^{\infty} \frac{(-1)^{n-1}}{n} x^n
\quad \text{for } |x| < 1.
\]

\medskip

\textbf{Geometric series:}
\[
\frac{1}{1-x} = \sum_{n=0}^{\infty} x^n
\quad \text{for } |x| < 1.
\]

\medskip

\textbf{Binomial series:}
\[
(1+x)^\alpha = \sum_{n=0}^{\infty} \binom{\alpha}{n} x^n
\quad \text{for } |x| < 1,
\]
valid for all complex $\alpha$.

\medskip

\textbf{Trigonometric functions:}
\begin{align}
\sin(x) &= \sum_{n=1}^{\infty} \frac{(-1)^{n-1}}{(2n-1)!} x^{2n-1}
\quad \text{for all } x, \\[4pt]
\cos(x) &= \sum_{n=0}^{\infty} \frac{(-1)^n}{(2n)!} x^{2n}
\quad \text{for all } x, \\[4pt]
\tan(x) &= \sum_{n=1}^{\infty} \frac{B_{2n} (-4)^n (1-4^n)}{(2n)!} x^{2n-1}
\quad \text{for } |x| < \tfrac{\pi}{2}, \\[4pt]
\sec(x) &= \sum_{n=0}^{\infty} \frac{(-1)^n E_{2n}}{(2n)!} x^{2n}
\quad \text{for } |x| < \tfrac{\pi}{2}, \\[4pt]
\arcsin(x) &= \sum_{n=0}^{\infty}
\frac{(2n)!}{4^n (n!)^2 (2n+1)} x^{2n+1}
\quad \text{for } |x| < 1, \\[4pt]
\arctan(x) &= \sum_{n=0}^{\infty}
\frac{(-1)^n}{2n+1} x^{2n+1}
\quad \text{for } |x| < 1.
\end{align}

\medskip

\textbf{Hyperbolic functions:}
\begin{align}
\sinh(x) &= \sum_{n=1}^{\infty} \frac{x^{2n-1}}{(2n-1)!}
\quad \text{for all } x, \\[4pt]
\cosh(x) &= \sum_{n=0}^{\infty} \frac{x^{2n}}{(2n)!}
\quad \text{for all } x, \\[4pt]
\tanh(x) &= \sum_{n=1}^{\infty}
\frac{B_{2n} 2^{2n}(2^{2n}-1)}{(2n)!} x^{2n-1}
\quad \text{for } |x| < \tfrac{\pi}{2}, \\[4pt]
\operatorname{arsinh}(x) &= \sum_{n=0}^{\infty}
\frac{(-1)^n (2n)!}{2^{2n}(n!)^2 (2n+1)} x^{2n+1}
\quad \text{for } |x| < 1, \\[4pt]
\operatorname{artanh}(x) &= \sum_{n=1}^{\infty}
\frac{x^{2n-1}}{2n-1}
\quad \text{for } |x| < 1.
\end{align}

\medskip

\textbf{Lambert W function:}
\[
W_0(x) = \sum_{n=1}^{\infty}
\frac{(-n)^{n-1}}{n!} x^n
\quad \text{for } |x| < \tfrac{1}{e}.
\]

\medskip

The numbers $B_k$ appearing in the expansions of $\tan(x)$ and $\tanh(x)$ are the \textbf{Bernoulli numbers}.  
The coefficients $\binom{\alpha}{n}$ are the generalized binomial coefficients.  
The numbers $E_k$ in the expansion of $\sec(x)$ are the \textbf{Euler numbers}.

%-------------------------------------------------------------------------------
\subsubsection{Multiple Dimensions}
\label{sss:mult-7324}
%-------------------------------------------------------------------------------

The Taylor series generalizes to functions of several variables. For a function
$f(x_1,\dots,x_d)$ expanded about $(a_1,\dots,a_d)$:
\[
f(x_1,\dots,x_d)
=
\sum_{n_1=0}^{\infty} \cdots \sum_{n_d=0}^{\infty}
\frac{\partial^{n_1+\cdots+n_d} f}{\partial x_1^{n_1}\cdots \partial x_d^{n_d}}(a_1,\dots,a_d)
\,
\frac{(x_1-a_1)^{n_1}\cdots (x_d-a_d)^{n_d}}{n_1!\cdots n_d!}.
\]

%-------------------------------------------------------------------------------
\subsubsection{History}
\label{sss:hist-3cd1}
%-------------------------------------------------------------------------------

The Taylor series is named after the mathematician \textbf{Brook Taylor}, who first published the general formula in 1715.

%-------------------------------------------------------------------------------
\subsubsection{Constructing a Taylor Series}
\label{sss:cons-9461}
%-------------------------------------------------------------------------------


Several methods exist for the calculation of Taylor series of a large number of functions. One can attempt to use the Taylor series as-is and generalize the form of the coefficients, or one can use manipulations such as substitution, multiplication or division, addition or subtraction of standard Taylor series (such as those above) to construct the Taylor series of a function, by virtue of Taylor series being power series. In some cases, one can also derive the Taylor series by repeatedly applying integration by parts. The use of computer algebra systems to calculate Taylor series is common, since it eliminates tedious substitution and manipulation.


%-------------------------------------------------------------------------------
\subsubsection{Example 1}
\label{sss:exam-acc3}
%-------------------------------------------------------------------------------

Consider the function
\[
f(x) = \ln\big(1+\cos(x)\big),
\]
for which we want a Taylor series at $0$.

\medskip

We recall:
\[
\ln(1+u) = \sum_{n=1}^{\infty} \frac{(-1)^{n+1}}{n} u^n
= u - \frac{u^2}{2} + \frac{u^3}{3} - \frac{u^4}{4} + \cdots,
\quad |u|<1,
\]
and
\[
\cos(x) = \sum_{n=0}^{\infty} \frac{(-1)^n}{(2n)!} x^{2n}
= 1 - \frac{x^2}{2!} + \frac{x^4}{4!} - \cdots.
\]

\medskip

To use the logarithm expansion, write:
\[
1+\cos(x) = 2 + \left(-\frac{x^2}{2} + \frac{x^4}{24} - \cdots\right)
= 2\left(1 + u(x)\right),
\]
where
\[
u(x) = -\frac{x^2}{4} + \frac{x^4}{48} - \cdots.
\]

Thus,
\[
\ln(1+\cos(x)) = \ln(2) + \ln(1+u(x)).
\]

Now expand:
\[
\ln(1+u(x)) = u(x) - \frac{u(x)^2}{2} + \frac{u(x)^3}{3} - \cdots.
\]

Substituting and simplifying gives:
\[
\ln(1+\cos(x))
= \ln(2)
- \frac{x^2}{4}
- \frac{x^4}{96}
- \frac{x^6}{1440}
- \frac{17x^8}{322560}
- \cdots.
\]

\medskip

Since $\cos(x)$ is even, $f(x)$ is also even, so all odd powers vanish.

\medskip

The general coefficients can be expressed using Faà di Bruno's formula, though this is typically not illuminating.

%-------------------------------------------------------------------------------
\subsubsection{Example 2}
\label{sss:exam-5bdd}
%-------------------------------------------------------------------------------

Suppose we want the Taylor series at $0$ of the function
\[
g(x) = \frac{e^x}{\cos(x)}.
\]

We use:
\[
e^x = \sum_{n=0}^{\infty} \frac{x^n}{n!}
= 1 + x + \frac{x^2}{2!} + \frac{x^3}{3!} + \cdots,
\]
\[
\cos(x) = 1 - \frac{x^2}{2!} + \frac{x^4}{4!} - \cdots.
\]

Assume:
\[
\frac{e^x}{\cos(x)} = \sum_{n=0}^{\infty} c_n x^n.
\]

Multiplying both sides by $\cos(x)$:
\[
e^x = \left(\sum_{n=0}^{\infty} c_n x^n\right)
\left(1 - \frac{x^2}{2} + \frac{x^4}{24} - \cdots\right).
\]

Expanding up to fourth order:
\[
\begin{aligned}
e^x
&= c_0 + c_1 x + c_2 x^2 + c_3 x^3 + c_4 x^4 \\
&\quad - \frac{c_0}{2} x^2 - \frac{c_1}{2} x^3 - \frac{c_2}{2} x^4 \\
&\quad + \frac{c_0}{24} x^4 + \cdots.
\end{aligned}
\]

Collecting terms:
\[
\begin{aligned}
e^x
&= c_0
+ c_1 x
+ \left(c_2 - \frac{c_0}{2}\right)x^2 \\
&\quad + \left(c_3 - \frac{c_1}{2}\right)x^3 \\
&\quad + \left(c_4 - \frac{c_2}{2} + \frac{c_0}{24}\right)x^4
+ \cdots.
\end{aligned}
\]

Matching coefficients with
\[
e^x = 1 + x + \frac{x^2}{2} + \frac{x^3}{6} + \frac{x^4}{24} + \cdots,
\]
we obtain:
\[
c_0 = 1,\quad c_1 = 1,
\]
\[
c_2 - \frac{1}{2} = \frac{1}{2} \;\Rightarrow\; c_2 = 1,
\]
\[
c_3 - \frac{1}{2} = \frac{1}{6} \;\Rightarrow\; c_3 = \frac{2}{3},
\]
\[
c_4 - \frac{1}{2} + \frac{1}{24} = \frac{1}{24}
\;\Rightarrow\; c_4 = \frac{1}{2}.
\]

Thus,
\[
\frac{e^x}{\cos(x)}
= 1 + x + x^2 + \frac{2}{3}x^3 + \frac{1}{2}x^4 + \cdots.
\]


%-------------------------------------------------------------------------------
\subsubsection{Resources}
\label{sss:reso-55b5}
%-------------------------------------------------------------------------------


%-------------------------------------------------------------------------------
\subsubsection{Taylor and Maclaurin Series}
\label{sss:tayl-313c}
%-------------------------------------------------------------------------------


\begin{itemize}
    \item \href{https://www.youtube.com/watch?v=cwpJxfHkR1o}{Taylor and Maclaurin Series} by James Sousa, Math is Power 4U
    \item \href{https://www.youtube.com/watch?v=jKq-TdR9GG8}{Ex: Find the Taylor Series of x^3} by James Sousa, Math is Power 4U
    \item \href{https://www.youtube.com/watch?v=vzn7EdBQgBk}{Ex: Find the Taylor Series of e^x} by James Sousa, Math is Power 4U
    \item \href{https://www.youtube.com/watch?v=Zt7QQt-3-xA}{Determine the Maclaurin Series and Polynomial} by James Sousa, Math is Power 4U
    \item \href{https://www.youtube.com/watch?v=xgxvpkN-jrU}{Determine the Maclaurin Series and Polynomial of acos(bx^2)} by James Sousa, Math is Power 4U
    \item \href{https://www.youtube.com/watch?v=vdFqHLY6elE}{Determine the Maclaurin Series and Polynomial of ax^2e^(bx)} by James Sousa, Math is Power 4U
    \item \href{https://www.youtube.com/watch?v=TfipAnX2sWs}{Determine the Maclaurin Series and Polynomial of ax^2sin(bx)} by James Sousa, Math is Power 4U
    \item \href{https://www.youtube.com/watch?v=GUtLtRDox3c}{Taylor/Maclaurin Series Expansion, Proof of the Formula} by Patrick JMT
    \item \href{https://www.youtube.com/watch?v=cjPoEZ0I5wQ}{Taylor and Maclaurin Series, Example 1} by Patrick JMT
    \item \href{https://www.youtube.com/watch?v=Os8OtXFBLkY}{Taylor and Maclaurin Series, Example 2} by Patrick JMT
    \item \href{https://www.youtube.com/watch?v=OADFAVC9EgQ}{Finding a Maclaurin Series Expansion - Another Example} by Patrick JMT
    \item \href{https://www.youtube.com/watch?v=jGAWFkLRr3s}{Finding a New Power Series by Manipulating a Known Power Series Ex 1} by Patrick JMT
    \item \href{https://www.youtube.com/watch?v=cT5IXkzcXJQ}{Finding a New Power Series by Manipulating a Known Power Series Ex 2} by Patrick JMT
    \item \href{https://www.youtube.com/watch?v=GJFKLyC48-8}{Maclaurin Series} by Krista King
    \item \href{https://www.youtube.com/watch?v=lv9Kqk7rxJw}{Maclaurin Series Radius of Convergence} by Krista King
    \item \href{https://www.youtube.com/watch?v=LDBnS4c7YbA}{Taylor Series and Maclaurin Series} by The Organic Chemistry Tutor
    \item \href{https://www.youtube.com/watch?v=3d6DsjIBzJ4}{Taylor Series} by 3Blue1Brown
\end{itemize}


%-------------------------------------------------------------------------------
\subsubsection{Taylor and Maclaurin Polynomials}
\label{sss:tayl-487b}
%-------------------------------------------------------------------------------


\begin{itemize}
    \item \href{https://www.youtube.com/watch?v=lShzrGyiueg}{Taylor Polynomials} by James Sousa, Math is Power 4U
    \item \href{https://www.youtube.com/watch?v=SC6TsVsP90U}{Find a Degree One and Degree Two Maclaurin Polynomial} by James Sousa, Math is Power 4U
    \item \href{https://www.youtube.com/watch?v=39s-mgfuuno}{Determine a Taylor Polynomial of the Square Root Function} by James Sousa, Math is Power 4U
    \item \href{https://www.youtube.com/watch?v=6eNE38s9Fv0}{Finding a Maclaurin Polynomial, Ex 1} Patrick JMT
    \item \href{https://www.youtube.com/watch?v=tYl7tXKKkSw}{Finding a Maclaurin Polynomial, Ex 2} Patrick JMT
    \item \href{https://www.youtube.com/watch?v=WKvBJdl1qrM}{Finding a Taylor Polynomial to Approximate a Function, Ex 1} by Patrick JMT
    \item \href{https://www.youtube.com/watch?v=w_oDLT5VnBI}{Finding a Taylor Polynomial to Approximate a Function, Ex 2} by Patrick JMT
    \item \href{https://www.youtube.com/watch?v=UINFWG0ErSA}{Finding a Taylor Polynomial to Approximate a Function, Ex 3} by Patrick JMT
    \item \href{https://www.youtube.com/watch?v=em4YU8rUhco}{Finding a Taylor Polynomial to Approximate a Function, Ex 4} by Patrick JMT
    \item \href{https://www.youtube.com/watch?v=C1vcVMotnkg}{Taylor Polynomial Example 1 Part 1} by Krista King
    \item \href{https://www.youtube.com/watch?v=JMU-4o-Ffus}{Taylor Polynomial Example 1 Part 2} by Krista King
    \item \href{https://www.youtube.com/watch?v=h2EGCv0PyZ8}{Taylor Polynomial Example 2 Part 1} by Krista King
    \item \href{https://www.youtube.com/watch?v=Fw9u5h4MoJo}{Taylor Polynomial Example 2 Part 2} by Krista King
    \item \href{https://www.youtube.com/watch?v=kd5nuHkKsr4}{Taylor Polynomial Example 3 Part 1} by Krista King
    \item \href{https://www.youtube.com/watch?v=l5eqQFfMp-0}{Taylor Polynomial Example 3 Part 2} by Krista King
    \item \href{https://www.youtube.com/watch?v=NWl-dwL7_Do}{Taylor Polynomial Example 3 Part 3} by Krista King
    \item \href{https://www.youtube.com/watch?v=urPIxvNBXF0}{Taylor Polynomials and Maclaurin Polynomials with Approximations} by The Organic Chemistry Tutor
\end{itemize}


%-------------------------------------------------------------------------------
\subsubsection{Taylor's Inequality / Taylor's Remainder Theorem}
\label{sss:tayl-341c}
%-------------------------------------------------------------------------------


\begin{itemize}
    \item \href{https://www.youtube.com/watch?v=DP_pGQaNGdw}{Taylor's Theorem with Remainder} by James Sousa, Math is Power 4U
    \item \href{https://www.youtube.com/watch?v=wqgAG9TtNAA}{Ex: Find a Maclaurin Polynomial and Error of an Approximation} by James Sousa, Math is Power 4U
    \item \href{https://www.youtube.com/watch?v=CExJJShnzzg}{Ex: Find a Maclaurin Polynomial and the Interval for a Given Error} by James Sousa, Math is Power 4U
    \item \href{https://www.youtube.com/watch?v=lwEU4ZkwcJA}{Ex: Find a Maclaurin Polynomial and the Interval for a Given Error} by James Sousa, Math is Power 4U
    \item \href{https://www.youtube.com/watch?v=9LnCcAoEHG4}{Taylor's Inequality} by Patrick JMT
    \item \href{https://www.youtube.com/watch?v=yJWFAG1g4Jw}{Taylor's Remainder Theorem, Finding the Remainder Ex 1} by Patrick JMT
    \item \href{https://www.youtube.com/watch?v=30n1eiu8lFs}{Taylor's Remainder Theorem, Finding the Remainder Ex 2} by Patrick JMT
    \item \href{https://www.youtube.com/watch?v=pZBRc5uegEs}{Taylor's Remainder Theorem, Finding the Remainder Ex 3} by Patrick JMT
    \item \href{https://www.youtube.com/watch?v=aa0hyXZH9vw}{Taylor's Remainder Theorem, Finding the Remainder Ex 4} by Patrick JMT
    \item \href{https://www.youtube.com/watch?v=xsCTDZQULJ4}{Taylor's Inequality} by Krista King
    \item \href{https://www.youtube.com/watch?v=lY0LzJXTgeo}{Taylor's Remainder Theorem} by The Organic Chemistry Tutor
    \item \href{https://www.youtube.com/watch?v=Q9y43_02vFw}{Lagrange Error Bound to Find Error when Using Taylor Polynomials} by Patrick JMT
\end{itemize}


%-------------------------------------------------------------------------------
\subsubsection{Applications of Taylor and Maclaurin Series}
\label{sss:appl-916e}
%-------------------------------------------------------------------------------


\begin{itemize}
    \item \href{https://www.youtube.com/watch?v=Ifnnuk6UeNE}{Ex: Use a Maclaurin Polynomial to Approximate an Integral} by James Sousa, Math is Power 4U
    \item \href{https://www.youtube.com/watch?v=vfrwoCGdXSo}{Ex: Determine a Simplified Power Series for a Function Involving e^(ax)} by James Sousa, Math is Power 4U
    \item \href{https://www.youtube.com/watch?v=5RpUtQBFNdM}{Find the Sum of an Infinite Series Using a Known Power Series} by James Sousa, Math is Power 4U
    \item \href{https://www.youtube.com/watch?v=Wv0SuLhHa4k}{Find the Sum of an Infinite Series Using a Known Power Series} by James Sousa, Math is Power 4U
    \item \href{https://www.youtube.com/watch?v=XN_CqqH89eU}{Find the Sum of an Infinite Series Using a Known Power Series} by James Sousa, Math is Power 4U
    \item \href{https://www.youtube.com/watch?v=3ZOS69YTxQ8}{Using Maclaurin/Taylor Series to Approximate a Definite Integral to a Desired Accuracy} by Patrick JMT
    \item \href{https://www.youtube.com/watch?v=7TFpx8DbfUQ}{Maclaurin Series to Evaluate a Limit} by Krista King
    \item \href{https://www.youtube.com/watch?v=dZhqQbqLxks}{Sum of the Maclaurin Series} by Krista King
    \item \href{https://www.youtube.com/watch?v=Tq6pNCwI0cM}{Maclaurin Series to Estimate a Definite Integral} by Krista King
    \item \href{https://www.youtube.com/watch?v=4WEGEC7mpQU}{Maclaurin Series to Estimate an Indefinite Integral} by Krista King
\end{itemize}


%-------------------------------------------------------------------------------
\subsubsection{Licensing}
\label{sss:lice-7352}
%-------------------------------------------------------------------------------


Content obtained and/or adapted from:


\begin{itemize}
    \item \href{https://en.wikibooks.org/wiki/Calculus/Taylor_series}{Taylor series, Wikibooks:
\end{itemize}
    Calculus}

    under a CC BY-SA license

\subsection{Learning Objectives}

\begin{enumerate}
    \item \textbf{Derive Taylor Series:}
    \item \textbf{Outcome:} Students will be able to derive the Taylor series for a given function at a specific point.
    \item \textbf{Details:} This includes finding the coefficients of the Taylor series by computing the derivatives of the function at the point and using the Taylor series formula.
    \item \textbf{Convergence of Taylor Series:}
    \item \textbf{Outcome:} Students will be able to determine the interval of convergence for a given Taylor series.
    \item \textbf{Details:} This involves using methods such as the ratio test or root test to identify the values of $x$ for which the Taylor series converges to the function.
    \item \textbf{Applications of Taylor Series:}
    \item \textbf{Outcome:} Students will be able to apply Taylor series to approximate functions and solve problems in physics and engineering.
    \item \textbf{Details:} This includes using Taylor series to approximate values of functions near a given point and applying these approximations to solve differential equations and other applied problems.
\end{enumerate}

\subsection{Assessment Instrument}
\label{ssec:tayl-02c1}

Assessment covers derivation of Taylor/Maclaurin series from the definition, manipulation of known series to derive new ones, determining convergence using Taylor's Remainder Theorem, and applying series to evaluate limits and approximate integrals.

\textbf{Example problems.}

\begin{enumerate}
    \item Derive the Maclaurin series for $\cos x$ from the definition and determine its interval of convergence.
    \item Use the known series for $e^x$ to write a power series for $e^{-x^2}$ and use it to approximate $\int_0^{0.5}e^{-x^2}\,dx$ to four decimal places.
    \item Find the Taylor series for $f(x) = \ln x$ centred at $a=1$ and use Taylor's Inequality to bound the error when truncating after three non-zero terms on $[0.9, 1.1]$.
\end{enumerate}

%===============================================================================
\section{Parametric Equations}
\label{sec:parametric-equations}
%===============================================================================

This wasn't listed in the Wiki under MAT 1224.